\xchapter{Third Session, held February 4, 1546.}

Sessio Tertia,

Third Session,

\emph{celebrata die IV. Februarii} 1546.

\emph{held February} 4, 1546.

DECRETUM DE SYMBOLO FIDEI.

DECREE TOUCHING THE SYMBOL OF FAITH.

\emph{In nomine sanctæ et individuæ Trinitatis, Patris, et Filii, et Spiritus sancti.}

In the name of the Holy and Undivided Trinity, Father, and Son, and Holy Ghost.

\emph{Hæc sacrosancta, œcumenica, et generalis tridentina synodus, in Spiritu sancto legitime congregata, in ea præsidentibus eisdem tribus apostolicæ sedis legatis, magnitudinem rerum tractandarum considerans, præsertim earum, quæ duobus illis capitibus, de extirpandis hæresibus, et moribus reformandis, continentur, quorum causa præcipue est congregata; agnoscens autem cum apostolo, non esse sibi colluctationem adversus carnem et sanguimem, sed adversus spirituales neguitias in cœlestibus, cum eodem omnes et singulos in primis}

This sacred and holy, œcumenical, and general Synod of Trent,—lawfully assembled in the Holy Ghost, the same three legates of the Apostolic See presiding therein,—considering the magnitude of the matters to be treated of, especially of those comprised under the two heads, of the extirpating of heresies, and the reforming of manners, for the sake of which chiefly it is assembled, and recognizing with the apostles, that its \emph{wrestling is not against flesh and blood, but against the spirits of wickedness in the high places},\footnote{\par   Ephes. vi. 12.} exhorts, with the same apostle, all and each, above all

\emph{hortatur, ut confortentur in Domino, et in potential, virtutis eius, in omnibus sumentes scutum fidei, in quo possint omnia tela nequissimi ignea extinguere, atque galeam spei salutis accipiant cum gladio spiritus quod est verbum Dei. Itaque, ut hæc pia eius sollicitudo principium et progressum suum per Dei gratiam habeat, ante omnia statuit et decernit præmittendam esse confessionem fidei, patrum exempla in hoc secuta, qui sacratioribus conciliis hoc scutum contra omnes hæreses in principio suarum actionum apponere consuevere: quo solo aliquando et infideles ad fidem traxerunt, hæreticos expugnarunt, et fideles confirmarunt. Quare symbolum fidei, quo sancta romana ecclesia utitur, tanquam principium illud, in quo omnes, qui fidem Christi profitentur, necessario conveniunt, ac fundamentum firmum et unicum, contra quod portæ inferi nunquam prævalebunt, totidem verbis, quibus in omnibus ecclesiis legitur, experimendum esse censuit; quod quidem eiusmodi est:}

things, to be \emph{strengthened in the Lord, and in the might of his power, in all things taking the shield of faith, wherewith they may be able to extinguish all the fiery darts of the most wicked one, and to take the helmet of salvation, with the sword of the Spirit, which is the word of God?}\footnote{\par   Ephes. vi. 16, 17.} Wherefore, that this its pious solicitude may begin and proceed by the grace of God, it ordains and decrees that, before all other things, a confession of faith is to be set forth; following herein the examples of the Fathers, who have been wont, in the most sacred councils, at the beginning of the Actions thereof, to oppose this shield against heresies; and with this alone, at times, have they drawn the unbelieving to the faith, overthrown heretics, and confirmed the faithful. For which cause, this Council has thought good, that the Symbol of faith which the holy Roman Church makes use of,—as being that principle wherein all who profess the faith of Christ necessarily agree, and that firm and alone foundation \emph{against which the gates of hell shall never prevail?}\footnote{\par   Matt. xvi. 18.} —be expressed in the very same words in which it is read in all the churches. Which Symbol is as follows:

\emph{Credo in unum Deum Patrem omnipotentem, factorem cœli et terra, visibilium omnium et invisibilium; et in unum Dominum Iesum Christum, Filium Dei unigenitum, et ex Patre natum ante omnia sæcula; Deum de Deo, lumen de lumine, Deum verum de Deo vero; genitum, non factum, consubstantialem Patri, per quem omnia facta sunt: qui propter nos homines et propter nostram salutem descendit de cœlis, et incarnatas est de Spiritu Sancto ex Maria virgine, et homo factus est: crucifixus etiam pro nobis sub Pontio Pilato, passus, et sepultus est: et resurrexit tertia die secundum Scripturas, et ascendit in cœlum, sedet ad dexteram Patris, et iterum venturus est cum gloria iudicare vivos et mortuos; cuius regni non erit finis: et in Spiritum Sanctum, Dominum et vivificantem, qui ex Patre Filioque procedit; qui cum Patre et Filio simul adoratur et conglorificatur; qui locutus est per prophetas: et unam sanctam catholicam et apostolicam ecclesiam. Confiteor unum baptisma in remissionem peccatorum: et expecto resurrectionem mortuorum et vitam venturi sæculi. Amen. }

I believe in one God, the Father Almighty, Maker of heaven and earth, of all things visible and invisible; and in one Lord Jesus Christ, the only-begotten Son of God, and born of the Father before all ages; God of God, light of light, true God of true God; begotten, not made, consubstantial with the Father, by whom all things were made: who for us men, and for our salvation, came down from the heavens, and was incarnate by the Holy Ghost of the Virgin Mary, and was made man: crucified also for us under Pontius Pilate, he suffered and was buried; and he rose again on the third day, according to the Scriptures; and he ascended into heaven, sitteth at the right hand of the Father; and again he will come with glory to judge the living and the dead; of whose kingdom there shall be no end: and in the Holy Ghost, the Lord, and the giver of life, who proceedeth from the Father and the Son; who with the Father and the Son together is adored and glorified; who spoke by the prophets: and one holy Catholic and Apostolic Church. I confess one baptism for the remission of sins; and I look for the resurrection of the dead, and the life of the world to come. Amen.

\xchapter{Fourth Session, held April 8, 1546.}

Sessio Quarta,

Fourth Session,

\emph{celebrata die VIII. Aprilis}, 1546.

\emph{held April} 8, 1546.

DECRETUM DE CONONICIS SCRIPTURIS.

DECREE CONCERNING THE CANONICAL SCRIPTURES.

\emph{Sacrosancta, œcumenica, et generalis tridentina synodus, in Spiritu Sancto legitime congregata, præsidentibus in ea eisdem tribus apostolicæ sedis legatis, hoc sibi perpetuo ante oculos proponens, ut, sublatis erroribus, puritas ipsa evangelii in ecclesia conservetur; quod promissum }

The sacred and holy, œcumenical, and general Synod of Trent,—lawfully assembled in the Holy Ghost, the same three legates of the Apostolic See presiding therein,—keeping this always in view, that, errors being removed, the purity itself of the Gospel be preserved in the Church; which (Gospel), before

\emph{ante per prophetas in Scripturis sanctis, Dominus noster Iesus Christus, Dei Filius, proprio ore primum promulgavit, deinde per suos apostolos, tanquam fontem omnis et salutaris veritatis et morum disciplinæ, omni creaturæ prædicari iussit; perspiciensque hanc veritatem et disciplinam contineri in libris scriptis et sine scripto traditionibus, quæ ab ipsius Christi ore ab apostolis acceptæ, aut ab ipsis apostolis, Spiritu Sancto dictante, quasi per manus traditæ, ad nos usque pervenerunt: orthodoxorum patrum exempla secuta, omnes libros tam Veteris quam Novi Testamenti, cum utriusque unus Deus sit auctor, necnon traditiones ipsas, tum ad fidem, tum ad mores pertinentes, tanquam vel oretenus a Christo vel a Spiritu Sancto dictatas, et continua successione in ecclesia catholica conservatas, pari pietatis affectu ac reverentia suscipit et veneratur. }

promised through the prophets in the holy Scriptures, our Lord Jesus Christ, the Son of God, first promulgated with His own mouth, and then commanded to be preached by His Apostles to every creature, as the fountain of all, both saving truth, and moral discipline; and seeing clearly that this truth and discipline are contained in the written books, and the unwritten traditions which, received by the Apostles from the mouth of Christ himself, or from the Apostles themselves, the Holy Ghost dictating, have come down even unto us, transmitted as it were from hand to hand: [the Synod] following the examples of the orthodox Fathers, receives and venerates with an equal affection of piety and reverence, all the books both of the Old and of the New Testament—seeing that one God is the author of both—as also the said traditions, as well those appertaining to faith as to morals, as having been dictated, either by Christ's own word of mouth, or by the Holy Ghost, and preserved in the Catholic Church by a continuous succession.

\emph{Sacrorum vero librorum indicem huic decreto adscribendum censuit, ne cui dubitatio suboriri possit, quinam sint, qui ab ipsa synodo suscipiuntur. Sunt vero }

And it has thought it meet that a list of the sacred books be inserted in this decree, lest a doubt may arise in any one's mind, which are the books that are received by this

\emph{infrascripti. Testamenti veteris, quinque Moysis, id est, Genesis, Exodus, Leviticus, Numeri, Deuteronomium; Iosuæ, Iudicum, Ruth, quatuor Regum, duo Paralipomenon, Esdræ primus et secundus, qui dicitur Nehemias, Tobias, Iudith, Esther, Iob, Psalterium davidicum centum quinquaginta psalmorum, Parabolæ, Ecclesiastes, Canticum canticorum, Sapientia, Ecclesiasticus, Isaias, Ieremias cum Baruch, Ezechiel, Daniel, duodecim prophetæ minores, id est: Osea, Ioel, Amos, Abdias, Ionas, Michæas, Nahum, Habacuc, Sophonias, Aggæus, Zacharias, Malachias, duo Machabæorum, primus et secundus. Testamenti novi: quatuor evangelia, secundum Mathæum, Marcum, Lucam, et Ioannem; actus apostolorum a Luca evangelista conscripti; quatuordecim epistolæ Pauli apostoli, ad Romanos, duæ ad Corinthios, ad Galatas, ad Ephesios, ad Philippenses, ad Colossenses, duæ ad Thessalonicenses, duæ ad Timotheum, ad Titum, ad Philemonem, ad Hebræos; Petri apostoli duæ, Ioannis apostoli tres, Iacobi apostoli una, Iudæ apostoli una, et apocalypsis Ioannis apostoli.}

Synod. They are as set down here below: of the Old Testament: the five books of Moses, to wit, Genesis, Exodus, Leviticus, Numbers, Deuteronomy; Josue, Judges, Ruth, four books of Kings, two of Paralipomenon, the first book of Esdras, and the second which is entitled Nehemias; Tobias, Judith, Esther, Job, the Davidical Psalter, consisting of a hundred and fifty psalms; the Proverbs, Ecclesiastes, the Canticle of Canticles, Wisdom, Ecclesiasticus, Isaias, Jeremias, with Baruch; Ezechiel, Daniel; the twelve minor prophets, to wit, Osee, Joel, Amos, Abdias, Jonas, Micheas, Nahum, Habacuc, Sophonias, Aggæus, Zacharias, Malachias; two books of the Machabees, the first and the second. Of the New Testament: the four Gospels, according to Matthew, Mark, Luke, and John; the Acts of the Apostles written by Luke the Evangelist; fourteen epistles of Paul the apostle, (one) to the Romans, two to the Corinthians, (one) to the Galatians, to the Ephesians, to the Philippians, to the Colossians, two to the Thessalonians, two to Timothy, (one) to Titus, to Philemon, to the Hebrews; two of Peter the apostle, three of John the apostle, one of the apostle James, one of Jude the apostle, and the Apocalypse of John the apostle.

\emph{Si quis autem libros ipsos integros cum omnibus suis partibus, prout in ecclesia catholica legi consueverunt, et in veteri Vulgata Latina editione habentur, pro sacris, et canonicis non susceperit, et traditiones prædictas sciens et prudens contempserit, anathema sit. Omnes itaque intelligant, quo ordine et via ipsa synodus, post jactum fidei confessionis fundamentum, sit progressura, et quibus potissimum testimoniis ac præsidiis in confirmandis dogmatibus et instaurandis in ecclesia moribus sit usura.}

But if any one receive not, as sacred and canonical, the said books entire with all their parts, as they have been used to be read in the Catholic Church, and as they are contained in the old Latin vulgate edition; and knowingly and deliberately contemn the traditions aforesaid; let him be anathema. Let all, therefore, understand, in what order, and in what manner, the said Synod, after having laid the foundation of the Confession of faith, will proceed, and what testimonies and authorities it will mainly use in confirming dogmas, and in restoring morals in the Church.

DECRETUM DE EDITIONE, ET USU SACRORUM LIBRORUM.

DECREE CONCERNING THE EDITION, AND THE USE, OF THE SACRED BOOKS.

\emph{Insuper eadem sacrosancta synodus considerans, non parum utilitatis accedere posse ecclesiæ Dei, si ex omnibus Latinis editionibus, quæ circumferuntur, sacrorum librorum, quænam pro authentica habenda sit, innotescat; statuit et declarat, ut hæc ipsa vetus et vulgata editio, quæ longo tot sæculorum usu in ipsa ecclesia probata est, in publicis lectionibus, disputationibus, prædicationibus et expositionibus pro authentica habeatur; et ut nemo illam rejicere quovis prætextu audeat vel præsumat.}

Moreover, the same sacred and holy Synod,—considering that no small utility may accrue to the Church of God, if it be made known which out of all the Latin editions, now in circulation, of the sacred books, is to be held as authentic,—ordains and declares, that the said old and vulgate edition, which, by the lengthened usage of so many ages, has been approved of in the Church, be, in public lectures, disputations, sermons, and expositions, held as authentic; and that no one is to dare, or presume to reject it under any pretext whatever.

\emph{Præterea, ad coercenda petulantia ingenia, decernit, ut nemo, suæ prudentiæ innixus, in rebus fidei, et morum ad ædificationem doctrinæ christianæ pertinentium, sacram scripturam ad suos sensus contorquens, contra eum sensum, quem tenuit et tenet sancta mater ecclesia, cuius est judicare de vero sensu, et interpretatione scripturarum sanctarum, aut etiam contra unanimem consensum patrum ipsam scripturam sacram interpretari audeat, etiamsi hujusmodi interpretationes nullo unquam tempore in lucem edendæ forent. Qui contravenerint, per ordinarios declarentur, et pœnis a jure statutis puniantur.}

Furthermore, in order to restrain petulant spirits, it decrees, that no one, relying on his own skill, shall,—in matters of faith, and of morals pertaining to the edification of Christian doctrine,—wresting the sacred Scripture to his own senses, presume to interpret the said sacred Scripture contrary to that sense which holy mother Church,—whose it is to judge of the true sense and intrepretation of the holy Scriptures,—hath held and doth hold; or even contrary to the unanimous consent of the Fathers; even though such interpretations were never [intended] to be at any time published. Contraveners shall be made known by their Ordinaries, and be punished with the penalties by law established.

\xchapter{Fifth Session, held June 17, 1546.}

Sessio Quinta,

Fifth Session,

\emph{celebrata die XVII. Junii}, 1546.

\emph{held June} 17, 1546.

DECRETUM DE PECCATO ORIGINALI.

DECREE CONCERNING ORIGINAL SIN.

\emph{Ut fides nostra catholica, sine qua impossibile est placere Deo, purgatis erroribus, in sua sinceritate integra et illibata permaneat; et ne populus christianus omni vento doctrinæ circumferatur; cum serpens ille antiquus, humani generis perpetuus hostis,}

That our Catholic \emph{faith, without which it is impossible to please God},\footnote{\par   Heb. xi. 6.} may, errors being purged away, continue in its own perfect and spotless integrity, and that the Christian people may not \emph{be carried about with every wind of doctrine};\footnote{\par   Ephes. iv. 14.} whereas that old serpent, the perpetual

\emph{inter plurima mala, quibus ecclesia Dei his nostris temporibus perturbatur, etiam de peccato originali ejusque remedio non solum nova, sed vetera etiam dissidia excitaverit: sacrosancta œcumenica et generalis Tridentina synodus, in Spiritu Sancto legitime congregata, præsidentibus in ea eisdem tribus apostolicæ sedis legatis, jam ad revocandos errantes et nutantes confirmandos accedere volens, sacrarum scripturarum et sanctorum patrum ac probatissimorum conciliorum testimonia et ipsius ecclesiæ judicium et consensum secuta, hæc de ipso peccato originali statuit, fatetur ac declarat.}

enemy of mankind, amongst the very many evils with which the Church of God is in these our times troubled, has also stirred up not only new, but even old, dissensions touching original sin, and the remedy thereof; the sacred and holy, œcumenical and general Synod of Trent,—lawfully assembled in the Holy Ghost, the three same legates of the Apostolic See presiding therein,—wishing now to come to the reclaiming of the erring, and the confirming of the wavering,—following the testimonies of the sacred Scriptures, of the holy Fathers, of the most approved councils, and the judgment and consent of the Church itself, ordains, confesses, and declares these things touching the said original sin:

1. \emph{Si quis non confitetur, primum hominem Adam, cum mandatum Dei in paradiso fuisset transgressus, statim sanctitatem et justitiam, in qua constitutus fuerat, amisisse incurrisseque per offensam prævaricationis hujusmodi iram et indignationem Dei, atque ideo mortem, quam antea illi comminatus fuerat Deus, et cum morte captivitatem sub ejus potestate, qui mortis deinde habuit imperium, hoc est, diaboli, totumque Adam, per illam}

1. If any one does not confess that the first man, Adam, when he had transgressed the commandment of God in Paradise, immediately lost the holiness and justice wherein he had been constituted; and that he incurred, through the offense of that prevarication, the wrath and indignation of God, and consequently death, with which God had previously threatened him, and, together with death, captivity under his power who thenceforth \emph{had the empire of death, that is to say, the devil},\footnote{\par   Heb. ii. 14.}

\emph{prævaricationis offensam, secundum corpus et animam in deterius commutatum fuisse; anathema sit.}

and that the entire Adam, through that offense of prevarication was changed, in body and soul, for the worse; let him be anathema.

2. \emph{Si quis Adæ prævaricationem sibi soli, et non eius propagini asserit nocuisse; et acceptam a Deo sanctitatem et justitiam, quam perdidit, sibi soli et non nobis etiam eum perdidisse; aut inquinatum illum per inobedientiæ peccatum, mortem et pœnas corporis tantum in omne genus humanum transfudisse, non autem et peccatum, quod mors est animæ; anathema sit: cum contradicat apostolo dicenti: Per unum hominem peccatum intravit in mundum et per peccatum mors, et ita in omnes homines mors pertransiit, in quo omnes peccaverunt.}

2. If any one asserts, that the prevarication of Adam injured himself alone, and not his posterity; and that the holiness and justice, received of God, which he lost, he lost for himself alone, and not for us also; or that he, being defiled by the sin of disobedience, has only transfused death and pains of the body into the whole human race, but not sin also, which is the death of the soul; let him be anathema:—whereas he contradicts the apostle who says: \emph{By one man sin entered into the world, and by sin death, and so death passed upon all men, in whom all have sinned.}\footnote{\par   Rom. v. 12.}

3. \emph{Si quis hoc Adæ peccatum, quod origine unum est et propagatione, non imitatione transfusum omnibus, inest unicuique proprium, vel per humanæ naturæ vires, vel per aliud remedium asserit tolli, quam per meritum unius mediatoris Domini nostri Iesu Christi, qui nos Deo reconciliavit in sanguine suo, factus nobis justitia, sanctificatio et redemptio; aut negat}

3. If any one asserts, that this sin of Adam,—which in its origin is one, and being transfused into all by propagation, not by imitation, is in each one as his own,—is taken away either by the powers of human nature, or by any other remedy than the merit of the \emph{one mediator, our Lord Jesus Christ},\footnote{\par   1 Tim. ii. 5.} \emph{who hath reconciled us to God in his own blood, being made unto us justice, sanctification, and redemption;}\footnote{\par   1 Cor. i. 30.}

\emph{ipsum Christi Iesu meritum per baptismi sacramentum in forma ecclesiæ rite collatum, tam adultis quam parvulis applicari; anathema sit: quia non est aliud nomen sub cœlo datum hominibus, in quo oporteat nos salvos fieri. Unde illa vox: Ecce agnus Dei; ecce qui tollit peccata mundi; et illa: Quicumque baptizati estis, Christum induistis.}

or if he denies that the said merit of Jesus Christ is applied, both to adults and to infants, by the sacrament of baptism rightly administered in the form of the Church; let him be anathema: \emph{For there is no other name under heaven given to men, whereby we must be saved.}\footnote{\par   Acts iv. 2.} Whence that voice: \emph{Behold the lamb of God, behold him who taketh away the sins of the world;}\footnote{\par   John i. 29.} and that other: \emph{As many as have been baptized, have put on Christ.}\footnote{\par   Gal. iii. 27.}

4. \emph{Si quis parvulos recentes ab uteris matrum baptizandos negat, etiam si fuerint a baptizatis parentibus orti; aut dicit in remissionem quidem peccatorum eos baptizari, sed nihil ex Adam trahere originalis peccati, quod regenerationis lavacro necesse sit expiari ad vitam æternam consequendam; unde fit consequens, ut in eis forma baptismatis in remissionem peccatorum non vera, sed falsa intelligatur; anathema sit; quoniam non aliter intelligendum est id, quod dixit apostolus: Per unum hominem peccatum intravit in mundum, et per peccatum mors, et ita in omnes homines mors pertransiit, in quo omnes peccaverunt, nisi quemadmodum ecclesia }

4. If any one denies, that infants, newly born from their mothers' wombs, even though they be sprung from baptized parents, are to be baptized; or says that they are \emph{baptized} indeed \emph{for the remission of sins},\footnote{\par   Acts ii. 38.} but that they derive nothing of original sin from Adam, which has need of being expiated by the laver of regeneration for obtaining life everlasting,—whence it follows as a consequence, that in them the form of baptism, \emph{for the remission of sins}, is understood to be not true, but false,—let him be anathema. For that which the apostle has said, \emph{By one man sin entered into the world, and by sin death, and so death passed upon all men, in whom all have sinned},\footnote{\par   Rom. v. 12.} is not to be understood otherwise

\emph{catholica ubique diffusa semper intellexit. Propter hanc enim regulam fidei ex traditione apostolorum etiam parvuli, qui nihil peccatorum in semetipsis adhuc committere potuerunt, ideo in remissionem peccatorum veraciter baptizantur, ut in eis regeneratione mundetur, quod generatione contraxerunt. Nisi enim quis renatus fuerit ex aqua et Spiritu Sancto, non potest introire in regnum Dei.}

than as the Catholic Church spread every where hath always understood it. For, by reason of this rule of faith, from a tradition of the apostles, even infants, who could not as yet commit any sin of themselves, are for this cause truly \emph{baptized for the remission of sins}, that in them that may be cleansed away by regeneration, which they have contracted by generation. \emph{For, unless a man be born again of water and the Holy Ghost, he can not enter into the kingdom of God.}\footnote{\par   John iii. 5.}

5. \emph{Si quis per Iesu Christi Domini nostri gratiam, quæ in baptismate confertur, reatum originalis peccati remitti negat; aut etiam asserit non tolli totum id quod veram et propriam peccati rationem habet; sed illud dicit tantum radi, aut non imputari; anathema sit. In renatis enim nihil odit Deus; quia nihil est damnationis iis, qui vere consepulti sunt cum Christo per baptisma in mortem; qui non secundum carnem ambulant, sed veterem hominem exuentes, et novum, qui secundum Deum creatus est, induentes, innocentes, immaculati, puri, innoxii, ac Deo dilecti effecti sunt, heredes quidem Dei, coheredes autem}

5. If any one denies, that, by the grace of our Lord Jesus Christ, which is conferred in baptism, the guilt of original sin is remitted; or even asserts that the whole of that which has the true and proper nature of sin is not taken away; but says that it is only rased, or not imputed; let him be anathema. For, in those who are \emph{born again}, there is nothing that God hates; because, \emph{There is no condemnation to those who are} truly \emph{buried together with Christ by baptism into death;}\footnote{\par   Rom. viii. 1;  vi. 4} \emph{who walk not according to the flesh}, but, \emph{putting off the old man, and putting on the new who is created according to God},\footnote{\par   Ephes. iv. 22, 24.} are made innocent, immaculate, pure, harmless, and beloved of God, \emph{heirs}

\emph{Christi; ita ut nihil prorsus eos ab ingressu cœli remoretur. Manere autem in baptizatis concupiscentiam vel fomitem, hæc sancta synodus fatetur et sentit: quæ cum ad agonem relicta sit, nocere non consentientibus, sed viriliter per Christi Iesu gratiam repugnantibus non valet: quinimmo qui legitime certaverit, coronabitur. Hanc concupiscentiam, quam aliquando apostolus peccatum appellat, sancta synodus declarat, ecclesiam catholicam nunquam intellexisse peccatum appellari, quod vere et proprie in renatis peccatum sit, sed quia ex peccato est et ad peccatum inclinat. Si quis autem contrarium senserit, anathema sit.}

\emph{indeed of God, but joint heirs with Christ;}\footnote{\par   Rom. viii.17.} so that there is nothing whatever to retard their entrance into heaven. But this holy synod confesses and is sensible, that in the baptized there remains concupiscence, or an incentive (to sin); which, whereas it is left for our exercise, can not injure those who consent not, but resist manfully by the grace of Jesus Christ; yea, he who shall have \emph{striven lawfully shall be crowned.}\footnote{\par   2 Tim. ii. 5.} This concupiscence, which the apostle sometimes calls sin,\footnote{\par   Rom. vi. 12;  vii. 8.} the holy Synod declares that the Catholic Church has never understood it to be called sin, as being truly and properly sin in those \emph{born again}, but because it is of sin, and inclines to sin. And if any one is of a contrary sentiment, let him be anathema.

\emph{Declarat tamen hæc ipsa sancta synodus, non esse suæ intentionis comprehendere in hoc decreto, ubi de peccato originali agitur, beatam et immaculatam virginem Mariam, Dei genitricem; sed observandas esse constitutiones felicis recordationis Sixti papæ IV. sub pænis in eis constitutionibus contentis, quas innovat.} \footnote{\par [This indirect exemption of the \emph{ immaculata Virgo Maria }from original sin is a very near approach to the positive definition of the \emph{ immaculata conceptio Virginis Mariæ }in 1854.—P. S.]}

This same holy Synod doth nevertheless declare, that it is not its intention to include in this decree, where original sin is treated of, the blessed and immaculate Virgin Mary, the mother of God; but that the constitutions of Pope Sixtus IV., of happy memory, are to be observed, under the pains contained in the said constitutions, which it renews.

\xchapter{Sixth Session, held January 13, 1547.}

Sessio Sexta,

Sixth Session,

\emph{celebrata die XIII. Januarii} 1547.

\emph{held January} 13, 1547.

DECRETUM DE JUSTIFICATIONE.

DECREE ON JUSTIFICATION.

Caput I.

Chapter I.

\emph{De naturæ et legis ad justificandos homines imbecillitate.}

\emph{On the Inability of Nature and of the Law to justify Man.}

\emph{Primum declarat sancta synodus, ad justificationis doctrinam probe et sincere intelligendam oportere, ut unusquisque agnoscat et fateatur, quod cum omnes homines in prævaricatione Adæ innocentiam perdidissent; facti immundi et ut apostolus inquit, natura filii iræ, quemadmodum in decreto de peccato originali exposuit, usque adeo servi erant peccati et sub potestate diaboli ac mortis, ut non modo gentes per vim naturæ, sed ne Iudæi quidem per ipsam etiam litteram legis Moysi, inde liberari aut surgere possent; tametsi in eis liberum arbitrium minime extinctum esset, viribus licet attenuatum et inclinatum.}

The holy Synod declares first, that, for the correct and sound understanding of the doctrine of Justification, it is necessary that each one recognize and confess, that, whereas all men had lost their innocence in the prevarication of Adam,—having become unclean,\footnote{\par   Isa. lxiv. 6.} and as the apostle says, \emph{by nature children of wrath},\footnote{\par   Ephes. ii. 3.} as (this Synod) has set forth in the decree on original sin,—they were so far \emph{the servants of sin},\footnote{\par   Rom. vi. 17, 20.} and under the power of the devil and of death, that not the Gentiles only by the force of nature, but not even the Jews by the very letter itself of the law of Moses, were able to be liberated, or to arise, therefrom; although freewill, attenuated as it was in its powers, and bent down, was by no means extinguished in them.

Caput II.

Chapter II.

\emph{De dispensatione et mysterio Adventus Christi.}

\emph{On the Dispensation and Mystery of Christ's Advent.}

\emph{Quo factum est, ut cœlestis Pater, Pater misericordiarum, et Deus totius consolationis},

Whence it came to pass, that the heavenly Father, \emph{the Father of mercies, and the God of all comfort},\footnote{\par   2Cor. i. 3.}

\emph{Christum Iesum, Filium suum, et ante legem et legis tempore multis sanctis patribus declaratum ac promissum, cum venit beata illa plenitudo temporis, ad homines miserit, ut et Iudæos, qui sub lege erant, redimeret, et gentes, quæ non sectabantur justitiam, justitiam apprehenderent, atque omnes adoptionem filiorum reciperent. Hunc proposuit Deus propitiatorem per fidem in sanguine ipsius pro peccatis nostris, non solum autem pro nostris, sed etiam pro totius mundi.}

when that blessed \emph{fullness of the time was come},\footnote{\par   Gal. iv. 4.} \emph{sent} unto men, Jesus Christ, \emph{his own Son}—who had been, both before the Law, and during the time of the Law, to many of the holy fathers announced and promised—\emph{that he might both redeem the Jews who were under the Law},\footnote{\par   Gal. v. 4.} and that \emph{the Gentiles, who followed not after justice}, might attain to justice,\footnote{\par   Rom. ix. 30.} and that all men might receive the adoption of sons. Him God hath \emph{proposed} as a propitiator, \emph{through faith in his blood},\footnote{\par   Rom. iii. 25.} \emph{for our sins, and not for our sins only, but also for those of the whole world.}\footnote{\par   1 John ii. 2.}

Caput III.

Chapter III.

\emph{Qui per Christum justificantur.}

\emph{Who are justified through Christ.}

\emph{Verum, etsi ille pro omnibus mortuus est, non omnes tamen mortis ejus beneficium recipiunt; sed ii dumtaxat, quibus meritum passionis ejus communicatur. Nam, sicut re vera homines, nisi ex semine Adæ propagati nascerentur, non nascerentur injusti; cum ea propagatione, per ipsum dum concipiuntur, propriam injustitiam contrahant: ita, nisi in Christo renascerentur, nunquam justificarentur; cum ea renascentia per meritum passionis ejus gratia},

But, though \emph{He died for all},\footnote{\par   2 Cor. v. 15.} yet do not all receive the benefit of his death, but those only unto whom the merit of his passion is communicated. For as in truth men, if they were not born propagated of the seed of Adam, would not be born unjust,—seeing that, by that propagation, they contract through him, when they are conceived, injustice as their own,—so, if they were not born again in Christ, they never would be justified; seeing that, in that new birth, there is bestowed upon them, through the

\emph{qua justi fiunt, illis tribuatur. Pro hoc beneficio apostolus gratias nos semper agere hortatur Patri, qui dignos nos fecit in partem sortis sanctorum in lumine, et eripuit de potestate tenebrarum, transtulitque in regnum Filii dilectionis suæ, in quo habemus redemptionem et remissionem peccatorum.}

merit of his passion, the grace whereby they are made just. For this benefit, the apostle exhorts us, evermore \emph{to give thanks to the Father, who hath made us worthy to be partakers of the lot of the saints in light, and hath delivered us from the power of darkness, and hath translated us into the Kingdom of the Son of his love, in whom we have redemption, and remission of sins.}\footnote{\par   Coloss. i. 12-14.}

Caput IV.

Chapter IV.

\emph{Insinuatur descriptio justifactionis impii, et modus ejus in statu gratiæ.}

\emph{A description is introduced of the Justification of the impious, and of the manner thereof in the state of grace.}

\emph{Quibus verbis justifications impii descriptio insinuatur, ut sit translatio ab eo statu, in quo homo nascitur filius primi Adæ, in statum gratiæ, et adoptionis filiorum Dei per secundum Adam Iesum Christum, salvatorem nostrum: quæ quidem translatio post evangelium promulgatum, sine lavacro regenerationis, aut ejus voto, fieri non potest; sicut scriptum est: Nisi quis renatus fuerit ex aqua et Spiritu Sancto, non potest introire in regnum Dei.}

By which words, a description of the Justification of the impious is indicated,—as being a translation, from that state wherein man is born a child of the first Adam, to the state of grace, and of \emph{the adoption of the sons of God},\footnote{\par   Rom. viii. 15, 16, 23.} through the second Adam, Jesus Christ, our Saviour. And this translation, since the promulgation of the Gospel, can not be effected, without the laver of regeneration, or the desire thereof, as it is written: \emph{unless a man be born again of water and the Holy Ghost, he can not enter into the Kingdom of God.}\footnote{\par   John iii. 5.}

Caput V.

Chapter V.

\emph{De necessitate præparationis ad justificationem in adultis, et unde sit.}

\emph{On the necessity, in adults, of preparation for Justification, and whence it proceeds.}

\emph{Declarat præterea, ipsius justificationis exordium in adultis a Dei per Christum Iesum præveniente gratia sumendum esse, hoc est, ab ejus vocatione, qua, nullis eorum existentibus meritis, vocantur; ut, qui per peccata a Deo aversi erant, per ejus excitantem atque adjuvantem gratiam ad convertendum se ad suam ipsorum justiftcationem, eidem gratiæ libere assentiendo et cooperando, disponantur: ita ut, tangente Deo cor hominis per Spiritus Sancti illuminationem, neque homo ipse nihil omnino agat, inspirationem illam recipiens, quippe qui illam et abjicere potest, neque tamen sine gratia Dei movere se ad justitiam coram illo libera sua voluntate possit. Unde in sacris litteris cum dicitur: Convertimini ad me, et ego convertar ad vos: libertatis nostræ admonemur. Cum respondemus: Converte nos, Domine, ad te, et convertemur: Dei nos gratia præveniri confitemur.}

The Synod furthermore declares, that, in adults, the beginning of the said Justification is to be derived from the prevenient grace of God, through Jesus Christ, that is to say, from his vocation, whereby, without any merits existing on their parts, they are called; that so they, who by sins were alienated from God, may be disposed through his quickening and assisting grace, to convert themselves to their own justification, by freely assenting to and co-operating with that said grace: in such sort that, while God touches the heart of man by the illumination of the Holy Ghost, neither is man himself utterly inactive while he receives that inspiration, forasmuch as he is also able to reject it; yet is he not able, by his own free will, without the grace of God, to move himself unto justice in his sight. Whence, when it is said in the sacred writings: \emph{Turn ye to me, and I will turn to you},\footnote{\par   Zach. i. 3.} we are admonished of our liberty; and when we answer: \emph{Convert us, O Lord, to thee, and we shall be converted},\footnote{\par   Lam. v. 21.} we confess that we are prevented (anticipated) by the grace of God.

Caput VI.

Chapter VI.

\emph{Modus præparationis.}

\emph{The manner of Preparation.}

\emph{Disponuntur autem ad ipsam justitiam, dum excitati divina gratia et adjuti, fidem ex auditu concipientes, libere moventur in Deum, credentes vera esse, quæ divinitus revelata et promissa sunt; atque illud in primis, a Deo justificari impium per gratiam ejus), per redemptionem, quæ est in Christo Iesu: et, dum peccatores se esse intelligentes, a divinæ justitiæ timore, quo utiliter concutiuntur, ad considerandam Dei misericordiam se convertendo, in spem eriguntur, fidentes Deum sibi propter Christum propitium fore; illumque, tamquam omnis justitiæ fontem diligere incipiunt; ac propterea moventur adversus peccata per odium aliquod et detestationem, hoc est, per eam pœnitentiam, quam ante baptismum agi oportet: denique dum proponunt suscipere baptismum inchoare novam vitam, et servare divina mandata. De hac dispositione scriptum est: Accedentem ad Deum oportet credere, quia est, et quod inquirentibus se remunerator sit: et, Confide, fili, remittuntur tibi peccata tua; et: Timor }

Now they [adults] are disposed unto the said justice, when, excited and assisted by divine grace, conceiving \emph{faith by hearing},\footnote{\par   Rom. x. 17.} they are freely moved towards God, believing those things to be true which God has revealed and promised—and this especially, that God justifies the impious \emph{by his grace, through the redemption that is in Christ Jesus};\footnote{\par   Rom. iii. 24.} and when, understanding themselves to be sinners, they, by turning themselves, from the fear of divine justice whereby they are profitably agitated, to consider the mercy of God, are raised unto hope, confiding that God will be propitious to them for Christ's sake; and they begin to love him as the fountain of all justice; and are therefore moved against sins by a certain hatred and detestation, to wit, by that penitence which must be performed before baptism: lastly, when they purpose to receive baptism, to begin a new life, and to \emph{keep the commandments} of God. Concerning this disposition it is written: \emph{He that cometh to God, must believe that he is, and is a rewarder to them that seek him};\footnote{\par   Heb. xi. 6.} and, \emph{Be of good faith, son, thy sins}

\emph{Domini expellit peccatum; et: Pœnitentiam agite, et baptizetur unusquisque vestrum in nomine Iesu Christi, in remissionem peccatorum vestrorum, et accipietis donum Spiritus Sancti; et: Euntes ergo docete omnes gentes, baptizantes eos in nomine Patris, et Filii et Spiritus Sancti, docentes eos servare quæcumque mandavi vobis; denique: Præparate corda vestra Domino.}

\emph{are forgiven thee};\footnote{\par   Matt. ii. 5.} and, \emph{The fear of the Lord driveth out sin};\footnote{\par   Eccles. i. 27.} and, \emph{Do penance, and be baptized every one of you in the name of Jesus Christ, for the remission of your sins, and you shall receive the gift of the Holy Ghost};\footnote{\par   Acts ii. 38.} and, \emph{Going, therefore, teach ye all nations, baptizing them in the name of the father, and of the Son, and of the Holy Ghost};\footnote{\par   Matt. xxviii. 19.} finally, \emph{Prepare your hearts unto the Lord.}\footnote{\par   1 Kings vii. 3.}

Caput VII.

Chapter VII.

\emph{Quid sit justificatio impii, et quæ ejus causæ.}

\emph{What the Justification of the impious is, and what the causes thereof.}

\emph{Hanc dispositionem, seu præparationem justificatio ipsa consequitur, quæ non est sola peccatorum remissio, sed et sanctificatio et renovatio interioris hominis per voluntariam susceptionem gratiæ et donorum, unde homo ex injusto fit justus, et ex inimico amicus, ut sit heres secundum spem vitæ, æternæ.}

This disposition, or preparation, is followed by Justification itself, which is not remission of sins merely, but also the sanctification and renewal of the inward man, through the voluntary reception of the grace, and of the gifts, whereby man of unjust becomes just, and of an enemy a friend, that so he may be \emph{an heir according to hope of life everlasting.}\footnote{\par   Titus iii. 7}

\emph{Hujus justificationis causæ sunt, finalis quidem: gloria Dei et Christi, ac vita æterna; efficiens vero; misericors Deus, qui gratuito abluit, et sanctificat signans, et ungens Spiritu promissionis Sancto, qui est pignus }

Of this Justification the causes are these: the final cause indeed is the glory of God and of Jesus Christ, and life everlasting; while the efficient cause is a merciful God who \emph{washes and sanctifies}\footnote{\par   1 Cor. vi. 11} gratuitously, \emph{signing}, and anointing with the

\emph{hereditatis nostræ; meritoria autem: dilectissimus unigenitus suus, Dominus noster Iesus Christus, qui cum essemus inimici, propter nimiam caritatem, qua dilexit nos, sua sanctissima passione in ligno crucis nobis justificationem meruit, et pro nobis Deo Patri satisfecit; instrumentalis item: sacramentum baptismi, quod est sacramentum fidei, sine qua nulli umquam contigit justificatio; demum unica formalis causa est justitia Dei; non qua ipse justus est, sed qua nos justos facit; qua videlicet ab eo donati, renovamur spiritu mentis nostræ, et non modo reputamur, sed vere justi nominamur et sumus, justitiam in nobis recipientes, unusquisque suam secundum mensuram, quam Spiritus Sanctus partitur singulis prout vult et secundum propriam cujusque dispositionem et cooperationem. Quamquam enim nemo possit esse justus, nisi cui merita passionis Domini nostri Iesu Christi communicantur: id tamen in hac impii justificatione fit, dum ejusdem sanctissimæ passionis merito per Spiritum Sanctum caritas Dei diffunditur in cordibus }

holy \emph{Spirit of promise, who is the pledge of our inheritance};\footnote{\par   Ephes. i. 13, 14.} but the meritorious cause is his most beloved only-begotten, our Lord Jesus Christ, who, when we were enemies, \emph{for the exceeding charity wherewith he loved us},\footnote{\par   Ephes. ii. 4.} merited Justification for us by his most holy Passion on the wood of the cross, and made satisfaction for us unto God the Father; the instrumental cause is the sacrament of baptism, which is the sacrament of faith, without which [faith] no man was ever justified;\footnote{\par   Heb. xi.} lastly, the alone formal cause is the justice of God, not that whereby he himself is just, but that whereby he maketh us just, that, to wit, with which \emph{we}, being endowed by him, \emph{are renewed in the spirit of our mind},\footnote{\par   Ephes. iv. 23} and we are not only reputed, but are truly called, and are just, receiving justice within us, each one according to his own measure, \emph{which the Holy Ghost distributes to every one as he wills},\footnote{\par   1 Cor. xii. 2.} and according to each one's proper disposition and co-operation. For, although no one can be just, but he to whom the merits of the Passion of our Lord Jesus Christ are communicated, yet is this done in the said justification of the impious, when by the merit of that same

\emph{eorum, qui justificantur, atque ipsis inhæret: unde in ipsa justificatione cum remissione peccatorum hæc omnia simul infusa accipit homo per Iesum Christum, cui inseritur, fidem, spem et caritatem: nam fides, nisi ad eam spes accedat, et caritas, neque unit perfecte cum Christo, neque corporis ejus vivum membrum efficit: qua ratione verissime dicitur, fidem sine operibus mortuam, et otiosam esse: et in Christo Iesu neque circumcisionem aliquid valere neque præputium, sed fidem, quæ per caritatem operatur. Hanc fidem ante baptismi sacramentum ex apostolorum traditione catechumeni ab ecclesia petunt, cum petunt fidem, vitam æternam præstantem: quam sine spe et caritate præstare fides non potest: unde et statim verbum Christi audiunt: Si vis ad vitam ingredi, serva mandata.}

most holy Passion, \emph{the charity of God is poured forth}, by the Holy Spirit, \emph{in the hearts}\footnote{\par   Rom. v. 5.} of those that are justified, and is inherent therein: whence, man, through Jesus Christ, in whom he is ingrafted, receives, in the said justification, together with the remission of sins, all these [gifts] infused at once, faith, hope, and charity. For faith, unless hope and charity be added thereto, neither unites man perfectly with Christ, nor makes him a living member of his body. For which reason it is most truly said, \emph{that Faith without works is dead} and profitless;\footnote{\par   James ii. 20.} and, \emph{In Christ Jesus neither circumcision availeth any thing nor uncircumcision, but faith which worketh by charity.}\footnote{\par   Gal. v. 6.} This faith, Catechumens beg of the Church—agreeably to a tradition of the apostles—previously to the sacrament of Baptism; when they beg for the faith which bestows life everlasting, which, without hope and charity, faith can not bestow: whence also do they immediately hear that word of Christ: \emph{If thou wilt enter into life, keep the commandments.}\footnote{\par   Matt. xix. 17.}

\emph{Itaque veram et Christianam justitiam accipientes, eam ceu primam stolam pro illa, quam Adam sua inobedienta sibi et nobis perdidit, }

Wherefore, when receiving true and Christian justice, they are bidden, immediately on being born again, to preserve it pure and spotless,

\emph{per Christum Iesum illis donatam, candidam et immaculatam jubentur statim renati conservare, ut eam perferant ante tribunal Domini nostri Iesu Christi, et habeant vitam æternam.}

as \emph{the first robe}\footnote{\par   Luke xv. 22.} given them through Jesus Christ in lieu of that which Adam, by his disobedience, lost for himself and for us, that so they may bear it before the judgment-seat of our Lord Jesus Christ, and may have life eternal.

Caput VIII.

Chapter VIII.

\emph{Quomodo intelligatur, impium per fidem et gratis justificari.}

\emph{In what manner it is to be understood, that the impious is justified by faith, and gratuitously.}

\emph{Cum vero Apostolus dicit, justificari hominem per fidem et gratis, ea verba in eo sensu intelligenda sunt, quem perpetuus ecclesiæ catholicæ consensus tenuit et expressit: ut scilicet per fidem ideo justificari dicamur, quia fides est humanæ salutis initium, fundamentum et radix omnis justificationis, sine qua impossibile est placere Deo et ad filiorum ejus consortium pervenire: gratis autem justificari ideo dicamur, quia nihil eorum, quæ, justificationem præcedunt, sive fides sive opera, ipsam justificationis gratiam promeretur: si enim gratia est, jam non ex operibus: alioquin, ut idem apostolus inquit, gratia jam non est gratia.}

And whereas the Apostle saith, that man is \emph{justified by faith} and \emph{freely},\footnote{\par   Rom. iii. 4.} those words are to be understood in that sense which the perpetual consent of the Catholic Church hath held and expressed; to wit, that we are therefore said to be \emph{justified by faith}, because faith is the beginning of human salvation, the foundation, and the root of all Justification; \emph{without which it is impossible to please God},\footnote{\par   Heb. xi. 6.} and to come unto the fellowship of his sons: but we are therefore said to be justified \emph{freely}, because that none of those things which precede justification—whether faith or works—merit the grace itself of justification. For, \emph{if it be a grace, it is not now by works}, otherwise, as the same Apostle says, \emph{grace is no more grace.}\footnote{\par   Rom. xi. 6.}

Caput IX.

Chapter IX.

\emph{Contra inanem hæreticorum fiduciam.}

\emph{Against the vain confidence of heretics.}

\emph{Quamvis autem necessarium sit credere, neque remitti, neque remissa unquam fuisse peccata, nisi gratis divina misericordia propter Christum: nemini tamen fiduciam, et certitudinem remissionis peccatorum suorum jactanti, et in ea sola quiescenti, peccata dimitti, vel dimissa esse dicendum est, cum apud hæreticos et schismaticos possit esse, imo nostra tempestate sit, et magna contra ecclesiam catholicam contentione prædicetur vana hæc et ab omni pietate remota fiducia. Sed neque illud asserendum est, oportere eos, qui vere justificati sunt, absque ulla omnino dubitatione apud semetipsos statuere, se esse justificatos, neminemque a peccatis absolvi ac justificari, nisi eum, qui certo credat se absolutum et justificatum esse; atque hac sola fide absolutionem et justificationem perfici, quasi qui hoc non credit, de Dei promissis, deque mortis et resurrectionis Christi efficacia dubitet. Nam, sicut nemo pius de Dei misericordia, de Christi merito deque sacramentorum virtute et efficacia dubitare debet: sic quilibet, dum se ipsum suamque propriam infirmitatem et }

But, although it is necessary to believe that sins neither are remitted, nor ever were remitted save gratuitously by the mercy of God for Christ's sake; yet is it not to be said, that sins are forgiven, or have been forgiven, to any one who boasts of his confidence and certainty of the remission of his sins, and rests on that alone; seeing that it may exist, yea does in our day exist, amongst heretics and schismatics; and with great vehemence is this vain confidence, and one alien from all godliness, preached up in opposition to the Catholic Church. But neither is this to be asserted—that they who are truly justified must needs, without any doubting whatever, settle within themselves that they are justified, and that no one is absolved from sins and justified, but he that believes for certain that he is absolved and justified; and that absolution and justification are effected by this faith alone: as though whoso has not this belief, doubts of the promises of God, and of the efficacy of the death and resurrection of Christ. For even as no pious person ought to doubt of the mercy of God, of the merit of Christ, and of the virtue and efficacy of the sacraments, even so

\emph{indispositionem respicit, de sua gratia formidare et timere potest; cum nullus scire valeat certitudine fidei, cui non potest subesse falsum, se gratiam Dei esse consecutum.}

each one, when he regards himself, and his own weakness and indisposition, may have fear and apprehension touching his own grace; seeing that no one can know with a certainty of faith, which can not be subject to error, that he has obtained the grace of God.

Caput X.

Chapter X.

\emph{De acceptæ justificationis incremento.}

\emph{On the increase of Justification received.}

\emph{Sic ergo justificati, et amici Dei ac domestici facti, euntes de virtute in virtutem, renovantur, ut apostolus inquit, de die in diem, hoc est, mortificando membra carnis suæ, et exhibendo ea arma justitiæ in sanctificationem; per observationem mandatorum Dei et ecclesiæ, in ipsa justitia per Christi gratiam accepta, cooperante fide bonis operibus, crescunt atque magis justificantur, sicut scriptum est: Qui justus est, justificetur adhuc; et iterum: Ne verearis usque ad mortem justificari; et rursus: Videtis, quoniam ex operibus justificatur homo, et non ex fide tantum. Hoc vero justitiæ incrementum petit sancta ecclesia, cum orat: Da }

Having, therefore, been thus justified, and made the friends and \emph{domestics of God},\footnote{\par   Ephes. ii. 19.} advancing \emph{from virtue to virtue},\footnote{\par   Psa. lxxxiii. 8.} they are \emph{renewed}, as the Apostle says, \emph{day by day};\footnote{\par   2 Cor. iv. 16.} that is, \emph{by mortifying the members} of their own flesh,\footnote{\par   Col. iii. 5.} and \emph{by presenting them as instruments of justice unto sanctification},\footnote{\par   Rom. vi. 13, 19.} they, through the observance of the commandments of God and of the Church, faith co-operating with good works, increase in that justice which they have received through the grace of Christ, and are still further justified, as it is written: \emph{He that is just, let him be justified still;}\footnote{\par   Apoc. xxii. 11.} and again, \emph{Be not afraid to be justified even to death;}\footnote{\par   Eccles. xviii. 22.} and also, \emph{Do you see that by works a man is justified, and not by faith only.}\footnote{\par   James ii. 24.} And this increase

\emph{nobis Domine fidei, spei, et caritatis augmentum.}

of justification holy Church begs, when she prays, 'Give unto us, O Lord, increase of faith, hope, and charity.'

Caput XI.

Chapter XI.

\emph{De observatione mandatorum, deque illius necessitate et possibilitate.}

\emph{On keeping the Commandments, and on the necessity and possibility thereof.}

\emph{Nemo autem, quantumvis justificatus, liberum se esse ab observatione mandatorum putare debet; nemo temeraria illa et a patribus sub anathemate prohibita voce uti, Dei præcepta homini justificato ad observandum esse impossibilia. Nam Deus impossibilia non jubet, sed jubendo monet et facere quod possis, et petere quod non possis, et adjuvat, ut possis. Cujus mandata gravia non sunt, cujus jugum suave est et onus leve. Qui enim sunt filii Dei, Christum diligunt; qui autem diligunt eum, ut ipsemet testatur, servant sermones ejus, quod utique cum divino auxilio præstare possunt. Licet enim in hac mortali vita quantumvis sancti et justi in levia saltem et quotidiana, quæ etiam venialia dicuntur, peccata quandoque cadant, non propterea desinunt esse justi; nam justorum illa vox est et humilis et verax: Dimitte }

But no one, how much soever justified, ought to think himself exempt from the observance of the commandments; no one ought to make use of that rash saying, one prohibited by the Fathers under an anathema,—that the observance of the commandments of God is impossible for one that is justified. For God commands not impossibilities, but, by commanding, both admonishes thee to do what thon art able, and to pray for what thou art not able (to do), and aids thee that thou mayest be able; \emph{whose commandments are not heavy;}\footnote{\par   1 John v. 3.} \emph{whose yoke is sweet and whose burthen light.}\footnote{\par   Matt. xi. 30.} For, whoso are the sons of God, love Christ; but \emph{they who love him, keep his commandments},\footnote{\par   John xiv. 15.} as himself testifies; which, assuredly, with the divine help, they can do. For, although, during this mortal life, men, how holy and just soever, at times fall into at least light and daily sins, which are also called venial, not therefore do they cease to

\emph{nobis debita nostra. Quo fit, ut justi ipsi eo magis se obligatos ad ambulandum in via justitiæ sentire debeant, quo liberati jam a peccato, servi autem facti Deo, sobrie, juste et pie viventes proficere possint per Christum Iesum, per quem accessum habuerunt in gratiam istam. Deus namque sua gratia semel justificatos non deserit, nisi ab eis prius deseratur. Itaque nemo sibi in sola fide blandiri debet, putans fide sola se heredem esse constitutum, hereditatemque consecuturum, etiam si Christo non compatiatur, ut et conglorificetur. Nam et Christus ipse, ut inquit apostolus, cum esset filius Dei, didicit ex iis, quæ passus est, obedientiam, et consummatus factus est omnibus obtemperantibus sibi causa salutis æternæ. Propterea apostolus ipse monet justificatos, dicens: Nescitis, quod ii, qui in stadio currunt, omnes quidem currunt, sed unus accipit bravium? Sic currite, ut comprehendatis. Ego igitur sic curro, non quasi in incertum, sic pugno, non quasi aërem verberans, sed castigo corpus meum, et in servitutem redigo, ne forte, }

be just. For that cry of the just, \emph{Forgive us our trespasses}, is both humble and true. And for this cause, the just themselves ought to feel themselves the more obliged to walk in the way of justice, in that, \emph{being} already \emph{freed from sins, but made servants of God},\footnote{\par   Rom. vi. 18.} they are able, \emph{living soberly, justly, and godly},\footnote{\par   Titus ii. 12.} to proceed onwards \emph{through Jesus Christ, by whom they have had access unto this grace.}\footnote{\par   Rom. v. 2.} For God forsakes not those who have been once justified by his grace, unless he be first forsaken by them. Wherefore, no one ought to flatter himself up with faith alone, fancying that by faith alone he is made an heir, and will obtain the inheritance, even though \emph{he suffer} not \emph{with Christ, that} so \emph{he may be also glorified with him.}\footnote{\par   Rom. viii. 17.} For even Christ himself, as the Apostle saith, \emph{Whereas he was the son of God, learned obedience by the things which he suffered, and being consummated, he became, to all who obey him, the cause of eternal salvation.}\footnote{\par  Heb. v. 8, 9.} For which cause the same Apostle admonishes the justified, saying: \emph{Know you not that they that run in the race, all run indeed, but one receiveth the prize? So run that you may obtain. I therefore so}

\emph{cum aliis prædicaverim, ipse reprobus efficiar. Item princeps apostolorum Petrus: Satagite, ut per bona opera certam vestram vocationem et electionem faciatis. Hæc enim facientes, non peccabitis aliquando. Unde constat eos orthodoxæ religionis doctrinæ adversari, qui dicunt, justum in omni bono opere saltem venialiter peccare, aut, quod intolerabilius est, pœnas æternas mereri, atque etiam eos, qui statuunt, in omnibus operibus justos peccare, si in illis suam ipsorum socordiam excitando, et sese ad currendum in stadio cohortando, cum hoc, ut in primis glorificetur Deus, mercedem quoque intuentur æternam; cum scriptum sit: Inclinavi cor meum ad faciendas justificationes tuas propter retributionem; et de Mose dicat apostolus, quod respiciebat in remunerationem.}

\emph{run, not as at an uncertainty: I so fight, not as one beating the air, but I chastise my body, and bring it into subjection; lest, perhaps, when I have preached to others, I myself should become a cast-away.}\footnote{\par  1 Cor. ix. 24, 26, 27.} So also the prince of the Apostles, Peter: \emph{Labor the more that by good works you may make sure your calling and election. For doing those things, you shall not sin at any time.}\footnote{\par   2 Peter i. 10.} From which it is plain, that those are opposed to the orthodox doctrine of religion, who assert that the just man sins, venially at least, in every good work; or, which is yet more insupportable, that he merits eternal punishments; as also those who state, that the just sin in all their works, if, in those works, they, together with this aim principally that God may be glorified, have in view also the eternal reward, in order to excite their sloth, and to encourage themselves \emph{to run in the course:} whereas it is writen, \emph{I have inclined my heart to do all thy justifications for the reward:}\footnote{\par   Psa. cxviii. 112.} and, concerning Moses, the Apostle saith, that \emph{he looked unto the reward.}\footnote{\par   Heb. xi. 26.}

Caput XII.

Chapter XII.

\emph{Prædestinationis temerariam præsumptionem cavendam esse.}

\emph{That a rash presumptuousness in the matter of Predestination is to be avoided.}

\emph{Nemo quoque, quamdiu in hoc mortalitate vivitur, de arcano divinæ prædestinationis mysterio usque adeo præsumere debet, ut certo statuat, se omnino esse in numero prædestinatorum, quasi verum esset, quod justificatus aut amplius peccare non possit, aut, si peccaverit, certam sibi resipiscentiam promittere debeat. Nam, nisi ex speciali revelatione, sciri non potest, quos Deus sibi elegerit.}

No one, moreover, so long as he is in this mortal life, ought so far to presume as regards the secret mystery of divine predestination, as to determine for certain that he is assuredly in the number of the predestinate; as if it were true, that he that is justified, either can not sin any more, or, if he do sin, that he ought to promise himself an assured repentance; for except by special revelation, it can not be known whom God hath chosen unto himself.

Caput XIII.

Chapter XIII.

\emph{De perseverantiæ munere.}

\emph{On the gift of Perseverance.}

\emph{Similiter de perseverantiæ munere, de quo scriptum est: Qui perseveraverit usque in finem, hic salvus erit; quod quidem aliunde haberi non potest, nisi ab eo, qui potens est eum, qui stat, statuere, ut perseveranter stet, et eum, qui cadit, restituere: nemo sibi certi aliquid absoluta certitudine polliceatur, tametsi in Dei auxilio firmissimam spem collocare et reponere omnes debent. Deus enim, nisi ipsi illius gratiæ defuerint, sicut cœpit opus bonum, ita perficiet, operans }

So also as regards the gift of perseverance, of which it is written, \emph{He that shall persevere to the end, he shall be saved;}\footnote{\par   Matt. xxiv. 13.} —which gift can not be derived from any other but Him, who is able to establish him who standeth\footnote{\par   Rom. xiv. 4.} that he stand perseveringly, and to restore him who falleth:—let no one herein promise himself any thing as certain with an absolute certainty; though all ought to place and repose a most firm hope in God's help. For God, unless men be themselves wanting in his grace, \emph{as he has begun the}

\emph{velle et perficere. Verumtamen, qui se existimant stare, videant ne cadant et cum timore, ac tremore salutem suam operentur in laboribus, in vigiliis, in eleemosynis, in orationibus et oblationibus, in jejuniis et castitate; formidare enim debent, scientes quod in spem gloriæ, et nondum in gloriam renati sunt, de pugna, quæ superest cum carne, cum mundo, cum diabolo; in qua victores esse non possunt, nisi cum Dei gratia apostolo obtemperent, dicenti: Debitores sumus non carni, ut secundum carnem vivamus; si enim secundum carnem vixeritis, moriemini; si autem spiritu facta carnis mortificaveritis, vivetis.}

\emph{good work, so will he perfect it, working} (in them) \emph{to will and to accomplish.}\footnote{\par  \par   Phil. i. 6; ii. 13.} Nevertheless, let those who \emph{think themselves to stand, take heed lest they fall},\footnote{\par   1Cor. x. 12.} and, \emph{with fear and trembling work out their salvation},\footnote{\par   Phil. ii. 12.} in labors, in watchings, in almsdeeds, in prayers and oblations, in fastings and chastity: for, knowing that \emph{they are born again unto a hope of glory},\footnote{\par   1 Peter i. 3.} but not as yet unto glory, they ought to fear for the combat which yet remains with the flesh, with the world, with the devil, wherein they can not be victorious, unless they be with God's grace, obedient to the Apostle, who says: \emph{We are debtors, not to the flesh, to live according to the flesh; for if you live according to the flesh, you shall die; but if by the spirit you mortify the deeds of the flesh, you shall live.}\footnote{\par   Rom. viii. 12, 13.}

Caput XIV.

Chapter XIV.

\emph{De lapsis, et eorum reparatione.}

\emph{On the fallen, and their restoration.}

\emph{Qui vero ab accepta justificationis gratia per peccatum exciderunt, rursus justificari poterunt, cum, excitante Deo, per pœnitentiæ sacramentum merito Christi amissam gratiam recuperare procuraverint; hic enim justificationis modus est lapsi }

As regards those who, by sin, have fallen from the received grace of Justification, they may be again justified, when, God exciting them, through the sacrament of Penance they shall have attained to the recovery, by the merit of Christ, of the grace lost: for this manner of

\emph{reparatio, quam secundam post naufragium deperditæ gratiæ tabulam sancti patres apte nucuparunt; etenim pro iis, qui post baptismum in peccata labuntur, Christus Iesus sacramentum instituit pœnitentiæ, cum dixit: Accipite Spiritum Sanctum: quorum remiseritis peccata, remittuntur eis; et quorum retinueritis, retenta sunt. Unde docendum est, Christiani hominis pœnitentiam post lapsum multo aliam esse a baptismali, eaque contineri non modo cessationem a peccatis, et eorum detestationem, aut cor contritum et humiliatum, verum etiam eorundem sacramentalem confessionem saltem in voto et suo tempore faciendam, et sacerdotalem absolutionem; itemque satisfactionem per jejunia, eleemosynas, orationes et alia pia spiritualis vitæ exercitia; non quidem pro pœna æterna, quæ vel sacramento, vel sacramenti voto una cum culpa remittitur; sed pro pœna temporali, quæ, ut sacræ litteræ docent, non tota semper, ut in baptismo fit, dimittitur illis, qui gratiæ Dei, quam acceperunt, ingrati, Spiritum Sanctum, contristaverunt, et templum Dei violare non sunt veriti. De qua }

Justification is of the fallen the reparation: which the holy Fathers have aptly called a second plank after the shipwreck of grace lost. For, on behalf of those who fall into sins after baptism, Christ Jesus instituted the sacrament of Penance, when he said, \emph{Receive ye the Holy Ghost, whose sins you shall forgive, they are forgiven them, and whose sins you shall retain, they are retained}.\footnote{\par   John xx. 22, 23.} Whence it is to be taught, that the penitence of a Christian, after his fall, is very different from that at (his) baptism; and that therein are included not only a cessation from sins, and a detestation thereof, or, \emph{a contrite and humble heart},\footnote{\par   Psa. l. 19.} but also the sacramental confession of the said sins,—at least in desire, and to be made in its season,—and sacerdotal absolution; and likewise satisfaction by fasts, alms, prayers, and the other pious exercises of a spiritual life; not indeed for the eternal punishment,—which is, together with the guilt, remitted, either by the sacrament, or by the desire of the sacrament,—but for the temporal punishment, which, as the sacred writings teach, is not always wholly remitted, as is done in baptism, to those who, ungrateful to the grace of God which they have received,

\emph{pœnitentia scriptum est: Memor esto, unde excideris, age pœnitentiam, et prima opera fac. Et iterum: quæ secundum Deum tristitia est, pœnitentiam in salutem stabilem operatur. Et rursus: Pœnitentiam agite, et facite fructus dignos pœnitentiæ.}

have \emph{grieved the Holy Spirit},\footnote{\par   Ephes. iv. 30.} and have not feared \emph{to violate the temple of God.}\footnote{\par   1 Cor. iii. 17.} Concerning which penitence it is written: \emph{Be mindful whence thou art fallen; do penance, and do the first works.}\footnote{\par   Apoc. ii. 5.} And again: \emph{The sorrow that is according to God worketh penance steadfast unto salvation.}\footnote{\par   2 Cor. vii. 10.} And again: \emph{Do penance}, and \emph{bring forth fruits worthy of penance.}\footnote{\par   Matt. iii. 2.}

Caput XV.

Chapter XV.

\emph{Quolibet mortali peccato amitti gratiam sed non fidem.}

\emph{That, by every mortal sin, grace is lost, but not faith.}

\emph{Adversus etiam hominum quorundam callida ingenia, qui per dulces sermones et benedictiones seducunt corda innocentium, asserendum est, non modo infidelitate, per quam et ipsa fides amittitur, sed etiam quocumque alio mortali peccato, quamvis non amittatur fides, acceptam justificationis gratiam amitti; divinæ legis doctrinam defendendo, quæ a regno Dei non solum infideles excludit, sed et fideles quoque, fornicarios, adulteros, molles, masculorum concubitores, fures, avaros, ebriosos, maledicos, rapaces, ceterosque omnes, qui letalia committunt peccata, a quibus cum divinæ gratiæ }

In opposition also to the subtle wits of certain men, who, \emph{by pleasing speeches and good words, seduce the hearts, of the innocent},\footnote{\par   Rom. xvi. 18.} it is to be maintained, that the received grace of Justification is lost, not only by infidelity whereby even faith itself is lost, but also by any other mortal sin whatever, though faith be not lost; thus defending the doctrine of the divine law, which excludes from the kingdom of God not only the unbelieving, but the faithful also [who are] \emph{fornicators, adulterers, effeminate, liers with mankind, thieves, covetous, drunkards, railers, extortioners},\footnote{\par  1 Cor. vi. 9, 10.} and all others who commit deadly sins; from which, with the

\emph{adiumento abstinere possunt, et pro quibus a Christi gratia separantur.}

help of divine grace, they can refrain, and on account of which they are separated from the grace of Christ.

Caput XVI.

Chapter XVI.

\emph{De fructu justificationis, hoc est, de merito bonorum operum, deque ipsius meriti ratione.}

\emph{On the fruit of Justification, that is, on the merit of good works, and on the nature of that merit.}

\emph{Hac igitur ratione justificatis hominibus, sive acceptam gratiam perpetuo conservaverint, sive amissam recuperaverint, proponenda sunt apostoli verba: Abundate in omni opere bono, scientes, quod labor vester non est inanis in Domino; non enim injustus est Deus, ut obliviscatur operis vestri et dilectionis, quam ostendistis in nomine ipsius; et: Nolite amittere confidentiam vestram, quæ magnam habet remunerationem. Atque ideo bene operantibus usque in finem, et in Deo sperantibus proponenda est vita æterna, et tanquam gratia filiis Dei per Christum Iesum misericorditer promissa, et tanquam merces ex ipsius Dei promissione bonis ipsorum operibus et mentis fideliter reddenda. Hæc est enim illa corona justitiæ, quam post suum certamen et cursum repositam sibi esse aiebat apostolus, a justo }

Before men, therefore, who have been justified in this manner,—whether they have preserved uninterruptedly the grace received, or whether they have recovered it when lost,—are to be set the words of the Apostle: \emph{Abound in every good work, knowing that your labor is not in vain in the Lord};\footnote{\par   1 Cor. xv. 58.} \emph{for God is not unjust, that he should forget your work, and the love which you have shown in his name};\footnote{\par   Heb. vi. 10.} and, \emph{do not lose your confidence, which hath, a great reward.}\footnote{\par   Heb. x. 35.} And, for this cause, life eternal is to be proposed to those working well \emph{unto the end},\footnote{\par   Matt. x. 22.} and hoping in God, both as a grace mercifully promised to the sons of God through Jesus Christ, and as a reward which is according to the promise of God himself, to be faithfully rendered to their good works and merits. For this is that \emph{crown of justice} which the Apostle declared was, after his \emph{fight} and \emph{course, laid up}

\emph{judice sibi reddendam; non solum autem sibi, sed et omnibus, qui diligunt adventum ejus: cum enim ille ipse Christus Iesus, tanquam caput in membra et tanquam vitis in palmites, in ipsos justificatos jugiter virtutem influat, quæ virtus bona eorum opera semper antecedit et comitatur et subsequitur, et sine qua nullo pacto Deo grata, et meritoria esse possent: nihil ipsis justificatis amplius deesse credendum est, quo minus plene illis quidem operibus, quæ in Deo sunt facta, divinæ legi pro hujus vitæ statu satisfecisse, et vitam æternam suo etiam tempore (si tamen in gratia decesserint), consequendam, vere promeruisse censeantur, cum Christus, Salvator noster, dicat: Si (quis biberit ex aqua, quam ego dabo ei, non sitiet in æternum, sed fiet in eo fons aquæ salientis in vitam æternam }

\emph{for him, to be rendered to him by the just Judge, and not only to him, but also to all that love his coming.}\footnote{\par   2 Tim. iv. 8.} For, whereas Jesus Christ himself continually infuses his virtue into the said justified,—as the head into the members, and the vine into the branches,—and this virtue always precedes and accompanies and follows their good works, which without it could not in any wise be pleasing and meritorious before God,—we must believe that nothing further is wanting to the justified, to prevent their being accounted to have, by those very works which have been done in God, fully satisfied the divine law according to the state of this life, and to have truly merited eternal life, to be obtained also in its (due) time, if so be, however, that they depart in grace: seeing that Christ, our Saviour, saith: \emph{If any one shall drink of the water that I will give him, he shall not thirst forever; but it shall become in him, a fountain of water springing up unto life everlasting.}\footnote{\par   John iv. 13, 14.}

\emph{Ita neque propria nostra justitia, tanquam ex nobis propria statuitur, neque ignoratur aut repudiatur justitia Dei; quæ enim justitia nostra dicitur, quia per eam nobis inhærentem justificamur, illa eadem }

Thus, neither is our own justice \emph{established as our own} as from ourselves;\footnote{\par   Rom. x. 3.} nor is the justice of God ignored or repudiated: for that justice which is called ours, because that we are justified from its being inherent in us, that same is (the justice) of

\emph{Dei est, quia a Deo nobis infunditur per Christi meritum. Neque vero illud omittendum est, quod licet bonis operibus in sacris litteris usque adeo tribuatur, ut etiam qui uni ex minimis suis potum aquæ frigidæ dederit, promittat Christus eum non esse sua mercede cariturum, et apostolus testetur, id quod in præsenti est momentaneum et leve tribulationis nostræ, supra modum in sublimitate æternum glorias pondus operari in nobis: absit tamen, ut Christianus homo in se ipso vel confidat vel glorietur, et non in Domino, cujus tanta est erga omnes homines bonitas, ut eorum velit esse merita, quæ sunt ipsius dona. Et quia in multis offendimus omnes, unusquisque, sicut misericordiam et bonitatem, ita severitatem et judicium ante oculos habere debet, neque se ipsum aliquis, etiam si nihil sibi conscius fuerit, judicare; quoniam omnis hominum vita non humano judido examinanda et judicanda est, sed Dei, qui illuminabit abscondita tenebrarum, et manifestabit consilia cordium: et tunc laus erit unicuique a Deo, qui, ut scriptum }

God, because that it is infused into us of God, through the merit of Christ. Neither is this to be omitted,—that although, in the sacred writings, so much is attributed to good works, that Christ promises, that even \emph{he that shall give a drink of cold water to one of his least ones, shall not lose his reward;}\footnote{\par   Matt. x. 42.}\emph{ }and the Apostle testifies that, \emph{That which is at present momentary and light of our tribulation, worketh for us above measure exceedingly an eternal weight of glory};\footnote{\par   2 Cor. iv. 17.} nevertheless God forbid that a Christian should either trust or glory in himself, and not in the Lord, whose bounty towards all men is so great, that he will have the things which are his own gifts be their merits. And forasmuch as \emph{in many things we all offend},\footnote{\par   James iii. 2.} each one ought to have before his eyes, as well the severity and judgment, as the mercy and goodness (of God); neither ought any one \emph{to judge himself, even though he be not conscious to himself of any thing};\footnote{\par   1 Cor. iv. 3, 4.} because the whole life of man is to be examined and judged, not by the judgment of man, but of God, \emph{who will bring to light the hidden things of darkness, and will make manifest the counsels of the hearts, and then shall every man}

\emph{est, reddet unicuique opera sua.}

\emph{have praise from God},\footnote{\par   1 Cor. iv. 6.} who, as it is written, \emph{will render to every man according to his works.}\footnote{\par   Matt. xvi. 27.}

\emph{Post hanc catholicam de justificatione doctrinam, quam nisi quisque fideliter firmiterque receperit, justificari non poterit, placuit sanctæ synodo hos canones subjungere, ut omnes sciant, non solum quid tenere et sequi, sed etiam quid vitare et fugere debeant.}

After this Catholic doctrine on Justification, which whoso receiveth not faithfully and firmly can not be justified, it hath seemed good to the holy Synod to subjoin these canons, that all may know not only what they ought to hold and follow, but also what to avoid and shun.

DE JUSTIFICATIONE.

ON JUSTIFICATION.

Canon I.—\emph{Si quis dixerit, hominem suis operibus, quæ vel per humanæ naturæ vires, vel per legis doctrinam fiant, absque divina per Iesum Christum gratia posse justificari coram Deo: anathema sit.}

CANON I.—If any one saith, that man may be justified before God by his own works, whether done through the teaching of human nature, or that of the law, without the grace of God through Jesus Christ: let him be anathema.

Canon II.—\emph{Si quis dixerit, ad hoc solum divinam gratiam per Christum Iesum dari, ut facilius homo juste vivere, ac vitam æternam promereri possit; quasi per liberum arbitrium sine gratia utrumque, sed ægre tamen et difficulter possit: anathema, sit.}

CANON II.—If any one saith, that the grace of God, through Jesus Christ, is given only for this, that man may be able more easily to live justly, and to merit eternal life, as if, by free-will without grace, he were able to do both, though hardly indeed and with difficulty: let him be anathema.

Canon III.—\emph{Si quis dixerit, sine præveniente Spiritus Sancti inspiratione atque ejus adjutorio hominem credere, sperare, diligere, aut pænitere posse, sicut}

CANON III.—If any one saith, that without the prevenient inspiration of the Holy Ghost, and without his help, man can believe, hope, love, or be penitent as he

\emph{oportet, ut ei justificationis gratia conferatur: anathema sit.}

ought, so that the grace of Justification may be bestowed upon him: let him be anathema.

Canon IV.—\emph{Si quis dixerit, liberum hominis arbitrium a Deo motum et excitatum nihil cooperari assentiendo Deo excitanti atque vocanti, quo ad obtinendam justificationis gratiam se disponat ac præparet; neque posse dissentire, si velit, sed veluti inanime quoddam nihil omnino agere, mereque passive se habere: anathema sit.}

CANON IV.—If any one saith, that man's free-will moved and excited by God, by assenting to God exciting and calling, nowise co-coperates towards disposing and preparing itself for obtaining the grace of Justification; that it can not refuse its consent, if it would, but that, as something inanimate, it does nothing whatever and is merely passive: let him be anathema.

Canon V.—\emph{Si quis liberum hominis arbitrium post Adæ peccatum amissum et extinctum esse dixerit, aut rem esse de solo titulo, imo titulum sine re, figmentum denique a Satana invectum in ecclesiam: anathema sit.}

CANON V.—If any one saith, that, since Adam's sin, the free-will of man is lost and extinguished; or, that it is a thing with only a name, yea a name without a reality, a figment, in fine, introduced into the Church by Satan: let him be anathema.

Canon VI.—\emph{Si quis dixerit, non esse in potestate hominis, vias suas malas facere, sed mala opera ita, ut bona, Deum operari, non permissive solum, sed etiam proprie et per se, adeo ut sit proprium ejus opus non minus proditio Iudæ, quam vocatio Pauli: anathema sit.}

CANON VI.—If any one saith, that it is not in man's power to make his ways evil, but that the works that are evil God worketh as well as those that are good, not permissively only, but properly, and of himself, in such wise that the treason of Judas is no less his own proper work than the vocation of Paul: let him be anathema.

Canon VII.—\emph{Si quis dixerit, opera omnia, quæ ante justiftcationem fiunt, quacumque ratione}

CANON VII.—If any one saith, that all works done before Justification, in whatsoever way they be

\emph{facta sint, vere esse peccata, vel odium Dei mereri, aut, quanto vehementius quis nititur se disponere ad gratiam, tanto eum gravius peccare: anathema sit.}

done, are truly sins, or merit the hatred of God; or that the more earnestly one strives to dispose himself for grace, the more grievously he sins: let him be anathema.

Canon VIII.—\emph{Si quis dixerit, gehennæ metum, per quem ad misericordiam Dei de peccatis dolendo confugimus vel a peccando abstinemus, peccatum esse, aut peccatores peiores facere: anathema sit.}

CANON VIII.—If any one saith, that the fear of hell,—whereby, by grieving for our sins, we flee unto the mercy of God, or refrain from sinning,—is a sin, or makes sinners worse: let him be anathema.

Canon IX.—\emph{Si quis dixerit, sola fide impium justificari, ita ut intelligat nihil aliud requiri, quod ad justificationis gratiam consequendam cooperetur, et nulla ex parte necesse esse, eum suæ voluntatis motu præparari atque disponi: anathema sit.}

CANON IX.—If any one saith, that by faith alone the impious is justified, in such wise as to mean, that nothing else is required to cooperate in order to the obtaining the grace of Justification, and that it is not in any way necessary, that he be prepared and disposed by the movement of his own will: let him be anathema.

Canon X.—\emph{Si quis dixerit, homines sine Christi justitia, per quam nobis meruit, justificari, aut per eam ipsam formaliter justos esse: anathema sit.}

CANON X.—If any one saith, that men are just without the justice of Christ, whereby he merited for us to be justified; or that it is by that justice itself that they are formally just: let him be anathema.

Canon XI.—\emph{Si quis dixerit, homines justificari, vel sola imputatione justitiæ Christi, vel sola peccatorum remissione, exclusa gratia et caritate, quæ in cordibus eorum per Spiritum Sanctum diffundatur atque illis}

CANON XI.—If any one saith, that men are justified, either by the sole imputation of the justice of Christ, or by the sole remission of sins, to the exclusion of the grace and \emph{the charity which is poured forth in their hearts by the Holy Ghost},\footnote{\par   Rom. v. 5.}

\emph{inhæreat; aut etiam gratiam, qua justificamur, esse tantum favorem Dei: anathema sit.}

and is inherent in them; or even that the grace, whereby we are justified, is only the favor of God: let him be anathema.

Canon XII.—\emph{Si quis dixerit, fidem justificantem nihil aliud esse, quam fiduciam divinæ misericordiæ peccata remittentis propter Christum; vel eam fiduciam solam esse, qua justificamur: anathema sit.}

CANON XII.—If any one saith, that justifying faith is nothing else but confidence in the divine mercy which remits sins for Christ's sake; or, that this confidence alone is that whereby we are justified: let him be anathema.

Canon XIII.—\emph{Si quis dixerit, omni homini ad remissionem peccatorum assequendam necessarium esse, ut credat certo, et absque ulla hæsitatione propriæ infirmitatis et indispositionis peccata sibi esse remissa: anathema sit.}

CANON XIII.—If any one saith, that it is necessary for every one, for the obtaining the remission of sins, that he believe for certain, and without any wavering arising from his own infirmity and indisposition, that his sins are forgiven him: let him be anathema.

Canon XIV.—\emph{Si quis dixerit, hominem a peccatis absolvi ac justificari ex eo quod se absolvi ac justificari certo credat; aut neminem vero esse justificatum, nisi qui credat se esse justificatum, et hac sola fide absolutionem et justiftcationem perfici: anathema sit.}

CANON XIV.—If any one saith, that man is truly absolved from his sins and justified, because that he assuredly believed himself absolved and justified; or, that no one is truly justified but he who believes himself justified; and that, by this faith alone, absolution and justification are effected: let him be anathema.

Canon XV.—\emph{Si quis dixerit, hominem renatum et justificatum teneri ex fide ad credendum, se certo esse in numero prædestinatorum: anathema sit.}

CANON XV.—If any one saith, that a man, who is born again and justified, is bound of faith to believe that he is assuredly in the number of the predestinate: let him be anathema.

Canon XVI.—\emph{Si quis magnum}

CANON XVI.—If any one saith,

\emph{illud usque in finem perseverantiæ donum se certo habiturum absoluta et infallibili certitudine dixerit, nisi hoc ex speciali revelatione didicerit: anathema sit.}

that he will for certain, of an absolute and infallible certainty, have that great gift of perseverance unto the end,—unless he have learned this by special revelation: let him be anathema.

Canon XVII.—\emph{Si quis justificationis gratiam non nisi prædestinatis ad vitam contingere dixerit, reliquos vero omnes, qui vocantur, vocari quidem, sed gratiam non accipere, utpote divina potestate prædestinatos ad malum: anathema sit.}

CANON XVII.—If any one saith, that the grace of Justification is only attained to by those who are predestined unto life; but that all others who are called, are called indeed, but receive not grace, as being, by the divine power, predestined unto evil: let him be anathema.

Canon XVIII.—\emph{Si quis dixerit, Dei præcepta homini etiam justificato et sub gratia constitute esse ad observandum impossibilia: anathema sit.}

CANON XVIII.—If any one saith, that the commandments of God are, even for one that is justified and constituted in grace, impossible to keep: let him be anathema.

Canon XIX.—\emph{Si quis dixerit, nihil præceptum esse in evangelio præter fidem, cetera esse indifferentia, neque præcepta, neque prohibita, sed libera; aut decem præcepta nihil pertinere ad Christianas: anathema sit.}

CANON XIX.—If any one saith, that nothing besides faith is commanded in the Gospel; that other things are indifferent, neither commanded nor prohibited, but free; or, that the ten commandments nowise appertain to Christians: let him be anathema.

Canon XX.—\emph{Si quis hominem justificatum et quantumlibet perfectum dixerit non teneri ad observantiam mandatorum Dei et ecclesiæ, sed tantum ad credendum, quasi vero evangelium sit nuda et absoluta promissio vitæ æternæ sine conditione observationis }

CANON XX.—If any one saith, that the man who is justified and how perfect soever, is not bound to observe the commandments of God and of the Church, but only to believe; as if indeed the Gospel were a bare and absolute promise of eternal life, without the condition of

\emph{mandatorum: anathema sit.}

observing the commandments: let him be anathema.

Canon XXI.—\emph{Si quis dixerit. Christum Iesum a Deo hominibus datum fuisse, ut redemptorem, cui fidant, non etiam ut legislatorem, cui obediant: anathema sit.}

CANON XXI.—If any one saith, that Christ Jesus was given of God to men, as a redeemer in whom to trust, and not also as a legislator whom to obey: let him be anathema.

Canon XXII.—\emph{Si quis dixerit, justificatum, vel sine speciali auxilio Dei in accepta justitia perseverare posse, vel cum eo non posse: anathema sit.}

CANON XXII.—If any one saith, that the justified, either is able to persevere, without the special help of God, in the justice received; or that, with that help, he is not able: let him be anathema.

Canon XXIII.—\emph{Si quis hominem semel justificatum dixerit amplius peccare non posse, neque gratiam amittere, atque ideo eum qui labitur et peccat, nunquam vere fuisse justificatum; aut contra, posse in tota vita peccata omnia, etiam venialia, vitare, nisi ex speciali Dei privilegio, quemadmodum de beata Virgine tenet ecclesia: anathema sit.}

CANON XXIII.—If any one saith, that a man once justified can sin no more, nor lose grace, and that therefore he that falls and sins was never truly justified; or, on the other hand, that he is able, during his whole life, to avoid all sins, even those that are venial,—except by a special privilege from God, as the Church holds in regard of the Blessed Virgin: let him be anathema.

Canon XXIV.—\emph{Si quis dixerit, justitiam acceptam non conservari, atque etiam non augeri coram Deo per bona opera; sed opera ipsa fructus solummodo et signa esse justificationis adeptæ, non autem ipsius augendæ causam: anathema sit.}

CANON XXIV.—If any one saith, that the justice received is not preserved and also increased before God through good works; but that the said works are merely the fruits and signs of Justification obtained, but not a cause of the increase thereof: let him be anathema.

Canon XXV.—\emph{Si quis in quolibet bono opere justum saltem venialiter peccare dixerit, aut, }

CANON XXV.—If any one saith, that, in every good work, the just sins venially at least, or—which is

\emph{quod intolerabilius est, mortaliter, atque ideo pœnas æternas mereri; tantumque ob id non damnari, quia Deus opera non imputet ad damnationem: anathema sit.}

more intolerable still—mortally, and consequently deserves eternal punishments; and that for this cause only he is not damned, that God does not impute those works unto damnation: let him be anathema.

Canon XXVI.—\emph{Si quis dixerit, justos non debere pro bonis operibus, quæ in Deo fuerint facta, expectare et sperare æternam retributionem a Deo per ejus misericordiam et Iesu Christi meritum, si bene agendo et divina mandata custodiendo usque in finem perseveraverint: anathema sit.}

CANON XXVI.—If any one saith, that the just ought not, for their good works done in God, to expect and hope for an eternal recompense from God, through his mercy and the merit of Jesus Christ, if so be that they persevere to the end in well doing and in keeping the divine commandments: let him be anathema.

Canon XXVII.—\emph{Si quis dixerit, nullum esse mortale peccatum, nisi infidelitatis; aut nullo alio, quantumvis gravi et enormi, præterquam infidelitatis, peccato, semel acceptam gratiam amitti: anathema sit.}

CANON XXVII.—If any one saith, that there is no mortal sin but that of infidelity; or, that grace once received is not lost by any other sin, however grievous and enormous, save by that of infidelity: let him be anathema.

Canon XXVIII.—\emph{Si quis dixerit, amissa per peccatum gratia, simul et fidem semper amitti; aut fidem, quæ remanet, non esse veram fidem, licet non sit viva; aut eum, qui fidem sine caritate habet, non esse Christianum: anathema sit.}

CANON XXVIII.—If any one saith, that, grace being lost through sin, faith also is always lost with it; or, that the faith which remains, though it be not a lively faith, is not a true faith; or, that he who has faith without charity is not a Christian: let him be anathema.

Canon XXIX.—\emph{Si quis dixerit, eum, qui post baptismum lapsus est, non posse per Dei gratiam resurgere; aut posse quidem, sed sola fide amissam justitiam recuperare }

CANON XXIX.—If any one saith, that he who has fallen after baptism is not able by the grace of God to rise again; or, that he is able indeed to recover the justice

\emph{sine sacramento pœnitentiæ, prout sancta romana et universalis ecclesia a Christo Domino et ejus apostolis edocta hucusque professa est, servavit et docuit: anathema sit.}

which he has lost, but by faith alone without the sacrament of Penance, contrary to what the holy Roman and universal Church—instructed by Christ and his Apostles—has hitherto professed, observed, and taught: let him be anathema.

Canon XXX.—\emph{Si quis post acceptam justificationis gratiam cuilibet peccatori pœnitenti ita culpam remitti et reatum æternæ pænæ deleri dixerit, ut nullus remaneat reatus pænæ temporalis exsolvendæ vel in hoc seculo, vel in futuro in purgatorio, antequam ad regna cœlorum aditus patere possit: anathema sit.}

CANON XXX.—If any one saith, that, after the grace of Justification has been received, to every penitent sinner the guilt is remitted, and the debt of eternal punishment is blotted out in such wise that there remains not any debt of temporal punishment to be discharged either in this world, or in the next in Purgatory, before the entrance to the kingdom of heaven can be opened [to him]: let him be anathema.

Canon XXXI.—\emph{Si quis dixerit, justificatum peccare, dum intuitu æternæ mercedis bene operatur; anathema sit.}

CANON XXXI.—If any one saith, that the justified sins when he performs good works with a view to an eternal recompense: let him be anathema.

Canon XXXII.—\emph{Si quis dixerit hominis justificati bona opera ita esse dona Dei, ut non sint etiam bona ipsius justificati merita; aut ipsum justificatum bonis operibus, quæ ab eo per Dei gratiam, et Iesu Christi meritum, cujus vivum membrum est, fiunt, non vere mereri augmentum gratiæ, vitam æternam, et ipsius vitæ æternæ, si tamen in gratiæ decesserit, consecutionem},

CANON XXXII.—If any one saith, that the good works of one that is justified are in such manner the gifts of God, that they are not also the good merits of him that is justified; or, that the said justified, by the good works which he performs through the grace of God and the merit of Jesus Christ, whose living member he is, does not truly merit increase of grace, eternal life, and the attainment of that eternal

\emph{atque etiam gloriæ augmentum: anathema sit.}

life,—if so be, however, that he depart in grace,—and also an increase of glory: let him be anathema.

Canon XXXIII.—\emph{Si quis dixerit, per hanc doctrinam catholicam de justificatione, a sancta synodo hoc præsenti decreto expressam, aliqua ex parte gloriæ Dei vel meritis Iesu Christi Domini nostri derogari, et non potius veritatem fidei nostræ, Dei denique, ac Christi Iesu gloriam illustrari: anathema sit.}

CANON XXXIII.—If any one saith, that, by the Catholic doctrine touching Justification, by this holy Synod set forth in this present decree, the glory of God, or the merits of our Lord Jesus Christ are in any way derogated from, and not rather that the truth of our faith, and the glory in fine of God and of Jesus Christ are rendered [more] illustrious: let him be anathema.

\xchapter{Seventh Session, held March 3, 1547.}

Sessio Septima,

Seventh Session,

\emph{celebrata die III. Martii} 1547.

\emph{held March} 3, 1547.

DECRETUM DE SACRAMENTIS.

DECREE ON THE SACRAMENTS.

\emph{Proœmium.}

\emph{Proëm.}

\emph{Ad consummationem salutaris de justificatione doctrinæ, quæ in præcedenti proxima sessione uno omnium patrum consensu promulgata fuit; consentaneum visum est de sanctissimis ecclesiæ sacramentis agere, per quæ omnis vera justitia vel incipit, vel cœpta augetur, vel amissa reparatur. Propterea sacrosancta, œcumenica et generalis Tridentina synodus, in Spiritu Sancto legitime congregata, præsidentibus in ea eisdem }

For the completion of the salutary doctrine on Justification, which was promulgated with the unanimous consent of the Fathers in the last preceding Session, it hath seemed suitable to treat of the most holy Sacraments of the Church, through which all true justice either begins, or being begun is increased, or being lost is repaired. With this view, in order to destroy the errors and to extirpate the heresies which have appeared in these our days on the subject of the said most holy sacraments,—

\emph{apostolicæ sedis legatis, ad errores eliminandos et extirpandas hæreses, quæ circa sanctissima ipsa sacramenta hac nostra tempestate, tum de damnatis olim a patribus nostris hæresibus suscitatæ, tum etiam de novo adinventæ sunt, quæ Catholicæ Ecclesiæ puritati et animarum saluti magnopere officiunt; sanctarum scripturarum doctrinæ, apostolicis traditionibus atque aliorum conciliorum et patrum consensui inhærendo, hos præsentes canones statuendos et decernendos censuit, reliquos, qui supersunt ad cœpti operis perfectionem, deinceps, divino Spiritu adjuvante, editura.}

as well those which have been revived from the heresies condemned of old by our Fathers, as also those newly invented, and which are exceedingly prejudicial to the purity of the Catholic Church, and to the salvation of souls,—the sacred and holy, œcumenical and general Synod of Trent, lawfully assembled in the Holy Ghost, the same legates of the Apostolic See presiding therein, adhering to the doctrine of the holy Scriptures, to the apostolic traditions, and to the consent of other councils and of the Fathers, has thought fit that these present canons be established and decreed; intending, the divine Spirit aiding, to publish later the remaining canons which are wanting for the completion of the work which it has begun.

DE SACRAMENTIS IN GENERE.

ON THE SACRAMENTS IN GENERAL.

Canon I.—\emph{Si quis dixerit, sacramenta novæ legis non fuisse omnia a Iesu Christo Domino nostro instituta; aut esse plura vel pauciora quam septem, videlicet: baptismum, confirmationem, eucharistiam, pœnitentiam, extremam unctionem, ordinem, et matrimonium; aut etiam aliquod horum septem non esse vere et proprie sacramentum: anathema sit.}

Canon I.—If any one saith, that the sacraments of the New Law were not all instituted by Jesus Christ, our Lord; or, that they are more, or less, than seven, to wit, Baptism, Confirmation, the Eucharist, Penance, Extreme Unction, Order, and Matrimony; or even that any one of these seven is not truly and properly a sacrament: let him be anathema.

Canon II.—\emph{Si quis dixerit, ea ipsa novæ legis sacramenta a sacramentis antiquæ legis non differre, nisi quia ceremoniæ sunt aliæ et alii ritus externi: anathema, sit.}

Canon II.—If any one saith, that these said sacraments of the New Law do not differ from the sacraments of the Old Law, save that the ceremonies are different, and different the outward rites: let him be anathema.

Canon III.—\emph{Si quis dixerit, hæc septem sacramenta ita esse inter se paria, ut nulla ratione aliud sit alio dignius; anathema sit.}

Canon III.—If any one saith, that these seven sacraments are in such wise equal to each other, as that one is not in any way more worthy than another: let him be anathema.

Canon IV.—\emph{Si quis dixerit, sacramenta novæ legis non esse ad salutem necessaria, sed superflua; et sine eis aut eorum voto per solam fidem homines a Deo gratiam justificationis adipisci; licet omnia singulis necessaria non sint: anathema sit.}

Canon IV.—If any one saith, that the sacraments of the New Law are not necessary unto salvation, but superfluous; and that, without them, or without the desire thereof, men obtain of God, through faith alone, the grace of justification;—though all [the sacraments] are not indeed necessary for every individual: let him be anathema.

Canon V.—\emph{Si quis dixerit, hæc sacramenta propter solam fidem nutriendam instituta fuisse: anathema sit.}

Canon V.—If any one saith, that these sacraments were instituted for the sake of nourishing faith alone: let him be anathema.

Canon VI.—\emph{Si quis dixerit, sacramenta novæ legis non continere gratiam, quam significant; aut gratiam ipsam non ponentibus obicem non conferre; quasi signa tantum externa sint acceptæ per fidem gratiæ, vel justitiæ, et notæ quædam Christianæ professionis, quibus apud homines }

Canon VI.—If any one saith, that the sacraments of the New Law do not contain the grace which they signify; or, that they do not confer that grace on those who do not place an obstacle thereunto; as though they were merely outward signs of grace or justice received through faith, and certain marks of the Christian

\emph{discernuntur fideles ab infidelibus; anathema sit.}

profession, whereby believers are distinguished amongst men from unbelievers: let him be anathema.

Canon VII.—\emph{Si quis dixerit, non dari gratiam per hujusmodi sacramenta semper et omnibus, quantum est ex parte Dei, etiam si rite ea suscipiant, sed aliquando et aliquibus: anathema sit.}

Canon VII.—If any one saith, that grace, as far as God's part is concerned, is not given through the said sacraments, always, and to all men, even though they receive them rightly, but [only] sometimes, and to some persons: let him be anathema.

Canon VIII.—\emph{Si quis dixerit, per ipsa novæ legis sacramenta ex opere operato non conferri gratiam, sed solam fidem divinæ promissionis ad gratiam consequendam sufficere: anathema sit.}

Canon VIII.—If any one saith, that by the said sacraments of the New Law grace is not conferred through the act performed, but that faith alone in the divine promise suffices for the obtaining of grace: let him be anathema.

Canon IX.—\emph{Si quis dixerit, in tribus sacramentis, baptismo scilicet, confirmatione et ordine, non imprimi characterem in anima, hoc est signum quoddam spirituale et indelebile, unde ea iterari non possunt: anathema sit.}

Canon IX.—If any one saith, that, in the three sacraments, to wit, Baptism, Confirmation, and Order, there is not imprinted in the soul a character, that is, a certain spiritual and indelible sign, on account of which they can not be repeated: let him be anathema.

Canon X.—\emph{Si quis dixerit, Christianos omnes in verbo, et omnibus sacramentis administrandis habere potestatem: anathema sit.}

Canon X.—If any one saith, that all Christians have power to administer the word, and all the sacraments: let him be anathema.

Canon XI.—\emph{Si quis dixerit, in ministris, dum sacramenta conficiunt et conferunt, non requiri intentionem saltem faciendi, quod facit ecclesia: anathema sit.}

Canon XI.—If any one saith, that, in ministers, when they effect, and confer the sacraments, there is not required the intention at least of doing what the Church does: let him be anathema.

Canon XII.—\emph{Si quis dixerit,}

Canon XII.—If any one saith,

\emph{ministrum in peccato mortali existentem, modo omnia essentialia, quæ ad sacramentum conficiendum aut conferendum pertinent, servaverit, non conficere aut conferre sacramentum: anathema sit.}

that a minister, being in mortal sin,—if so be that he observe all the essentials which belong to the effecting, or conferring of, the sacrament,—neither effects, nor confers the sacrament: let him be anathema.

Canon XIII.—\emph{Si quis dixerit, receptos et approbates Ecclesiæ Catholicæ ritus, in solemni sacramentorum administratione adhiberi consuetos, aut contemni, aut sine peccato a ministris pro libito omitti, aut in novos alios per quemcumque ecclesiarum pastorem mutari posse: anathema sit.}

Canon XIII.—If any one saith, that the received and approved rites of the Catholic Church, wont to be used in the solemn administration of the sacraments, may be contemned, or without sin be omitted at pleasure by the ministers, or be changed, by every pastor of the churches, into other new ones: let him be anathema.

DE BAPTISMO.

ON BAPTISM.

Canon I.—\emph{Si quis dixerit, baptismum Ioannis habuisse eamdem vim cum baptismo Christi: anathema sit.}

Canon I.—If any one saith, that the baptism of John had the same force as the baptism of Christ: let him be anathema.

Canon II.—\emph{Si quis dixerit, aquam veram et naturalem non esse de necessitate baptismi; atque ideo verba illa Domini nostri Iesu Christi: Nisi quis renatus fuerit ex aqua et Spiritu Sancto; ad metaphoram aliquam detorserit: anathema sit.}

Canon II.—If any one saith, that true and natural water is not of necessity for baptism, and, on that account, wrests, to some sort of metaphor, those words of our Lord Jesus Christ: \emph{Unless a man be born again of water and the Holy Ghost:}\footnote{\par   John iii. 5.} let him be anathema.

Canon III.—\emph{Si quis dixerit, in Ecclesiæ Romana, quæ omnium ecclesiarum mater est et magistra, non esse veram de baptismi }

Canon III.—If any one saith, that in the Roman Church, which is the mother and mistress of all churches, there is not the true doctrine concerning

\emph{sacramento doctrinam: anathema sit.}

the sacrament of baptism: let him be anathema.

Canon IV.—\emph{Si quis dixerit, baptismum, qui etiam datur ab hæreticis in nomine Patris, et Filii, et Spiritus Sancti, cum intentione faciendi, quod facit ecclesia, non esse verum baptismum: anathema sit.}

Canon IV.—If any one saith, that the baptism which is even given by heretics in the name of the Father, and of the Son, and of the Holy Ghost, with the intention of doing what the Church doth, is not true baptism: let him be anathema.

Canon V.—\emph{Si quis dixerit, baptismum liberum esse, hoc est, non necessarium ad salutem: anathema, sit.}

Canon V.—If any one saith, that baptism is free, that is, not necessary unto salvation: let him be anathema.

Canon VI.—\emph{Si quis dixerit, baptizatum non posse, etiam si velit, gratiam amittere, quantumcumque peccet, nisi nolit credere: anathema sit.}

Canon VI.—If any one saith, that one who has been baptized can not, even if he would, lose grace, let him sin ever so much, unless he will not believe: let him be anathema.

Canon VII.—\emph{Si quis dixerit, baptizatos per baptismum ipsum, solius tantum fidei debitores fieri, non autem universæ legis Christi servandæ: anathema sit.}

Canon VII.—If any one saith, that the baptized are, by baptism itself, \emph{made debtors} but to faith alone, and not to the observance of \emph{the whole law}\footnote{\par  Gal. v. 3.} of Christ: let him be anathema.

Canon VIII.—\emph{Si quis dixerit, baptizatos liberos esse ab omnibus sanctæ ecclesiæ præceptis, quæ vel scripta vel tradita sunt, ita ut ea observare non teneantur, nisi se sua sponte illis submittere voluerint: anathema sit.}

Canon VIII.—If any one saith, that the baptized are freed from all the precepts, whether written or transmitted, of holy Church, in such wise that they are not bound to observe them, unless they have chosen of their own accord to submit themselves thereunto: let him be anathema.

Canon IX.—\emph{Si quis dixerit, ita revocandos esse homines ad }

Canon IX.—If any one saith, that the remembrance of the baptism

\emph{baptismi suscepti memoriam, ut vota omnia, quæ post baptismum fiunt, vi promissionis in baptismo ipso jam factæ, irrita esse intelligant, quasi per ea et fidei, quam professi sunt, detrahatur et ipsi baptismo: anathema sit.}

which they have received is so to be recalled unto men, as that they are to understand that all vows made after baptism are void, in virtue of the promise already made in that baptism; as if, by those vows, they both derogated from that faith which they have professed, and from that baptism itself: let him be anathema.

Canon X.—\emph{Si quis dixerit, peccata omnia, quæ post baptismum fiunt, sola recordatione et fide suscepti baptismi vel dimitti, vel venialia fieri: anathema sit.}

Canon X.—If any one saith, that by the sole remembrance and the faith of the baptism which has been received, all sins committed after baptism are either remitted, or made venial: let him be anathema.

Canon XI.—\emph{Si quis dixerit, verum et rite collatum baptismum iterandum esse illi, qui apud infideles fidem Christi negaverit, cum ad pœnitentiam convertitur: anathema sit.}

Canon XI.—If any one saith, that baptism, which was true and rightly conferred, is to be repeated, for him who has denied the faith of Christ amongst Infidels, when he is converted unto penitence: let him be anathema.

Canon XII.—\emph{Si quis dixerit, neminem esse baptizandum, nisi ea ætate, qua Christus baptizatus est, vel in ipso mortis articulo: anathema sit.}

Canon XII.—If any one saith, that no one is to be baptized save at that age at which Christ was baptized, or in the very article of death: let him be anathema.

Canon XIII.—\emph{Si quis dixerit, parvulos, eo quod actum credendi non habent, suscepto baptismo inter fideles computandos non esse, ac propterea, cum ad annos discretionis pervenerint, esse rebaptizandos; aut præstare, omitti eorum baptisma,}

Canon XIII.—If any one saith, that little children, for that they have not actual faith, are not, after having received baptism, to be reckoned amongst the faithful; and that, for this cause, they are to be rebaptized when they have attained to years of discretion; or, that it is

\emph{quam eos non acta, proprio credentes, baptizari in sola fide ecclesiæ: anathema sit.}

better that the baptism of such be omitted, than that, while not believing by their own act, they should be baptized in the faith alone of the Church: let him be anathema.

Canon XIV.—\emph{Si quis dixerit, hujusmodi parvulos baptizatos, cum adoleverint, interrogandos esse, an ratum habere velint, quod patrini eorum nomine, dum baptizarentur, polliciti sunt; et, ubi se nolle responderint, suo esse arbitrio relinquendos; nec alia interim pœna ad Christianam vitam cogendos, nisi ut ab Eucharistiæ aliorumque sacramentorum perceptione arceantur, donec resipiscant: anathema sit.}

Canon XIV.—If any one saith, that those who have been thus baptized when children, are, when they have grown up, to be asked whether they will ratify what their sponsors promised in their names when they were baptized; and that, in case they answer that they will not, they are to be left to their own will; and are not to be compelled meanwhile to a Christian life by any other penalty, save that they be excluded from the participation of the Eucharist, and of the other sacraments, until they repent: let him be anathema.

DE CONFIRMATIONE.

ON CONFIRMATION.

Canon I.—\emph{Si quis dixerit, confirmationem baptizatorum otiosam ceremoniam esse, et non potius verum et proprium sacramentum; aut olim nihil aliud fuisse, quam catechesim quamdam, qua adolescentiæ, proximi fidei suæ rationem coram ecclesia exponebant: anathema sit.}

Canon I.—If any one saith, that the confirmation of those who have been baptized is an idle ceremony, and not rather a true and proper sacrament; or that of old it was nothing more than a kind of catechism, whereby they who were near adolescence gave an account of their faith in the face of the Church: let him be anathema.

Canon II.—\emph{Si quis dixerit, injurios esse Spiritui Sancto eos, qui sacro confirmationis chrismati}

Canon II.—If any one saith, that they who ascribe any virtue to the sacred chrism of confirmation, offer

\emph{virtutem aliquam tribuunt: anathema sit. }

an outrage to the Holy Ghost: let him be anathema.

Canon III.—\emph{Si quis dixerit, sanctæ confirmationis ordinarium ministrum non esse solum episcopum, sed quemvis simplicem sacerdotem: anathema sit.}

Canon III.—If any one saith, that the ordinary minister of holy confirmation is not the bishop alone, but any simple priest soever: let him be anathema.

\xchapter{Thirteenth Session, held October 11, 1551.}

Sessio Decimatertia,

Thirteenth Session,

\emph{celebrata die XI. Octobris} 1551.

\emph{held October} 11, 1551.

DECRETUM DE SANCTISSIMO EUCHARISTIÆ.

DECREE CONCERNING THE MOST HOLY SACRAMENT OF THE EUCHARIST.

Caput I.

Chapter I.

\emph{De reali præsentiæ Domini nostri Iesu Christi in sanctissimo Eucharistiæ sacramento.}

On the real presence of our Lord Jesus Christ in the most holy sacrament of the Eucharist.

\emph{Principio docet sancta synodus, et aperte ac simpliciter profitetur, in almo sanctæ Eucharistiæ sacramento, post panis, et vini consecrationem, Dominum nostrum Iesum Christum, verum Deum atque hominem, vere, realiter, ac substantialiter sub specie illarum rerum sensibilium contineri. Neque enim hæc inter se pugnant, ut ipse Salvator noster semper ad dexteram Patris in cœlis assideat juxta modum existendi naturalem, et ut multis nihilominus aliis in locis sacramentaliter præsens sua substantia nobis adsit, ea existendi ratione, quam etsi verbis exprimere vix possumus, }

In the first place, the holy Synod teaches, and openly and simply professes, that, in the august sacrament of the holy Eucharist, after the consecration of the bread and wine, our Lord Jesus Christ, true God and man, is truly, really, and substantially contained under the species of those sensible things. For neither are these things mutually repugnant,—that our Saviour himself always sitteth at the right hand of the Father in heaven, according to the natural mode of existing, and that, nevertheless, he be, in many other places, sacramentally present to us in his own substance, by a manner of existing, which, though we can scarcely express it in words, yet

\emph{possibilem tamen esse Deo, cogitatione per fidem illustrata assequi possumus, et constantissime credere debemus: ita enim majores nostri omnes, quotquot in vera Christi ecclesia fuerunt, qui de sanctissimo hoc sacramento disseruerunt, apertissime professi sunt, hoc tam admirabile sacramentum in ultima cæna redemptorem nostrum instituisse, cum post panis vinique benedictionem se suum ipsius corpus illis præbere, ac suum sanguinem, disertis et perspicuis verbis testatus est; quæ verba a sanctis evangelistis commemorata et a divo Paulo postea repetita, cum propriam illam et apertissimam significationem præ se ferant, secundum quam a patribus intellecta sunt; indignissimum sane flagitium est, ea a quibusdam contentiosis et pravis hominibus ad fictitios et imaginarios tropos, quibus veritas carnis et sanguinis Christi negatur, contra universum ecclesiæ sensum detorqueri; quæ, tamquam columna et firmamentum veritatis, hæc ab impiis hominibus excogitata commenta velut satanica detestata est, grato semper et memore animo præstantissimum hoc Christi beneficium agnoscens.}

can we, by the understanding illuminated by faith, conceive, and we ought most firmly to believe, to be possible unto God: for thus all our forefathers, as many as were in the true Church of Christ, who have treated of this most holy Sacrament, have most openly professed, that our Redeemer instituted this so admirable a sacrament at the last supper, when, after the blessing of the bread and wine, he testified, in express and clear words, that he gave them his own very body, and his own blood, words which,—recorded by the holy Evangelists, and afterwards repeated, by Saint Paul, whereas they carry with them that proper and most manifest meaning in which they were understood by the Fathers,—it is indeed a crime the most unworthy that they should be wrested, by certain contentious and wicked men, to fictitious and imaginary tropes, whereby the verity of the flesh and blood of Christ is denied, contrary to the universal sense of the Church, which, as \emph{the pillar and ground of truth}, has detested, as satanical, these inventions devised by impious men; she recognizing, with a mind ever grateful and unforgetting, the most excellent benefit of Christ.

Caput II.

Chapter II.

\emph{De ratione institutionis sanctissimi hujus sacramenti.}

On the reason of the institution of this most holy sacrament.

\emph{Ergo Salvator noster, discessurus ex hoc mundo ad Patrem, sacramentum hoc instituit, in quo divitias divini sui erga homines amoris velut effudit, memoriam faciens mirabilium suorum; et in illius sumptione colere nos sui memoriam præcepit, suamque annunciare mortem, donec ipse ad judicandum mundum veniat. Sumi autem voluit sacramentum hoc, tamquam spiritualem animarum cibum, quo alantur, et confortentur viventes vita illius, qui dixit: Qui manducat me, et ipse vivet propter me: et tamquam antidotum, quo liberemur a culpis quotidianis, et a peccatis mortalibus præservemur. Pignus præterea id esse voluit futuræ nostræ gloriæ, et perpetuæ felicitatis, adeoque symbolum unius illius corporis, cujus ipse caput existit, cuique nos, tamquam membra, arctissima fidei, spei et caritatis connexione adstrictos esse voluit, ut idipsum omnes diceremus, nec essent in nobis schismata.}

Wherefore, our Saviour, when about to depart out of this world to the Father, instituted this sacrament, in which he poured forth as it were the riches of his divine love towards men, \emph{making a remembrance of his wonderful works;}\footnote{\par   Psa. cx. 4.} and he commanded us, in the participation thereof, to venerate his \emph{memory}, and to \emph{show forth his death until he come}\footnote{\par   1 Cor. xi. 26.} to judge the world. And he would also that this sacrament should be received as the spiritual food of souls, whereby may be fed and strengthened those who live with his life who said, \emph{He that eateth me, the same also shall live by me;}\footnote{\par   John vi. 58.} and as an antidote, whereby we may be freed from daily faults, and be preserved from mortal sins. He would, furthermore, have it be a pledge of our glory to come, and everlasting happiness, and thus be a symbol of that one body whereof he is the head, and to which he would fain have us as members be united by the closest bond of faith, hope, and charity, \emph{that we might all speak the same things, and there might be no schisms amongst us.}\footnote{\par   1 Cor. i. 10.}

Caput III.

Chapter III.

\emph{De excellentia sanctissimæ Eucharistiæ super reliqua sacramenta.}

\emph{On the excellency of the most holy Eucharist over the rest of the sacraments.}

\emph{Commune hoc quidem est sanctissimæ Eucharistiæ cum ceteris sacramentis, symbolum esse rei sacræ, et invisibilis gratiæ formam visibilem; verum illud in ea excellens et singulare reperitur, quod reliqua sacramenta tunc primum sanctificandi vim habent, cum quis illis utitur: at in Eucharistia ipse sanctitatis auctor ante usum est. Nondum enim Eucharistiam de manu Domini apostoli susceperant, cum vere tamen ipse affirmaret corpus suum esse, quod præbebat.}

The most holy Eucharist has indeed this in common with the rest of the sacraments, that it is a symbol of a sacred thing, and is a visible form of an invisible grace; but there is found in the Eucharist this excellent and peculiar thing, that the other sacraments have then first the power of sanctifying when one uses them, whereas in the Eucharist, before being used, there is the Author himself of sanctity. For the apostles had not as yet received the Eucharist from the hand of the Lord, when nevertheless himself affirmed with truth that to be his own body which he presented [to them].

\emph{Et semper hæc fides in Ecclesia Dei fuit, statim post consecrationem verum Domini nostri corpus verumque ejus sanguinem sub panis et vini specie una cum ipsius anima et divinitate existere; sed corpus quidem sub specie panis et sanguinem sub vini specie ex vi verborum; ipsum autem corpus sub specie vini, et sanguinem sub specie panis, animamque sub utraque, vi naturalis illius connexionis et concomitantiæ, qua partes Christi Domini, qui jam ex mortuis resurrexit non amplius }

And this faith has ever been in the Church of God, that, immediately after the consecration, the veritable body of our Lord, and his veritable blood, together with his soul and divinity, are under the species of bread and wine; but the body indeed under the species of bread, and the blood under the species of wine, by the force of the words; but the body itself under the species of wine, and the blood under the species of bread, and the soul under both, by the force of that natural connection and concomitancy whereby the parts of Christ

\emph{moriturus, inter se copulantur, divinitatem porro propter admirabilem illam ejus cum corpore et anima hypostaticam unionem. Quapropter verissimum est, tantumdem sub alterutra specie atque sub utraque contineri: totus enim, et integer Christus sub panis specie et sub quavis ipsius speciei parte, totus item sub vini specie et sub ejus partibus existit.}

our Lord, \emph{who hath now risen from the dead, to die no more},\footnote{\par   1 Cor. vi. 9.} are united together; and the divinity, furthermore, on account of the admirable hypostatical union thereof with his body and soul. Wherefore it is most true, that as much is contained under either species as under both; for Christ whole and entire is under the species of bread, and under any part whatsoever of that species; likewise the whole (Christ) is under the species of wine, and under the parts thereof.

Caput IV.

Chapter IV.

De Transsubstantiatione.

On Transubstantiation.

\emph{Quoniam autem Christus, redemptor noster, corpus suum id, quod sub specie panis offerebat, vere esse dixit; ideo persuasum semper in Ecclesia Dei fuit, idque nunc denuo sancta hæc synodus declarat, per consecrationem panis et vini conversionem fieri totius substantiæ panis in substantiam corporis Christi Domini nostri, et totius substantiæ vini in substantiam sanguinis ejus: quæ conversio convenienter et proprie a sancta Catholica Ecclesia Transsubstantiatio est appellata.}

And because that Christ, our Redeemer, declared that which he offered under the species of bread to be truly his own body, therefore has it ever been a firm belief in the Church of God, and this holy Synod doth now declare it anew, that, by the consecration of the bread and of the wine, a conversion is made of the whole substance of the bread into the substance of the body of Christ our Lord, and of the whole substance of the wine into the substance of his blood; which conversion is, by the holy Catholic Church, suitably and properly called Transubstantiation.

Caput V.

Chapter V.

De cultu et veneratione huic sanctissimo Sacramento exhibenda.

\emph{On the cult and veneration to be shown to this most holy sacrament.}

\emph{Nullus itaque dubitandi locus relinquitur, quin omnes Christi fideles pro more in Catholica Ecclesia semper recepto latriæ cultum, qui vero Deo debetur, huic sanctissimo sacramento in veneratione exhibeant: neque enim ideo minus est adorandum, quod fuerit a Christo Domino, ut sumatur, institutum: nam illum eumdem Deum præsentem in eo adesse credimus, quem Pater æternus introducens in orbem terrarum dicit: Et adorent eum omnes angeli Dei; quem magi procidentes adoraverunt; quem denique in Galilæa ab apostolis adoratum fuisse, scriptura testatur.}

Wherefore, there is no room left for doubt, that all the faithful of Christ may, according to the custom ever received in the Catholic Church, render in veneration the worship of latria, which is due to the true God, to this most holy sacrament. For not therefore is it the less to be adored on this account, that it was instituted by Christ, the Lord, in order to be received; for we believe that same God to be present therein, of whom the eternal Father, when introducing him into the world, says: \emph{And let all the angels of God adore him;}\footnote{\par   Psa. xcvi. 7.} whom the Magi, \emph{falling down, adored;}\footnote{\par   Matt. ii. 11.} who, in fine, as the Scripture testifies, was adored by the apostles in Galilee.

\emph{Declarat præterea sancta synodus, pie et religiose admodum in Dei Ecclesiam inductum fuisse hunc morem, ut singulis annis peculiari quodam et festo die præcelsum hoc et venerabile sacramentum singulari veneratione ac solemniter celebraretur, utque in processionibus reverenter et honorifice illud per vias et loca publica circumferretur. Æquissimum est enim, sacros }

The holy Synod declares, moreover, that very piously and religiously was this custom introduced into the Church, that this sublime and venerable sacrament be, with special veneration and solemnity, celebrated, every year, on a certain day, and that a festival; and that it be borne reverently and with honor in processions through the streets and public places. For it is most just that there be certain appointed

\emph{aliquos statutos esse dies, cum Christiani omnes singulari ac rara quadam significatione gratos et memores testentur animos erga communem Dominum et Redemptorem pro tam ineffabili et plane divino beneficio, quo mortis ejus victoria et triumphus repræsentatur. Ac sic quidem oportuit victricem veritatem de mendacio et hæresi triumphum agere, ut ejus adversarii in conspectu tanti splendoris, et in tanta universæ ecclesiæ lætitia positi, vel debilitati et fracti tabescant, vel pudore affecti et confusi aliquando resipiscant.}

holy days, whereon all Christians may, with a special and unusual demonstration, testify that their minds are grateful and thankful to their common Lord and Redeemer for so ineffable and truly divine a benefit, whereby the victory and triumph of his death are represented. And so indeed did it behoove victorious truth to celebrate a triumph over falsehood and heresy, that thus her adversaries, at the sight of so much splendor, and in the midst of so great joy of the universal Church, may either \emph{pine away}\footnote{\par   Psa. cxi. 10.} weakened and broken; or, touched with shame and confounded, at length repent.

Caput VI.

Chapter VI.

De asservando sacræ Eucharistiæ Sacramento, et ad infirmos deferendo.

On reserving the sacrament of the sacred Eucharist, and bearing it to the sick.

\emph{Consuetudo asservandi in sacrario sanctam Eucharistiam adeo antiqua est, ut eam sæculum etiam Nicæni Concilii agnoverit. Porro deferri ipsam sacram Eucharistiam ad infirmos, et in hunc usum diligenter in ecclesiis conservari, præterquam quod cum summa æquitate et ratione conjunctum est, tum multis in conciliis præceptum invenitur et vetustissimo Catholicæ Ecclesiæ more est observatum.}

The custom of reserving the holy Eucharist in the sacrarium is so ancient, that even the age of the Council of Nicæa recognized that usage. Moreover, as to carrying the sacred Eucharist itself to the sick, and carefully reserving it for this purpose in churches, besides that it is exceedingly conformable to equity and reason, it is also found, enjoined in numerous councils, and is a very ancient observance of the Catholic Church.

\emph{Quare sancta hæc synodus retinendum omnino salutarem hunc et necessarium morem statuit.}

Wherefore, this holy Synod ordains that this salutary and necessary custom is to be by all means retained.

Caput VII.

Chapter VII.

De præparatione, quæ adhibenda est, ut digne quis sacram Eucharistiam percipiat.

On the preparation to be given that one may worthily receive the sacred Eucharist.

\emph{Si non decet ad sacras ullas functiones quempiam accedere nisi sancte, certe, quo magis sanctitas et divinitas cœlestis hujus sacramenti viro Christiano comperta est, eo diligentius cavere ille debet, ne absque magna reverentia et sanctitate ad id percipiendum accedat, præsertim cum illa plena formidinis verba apud apostolum legamus: Qui manducat et bibit indigne, judicium sibi manducat et bibit, non dijudicans corpus Domini. Quare communicare volenti revocandum est in memoriam ejus præceptum: Probet autem seipsum homo. Ecclesiastica autem consuetudo declarat, eam probationem necessariam esse, ut nullus sibi conscius peccati mortalis, quantumvis sibi contritus videatur, absque præmissæ sacramentali confessione ad sacram Eucharistiam accedere debeat. Quod a Christianis omnibus, etiam ab iis sacerdotibus, }

If it is unbeseeming for any one to approach to any of the sacred functions, unless he approach holily; assuredly, the more the holiness and divinity of this heavenly sacrament are understood by a Christian, the more diligently ought he to give heed that he approach not to receive it but with great reverence and holiness, especially as we read in the Apostle those words full of terror: \emph{He that eateth and drinketh unworthily, eateth and drinketh judgment to himself.}\footnote{\par   1 Cor. xi. 29.} Wherefore, he who would communicate, ought to recall to mind the precept of the Apostle: \emph{Let a man prove himself.}\footnote{\par   1 Cor. v. 28.} Now ecclesiastical usage declares that necessary proof to be, that no one, conscious to himself of mortal sin, how contrite soever he may seem to himself, ought to approach to the sacred Eucharist without previous sacramental confession. This the holy Synod hath decreed is to be invariably observed by all Christians,

\emph{quibus ex officio incubuerit celebrare, hæc sancta synodus perpetuo servandum esse decrevit, modo non desit illis copia confessoris. Quod si necessitate urgente sacerdos absque prævia confessione celebraverit, quamprimum confiteatur.}

even by those priests on whom it may be incumbent by their office to celebrate, provided the opportunity of a confessor do not fail them; but if, in an urgent necessity, a priest should celebrate without previous confession, let him confess as soon as possible.

Caput VIII.

Chapter VIII.

De usu admirabilis hujus sacramenti.

On the use of this admirable sacrament.

\emph{Quoad usum autem recte et sapienter Patres nostri tres rationes hoc sanctum sacramentum accipiendi distinxerunt. Quosdam enim docuerunt sacramentaliter dumtaxat id sumere ut peccatores; alios tantum spiritualiter, illos nimirum, qui voto propositum illum cœlestem panem edentes, fide viva, quæ per dilectionem operatur, fructum ejus et utilitatem sentiunt; tertios porro sacramentaliter simul et spiritualiter; hi autem sunt, qui ita se prius probant et instruunt, ut vestem nuptialem induti ad divinam hanc mensam accedant.}

Now as to the use of this holy sacrament, our Fathers have rightly and wisely distinguished three ways of receiving it. For they have taught that some receive it sacramentally only, to wit, sinners: others spiritually only, those to wit who eating in desire that heavenly bread which is set before them, are, by a lively \emph{faith which worketh by charity},\footnote{\par   Gal. v. 6.} made sensible of the fruit and usefulness thereof: whereas the third [class] receive it both sacramentally and spiritually, and these are they who so \emph{ prove} and prepare themselves beforehand, as to approach to this divine table \emph{clothed with the wedding garment.}\footnote{\par   Matt. xxii. 11, 12.} Now as to the reception of the sacrament, it was always the custom in the Church of God that laymen should receive the communion from priests; but that priests when celebrating should communicate themselves;

\emph{In sacramentali autem sumptione semper in Ecclesia Dei mos fuit, ut laici a sacerdotibus communionem acciperent; sacerdotes autem celebrantes seipsos communicarent, qui mos, }

\emph{tamquam ex traditione apostolica descendens, jure ac merito retineri debet.}

which custom, as coming down from an apostolic tradition, ought with justice and reason to be retained.

\emph{Demum autem paterno affectu admonet sancta synodus, hortatur, rogat et obsecrat per viscera misericordiæ Dei nostri, ut omnes et singuli, qui Christiano nomine censentur, in hoc unitatis signo, in hoc vinculo caritatis, in hoc concordiæ symbolo jam tandem aliquando conveniant et concordent, memoresque tantæ majestatis, et tam eximii amoris Iesu Christi, Domini nostri, qui dilectam animam suam in nostræ salutis pretium et carnem suam nobis dedit ad manducandum, hæc sacra mysteria corporis et sanguinis ejus ea fidei constantia et firmitate ea animi devotione, ea pietate et cultu credant et venerentur, ut panem illum supersubstantialem frequenter suscipere possint, et is vere eis sit animæ vita et perpetua sanitas mentis, cujus vigore confortati, ex hujus miseræ peregrinationis itinere ad cœlestem patriam pervenire valeant, eumdem panem angelorum, quem modo sub sacris velaminibus edunt, absque ullo velamine manducaturi.}

And finally this holy Synod, with true fatherly affection, admonishes, exhorts, begs, and beseeches, through the bowels of the mercy of our God, that all and each of those who bear the Christian name would now at length agree and be of one mind in this sign of unity, in this bond of charity, in this symbol of concord; and that, mindful of the so great majesty, and the so exceeding love of our Lord Jesus Christ, who gave his own beloved soul as the price of our salvation, and gave us his own flesh to eat, they would believe and venerate these sacred mysteries of his body and blood, with such constancy and firmness of faith, with such devotion of soul, with such piety and worship, as to be able frequently to receive that supersubstantial bread, and that it may be to them truly the life of the soul and the perpetual health of their mind; that being invigorated by the strength thereof, they may, after the journeying of this miserable pilgrimage, be able to arrive at their heavenly country, there to eat, without any veil, that same bread of angels which they now eat under the sacred veils.

\emph{Quoniam autem non est satis }

But forasmuch as it is not enough

\emph{veritatem dicere, nisi detegantur et refellantur errores: placuit sanctæ synodo hos canones subjungere, ut omnes, jam agnita Catholica doctrina, intelligant quoque, quæ ab illis hæreses caveri, vitarique debeant.}

to declare the truth, if errors be not laid bare and repudiated, it hath seemed good to the holy Synod to subjoin these canons, that all,—the Catholic doctrine being already recognized,—may now also understand what are the heresies which they ought to guard against and avoid.

DE SACROSANCTO EUCHARISTIÆ SACRAMENTO.

ON THE MOST HOLY SACRAMENT OF THE EUCHARIST.

Canon I.—\emph{Si quis negaverit, in sanctissimæ Eucharistiæ sacramento contineri vere, realiter et substantialiter corpus et sanguinem una cum anima et divinitate Domini nostri Jesu Christi, ac proinde totum Christum; sed dixerit, tantummodo esse in eo, ut in signo, vel figura, aut virtute: anathema sit.}

Canon I.—If any one denieth, that, in the sacrament of the most holy Eucharist, are contained truly, really, and substantially, the body and blood together with the soul and divinity of our Lord Jesus Christ, and consequently the whole Christ; but saith that he is only therein as in a sign, or in figure, or virtue: let him be anathema.

Canon II.—\emph{Si quis dixerit, in sacrosancto Eucharistiæ sacramento remanere substantiam panis et vini una cum corpore et sanguine Domini nostri Iesu Christi, negaveritque mirabilem illam et singularem conversionem totius substantiæ panis in corpus, et totius substantiæ vini in sanguinem, manentibus dumtaxat speciebus panis et vini; quam quidem conversionem Catholica Ecclesia aptissime }

Canon II.—If any one saith, that, in the sacred and holy sacrament of the Eucharist, the substance of the bread and wine remains conjointly with the body and blood of our Lord Jesus Christ, and denieth that wonderful and singular conversion of the whole substance of the bread into the body, and of the whole substance of the wine into the blood—the species only of the bread and wine remaining—which conversion indeed the Catholic

\emph{Transsubstantiationem appellat: anathema, sit.}

Church most aptly calls Transubstantiation: let him be anathema.

Canon III.—\emph{Si quis negaverit, in venerabili sacramento Eucharistiæ sub unaquaque specie, et sub singulis cujusque speciei partibus, separatione facta, totum Christum contineri: anathema sit.}

Canon III.—If any one denieth, that, in the venerable sacrament of the Eucharist, the whole Christ is contained under each species, and under every part of each species, when separated: let him be anathema.

Canon IV.—\emph{Si quis dixerit, peracta consecratione, in admirabili Eucharistiæ sacramento non esse corpus et sanguinem Domini nostri Iesu Christi, sed tantum in usu, dum sumitur, non autem ante vel post, et in hostiis seu particulis consecratis, quæ post communionem reservantur vel supersunt, non remanere verum corpus Domini: anathema sit.}

Canon IV.—If any one saith, that, after the consecration is completed, the body and blood of our Lord Jesus Christ are not in the admirable sacrament of the Eucharist, but [are there] only during the use, whilst it is being taken, and not either before or after; and that, in the hosts, or consecrated particles, which are reserved or which remain after communion, the true body of the Lord remaineth not: let him be anathema.

Canon V.—\emph{Si quis dixerit, vel præcipuum fructum sanctissimæ Eucharistiæ esse remissionem peccatorum, vel ex ea non alios effectus provenire: anathema sit.}

Canon V.—If any one saith, either that the principal fruit of the most holy Eucharist is the remission of sins, or that other effects do not result therefrom: let him be anathema.

Canon VI.—\emph{Si quis dixerit, in sancto Eucharistiæ sacramento Christum, unigenitum Dei Filium, non esse cultu latriæ etiam externo adorandum, atque ideo non festiva peculiari celebritate venerandum, neque in processionibus secundum laudabilem et universalem Ecclesiæ }

Canon VI.—If any one saith, that, in the holy sacrament of the Eucharist, Christ, the only-begotten Son of God, is not to be adored with the worship, even external of latria; and is, consequently, neither to be venerated with a special festive solemnity, nor to be solemnly borne about in procession, according

\emph{sanctæ ritum et consuetudinem solemniter circumgestandum, vel non publice, ut adoretur, populo proponendum, et ejus adoratores esse idololatras: anathema sit.}

to the laudable and universal rite and custom of holy Church; or, is not to be proposed publicly to the people to be adored, and that the adorers thereof are idolaters: let him be anathema.

Canon VII.—\emph{Si quis dixcerit, non licere sacram Eucharistiam in sacrario reservari, sed statim post consecrationem adstantibus necessario distribuendam; aut non licere, ut illa ad infirmos honorifice deferatur: anathema sit.}

Canon VII.—If any one saith, that it is not lawful for the sacred Eucharist to be reserved in the \emph{sacrarium}, but that, immediately after consecration, it must necessarily be distributed amongst those present; or, that it is not lawful that it be carried with honor to the sick: let him be anathema.

Canon VIII.—\emph{Si quis dixerit, Christum in Eucharistia exhibitum spiritualiter tantum manducari, et non etiam sacramentaliter ac realiter: anathema sit.}

Canon VIII.—If any one saith, that Christ, given in the Eucharist, is eaten spiritually only, and not also sacramentally and really: let him be anathema.

Canon IX.—\emph{Si quis negaverit, omnes et singulos Christi fideles utriusque sexus, cum ad annos discretionis pervenerint, teneri singulis annis, saltem in paschate, ad communicandum, juxta præceptum sanctæ matris Ecclesiæ: anathema sit.}

Canon IX.—If any one denieth, that all and each of Christ's faithful of both sexes are bound, when they have attained to years of discretion, to communicate every year, at least at Easter, in accordance with the precept of holy Mother Church: let him be anathema.

Canon X.—\emph{Si quis dixerit, non licere sacerdoti celebranti seipsum communicare: anathema sit.}

Canon X.—If any one saith, that it is not lawful for the celebrating priest to communicate himself: let him be anathema.

Canon XI.—\emph{Si quis dixerit, solam fidem esse sufficientem præparationem ad sumendum sanctissimæ Eucharistiæ sacramentum: }

Canon XI.—If any one saith, that faith alone is a sufficient preparation for receiving the sacrament of the most holy Eucharist: let him

\emph{anathema sit. Et, ne tantum sacramentum indigne, atque ideo in mortem et condemnationem sumatur, statuit atque declarat ipsa sancta synodus illis, quos conscientia peccati mortalis gravat, quantumcumque etiam se contritos existiment, habita copia confessoris, necessario præmittendam esse confessionem sacramentalem. Si quis autem contrarium docere, prædicare, vel pertinaciter asserere, seu etiam publice disputando defendere præsumpserit, eo ipso excommunicatus existat.}

be anathema. And for fear lest so great a sacrament may be received unworthily, and so unto death and condemnation, this holy Synod ordains and declares, that sacramental confession, when a confessor may be had, is of necessity to be made beforehand, by those whose conscience is burthened with mortal sin, how contrite even soever they may think themselves. But if any one shall presume to teach, preach, or obstinately to assert, or even in public disputation to defend the contrary, he shall be thereupon excommunicated.

\xchapter{Fourteenth Session, held November 25, 1551.}

Sessio Decimaquarta,

Fourteenth Session,

\emph{celebrata die XXV. Nov.} 1551.

\emph{held November} 25, 1551.

DE SANCTISSIMIS PŒNITENTIÆ ET EXTREMÆ UNCTIONIS SACRAMENTIS.

ON THE MOST HOLY SACRAMENTS OF PENANCE AND EXTREME UNCTION.

Caput I.

Chapter I.

De necessitate et institutione Sacramenti pœnitentiæ

\emph{On the necessity, and on the institution of the Sacrament of Penance.}

\emph{Si ea in regeneratis omnibus gratitudo erga Deum esset, ut justitiam in baptismo, ipsius beneficio et gratia susceptam constanter tuerentur, non fuisset opus, aliud ab ipso baptismo sacramentum ad peccatorum remissionem esse institutum. Quoniam autem Deus, dives in misericordia, cognovit figmentum nostrum, }

If such, in all the regenerate, were their gratitude towards God, as that they constantly preserved the justice received in baptism by his bounty and grace, there would not have been need for another sacrament, besides that of baptism itself, to be instituted for the remission of sins. But because God, \emph{rich in mercy, knows our frame},\footnote{\par   Psa. cii. 14.} he hath

\emph{illis etiam vitæ remedium contulit, qui se postea in peccati servitutem et dæmonis potestatem tradidissent, sacramentum videlicet pœnitentiæ, quo lapsis post baptismum beneficium mortis Christi applicatur. Fuit quidem pœnitentia universis hominibus, qui se mortali aliquo peccato inquinassent, quovis tempore ad gratiam et justitiam assequendam necessaria, illis etiam, qui baptismi sacramento ablui petivissent, ut, perversitate abjecta et emendata, tantam Dei offensionem cum peccati odio et pio animi dolore detestarentur; unde propheta ait; Convertimini, et agite pœnitentiam ab omnibus iniquitatibus vestris; et non erit vobis in ruinam iniquitas. Dominus etiam dixit: Nisi pœnitentiam egeritis, omnes similiter peribitis. Et princeps apostolorum Petrus peccatoribus baptismo initiandis pœnitentiam commendans dicebat: Pœnitentiam agite, et baptizetur unusquisque vestrum. Porro nec ante adventum Christi pœnitentia erat sacramentum, nec est post adventum illius cuiquam ante baptismum. Dominus autem sacramentum pœnitentiæ tunc præcipue }

bestowed a remedy of life even on those who may, after baptism, have delivered themselves up to the servitude of sin and the power of the devil,—the sacrament to wit of Penance, by which the benefit of the death of Christ is applied to those who have fallen after baptism. Penitence was indeed at all times necessary, in order to attain to grace and justice, for all men who had defiled themselves by any mortal sin, even for those who begged to be washed by the sacrament of Baptism; that so, their perverseness renounced and amended, they might, with a hatred of sin and a godly sorrow of mind, detest so great an offense of God. Wherefore the prophet says: \emph{Be converted and do penance for all your iniquities, and iniquity shall not be your ruin.}\footnote{\par   Ezek. xviii. 30.} The Lord also said: \emph{Except you do penance, you shall also likewise perish;}\footnote{\par   Luke xiii. 5.} and Peter, the prince of the apostles, recommending penitence to sinners who were about to be initiated by baptism, said: \emph{Do penance, and be baptized every one of you.}\footnote{\par   Acts ii. 38.} Nevertheless, neither before the coming of Christ was penitence a sacrament, nor is it such, since his coming, to any previously to baptism. But the Lord then principally instituted the sacrament

\emph{instituit, cum a mortuis excitatus insufflavit in discipulos suos, dicens: Accipite Spiritum Sanctum; quorum remiseritis peccata, remittuntur eis, et quorum retinueritis, retenta sunt. Quo tam insigni facto et verbis tam perspicuis potestatem remittendi et retinendi peccata, ad reconciliandos fideles post baptismum lapsos, apostolis et eorum legitimis successoribus fuisse communicatam, universorum patrum consensus semper intellexit, et Novatianos, remittendi potestatem olim pertinaciter negantes, magna ratione Ecclesia Catholica, tamquam hæreticos, explosit atque condemnavit. Quare verissimum hunc illorum verborum Domini sensum sancta hæc synodus probans et recipiens, damnat eorum commentitias interpretationes, qui verba illa ad potestatem prædicandi verbum Dei et Christi evangelium annuntiandi, contra hujusmodi sacramenti institutionem, falso detorquent.}

of penance, when, being raised from the dead, he breathed upon his disciples, saying: \emph{Receive ye the Holy Ghost: whose sins you shall forgive, they are forgiven them, and whose sins you shall retain, they are retained.}\footnote{\par  John xx. 23.} By which action so signal, and words so clear, the consent of all the Fathers has ever understood that the power of \emph{forgiving and retaining sins} was communicated to the apostles and their lawful successors, for the reconciling of the faithful who have fallen after baptism. And the Catholic Church with great reason repudiated and condemned as heretics the Novatians, who of old obstinately denied that power of forgiving. Wherefore, this holy Synod, approving of and receiving as most true this meaning of those words of our Lord, condemns the fanciful interpretations of those who, in opposition to the institution of this sacrament, falsely wrest those words to the power of preaching the Word of God, and of announcing the Gospel of Christ.

Caput II.

Chapter II.

De differentia Sacramenti pœnitentiæ et Baptismi.

On the difference between the Sacrament of Penance and that of Baptism.

\emph{Ceterum hoc sacramentum multis rationibus a baptismo differre }

For the rest, this sacrament is clearly seen to be different from

\emph{dignoscitur. Nam præterquam quod materia et forma, quibus sacramenti essentia perficitur, longissime dissidet: constat certe, baptismi ministrum judicem esse non oportere, cum Ecclesia in neminem judicium exerceat, qui non prius in ipsam per baptismi januam fuerit ingressus. Quid enim mihi, inquit apostolus, de iis, qui foris sunt, judicare? Secus est de domesticis fidei, quos Christus dominus lavacro baptismi sui corporis membra semel effecit; nam hos, si se postea crimine aliquo contaminaverint, non jam repetito baptismo ablui, cum id in Ecclesia Catholica nulla ratione liceat, sed ante hoc tribunal tamquam reos sisti voluit, ut per sacerdotum sententiam non semel, sed quoties ab admissis peccatis ad ipsum pœnitentes confugerint, possent liberari. Alius præterea est baptismi, et alius pœnitentiæ fructus; per baptismum enim Christum induentes, nova prorsus in illo efficimur creatura, plenam et integram peccatorum omnium remissionem consequentes: ad quam tamen novitatem, et integritatem per sacramentum pœnitentiæ, sine magnis nostris fletibus}

baptism in many respects: for besides that it is very widely different indeed in matter and form, which constitute the essence of a sacrament, it is beyond doubt certain that the minister of baptism need not be a judge, seeing that the Church exercises judgment on no one who has not entered therein through the gate of baptism. For, \emph{what have I}, saith the apostle, \emph{to do to judge them that are without?}\footnote{\par   1 Cor. v. 12.} It is otherwise with those who are of \emph{the household of the faith}, whom Christ our Lord has once, by the laver of baptism, made the members of his own body; for such, if they should afterwards have defiled themselves, by any crime, he would no longer have them cleansed by a repetition of baptism—that being nowise lawful in the Catholic Church—but be placed as criminals before this tribunal; that, by the sentence of the priests, they might be freed, not once, but as often as, being penitent, they should, from their sins committed, flee thereunto. Furthermore, one is the fruit of baptism, and another that of penance. For, by baptism \emph{putting on Christ},\footnote{\par   Gal. iii. 23.} we are made therein entirely a new creature, obtaining a full and entire remission of all sins; unto which newness and

\emph{et laboribus, divina id exigente justitia, pervenire nequaquam possumus, ut merito pœnitentia laboriosus quidam baptismus a sanctis patribus dictus fuerit. Est autem hoc sacramentum pœnitentiæ lapsis post baptismum ad salutem necessarium, ut nondum regeneratis ipse baptismus.}

entireness, however, we are no ways able to arrive by the sacrament of Penance, without many tears and great labors on our parts, the divine justice demanding this; so that penance has justly been called by holy Fathers a laborious kind of baptism. And this sacrament of Penance is, for those who have fallen after baptism, necessary unto salvation; as baptism itself is for those who have not as yet been regenerated.

Caput III.

Chapter III.

De partibus et fructibus hujus sacramenti.

On the parts and on the fruit of this sacrament.

\emph{Docet præterea sancta synodus, sacramenti pœnitentiæ formam, in qua præcipue ipsius vis sita est, in illis ministri verbis positam esse: Ego te absolvo, etc. Quibus quidem de Ecclesiæ sanctæ more preces quædam laudabiliter adjunguntur; ad ipsius tamen formæ essentiam nequaquam spectant, neque ad ipsius sacramenti administrationem sunt necessariæ. Sunt autem quasi materia hujus sacramenti ipsius pœnitentis actus, nempe contritio, confessio, et satisfactio. Qui quatenus in pœnitente ad integritatem sacramenti, ad plenamque et perfectam peccatorum remissionem ex Dei institutione requiruntur, }

The holy Synod doth furthermore teach, that the form of the sacrament of Penance, wherein its force principally consists, is placed in those words of the minister: \emph{I absolve thee}, etc.; to which words indeed certain prayers are, according to the custom of holy Church, laudably joined, which nevertheless by no means regard the essence of that form, neither are they necessary for the administration of the sacrament itself. But the acts of the penitent himself, to wit, contrition, confession, and satisfaction, are as it were the matter of this sacrament. Which acts, inasmuch as they are, by God's institution, required in the penitent for the integrity of the sacrament, and for the full and perfect

\emph{hac ratione pœnitentiæ partes dicuntur. Sane vero res et effectus hujus sacramenti, quantum ad ejus vim et efficaciam pertinet, reconciliatio est cum Deo, quam interdum in viris piis, et cum devotione hoc sacramentum percipientibus, conscientiæ pax ac serenitas cum vehementi spiritus consolatione consequi solet. Hæc de partibus et effectu hujus sacramenti sancta synodus tradens, simul eorum sententias damnat, qui pœnitentiæ partes incussos conscientiæ terrores et fidem esse contendunt.}

remission of sins, are for this reason called the parts of penance. But the thing signified indeed, and the effect of this sacrament, as far as regards its force and efficacy, is reconciliation with God, which sometimes, in persons who are pious and who receive this sacrament with devotion, is wont to be followed by peace and serenity of conscience, with exceeding consolation of spirit. The holy Synod, whilst delivering these things touching the parts and the effect of this sacrament, condemns at the same time the opinions of those who contend that the terrors which agitate the conscience, and faith, are the parts of penance.

Caput IV.

Chapter IV.

De Contritione.

On Contrition.

\emph{Contritio, quæ primum locum inter dictos pœnitentis actus habet, animi dolor ac detestatio est de peccato commisso, cum proposito non peccandi de cetero. Fuit autem quovis tempore ad impetrandam veniam peccatorum hic contritionis motus necessarius, et in homine post baptismum lapso ita demum præparat ad remissionem peccatorum, si cum fiducia divinæ misericordiæ et voto præstandi reliqua conjunctus sit, quæ ad rite suscipiendum }

Contrition, which holds the first place amongst the aforesaid acts of the penitent, is a sorrow of mind, and a detestation for sin committed, with the purpose of not sinning for the future. This movement of contrition was at all times necessary for obtaining the pardon of sins; and, in one who has fallen after baptism, it then at length prepares for the remission of sins, when it is united with confidence in the divine mercy, and with the desire of performing the other things which are required for rightly receiving this sacrament.

\emph{hoc sacramentum requiruntur. Declarat igitur sancta synodus, hanc contritionem non solum cessationem a peccato et vitæ novas propositum et inchoationem, sed veteris etiam odium continere, juxta illud: Projicite a vobis omnes iniquitates vestras, in quibus prævaricati estis, et facite vobis cor novum et spiritum novum. Et certe, qui illos sanctorum clamores consideraverit: Tibi soli peccavi, et malum coram te feci; Laboravi in gemitu meo, lavabo per singulas noctes lectum meum. Recogitabo tibi omnes annos meos in amaritudine animæ meæ; et alios hujus generis, facile intelliget, eos ex vehementi quodam anteactæ vitæ odio et ingenti peccatorum detestatione manasse. Docet præterea, etsi contritionem hanc aliquando caritate perfectam esse contingat, hominemque Deo reconciliare, priusquam hoc sacramentum actu suscipiatur, ipsam nihilominus reconciliationem ipsi contritioni sine sacramenti voto, quod in illa includitur, non esse adscribendam. Illam vero contritionem imperfectam, quæ attritio dicitur, }

Wherefore the holy Synod declares, that this contrition contains not only a cessation from sin, and the purpose and the beginning of a new life, but also a hatred of the old, agreeably to that saying: \emph{Cast away from you all your iniquities, wherein you have transgressed, and make to yourselves a new heart and a new spirit.}\footnote{\par   Ezek. xviii. 31.} And assuredly he who has considered those cries of the saints: \emph{To thee only have I sinned, and have done evil before thee;}\footnote{\par   Psa. l. 6.} \emph{I have labored in my groaning, every night I will wash my bed;}\footnote{\par  Psa. vi. 7.} \emph{I will recount to thee all my years, in the bitterness of my soul;}\footnote{\par   Isa. xxxviii. 15.} and others of this kind, will easily understand that they flowed from a certain vehement hatred of their past life, and from an exceeding detestation of sins. The Synod teaches moreover, that, although it sometimes happens that this contrition is perfect through charity, and reconciles man with God before this sacrament be actually received, the said reconciliation, nevertheless, is not to be ascribed to that contrition, independently of the desire of the sacrament which is included therein. And as to that imperfect contrition, which is called attrition, because

\emph{quoniam vel ex turpitudinis peccati consideratione vel ex gehennæ et pœnarum metu communiter concipitur, si voluntatem peccandi excludat cum spe veniæ, declarat non solum non facere hominem hypocritam et magis peccatorem, verum etiam donum Dei esse et Spiritus Sancti impulsum, non adhuc quidem inhabitantis, sed tantum moventis, quo pœnitens adjutus viam sibi ad justitiam parat. Et quamvis sine sacramento pœnitentiæ per se ad justificationem perducere peccatorem nequeat, tamen eum ad Dei gratiam in sacramento pœnitentiæ impetrandam disponit: hoc enim timore utiliter concussi Ninivitæ, ad Ionæ prædicationem, plenam terroribus pœnitentiam egerunt et misericordiam a Domino impetrarunt. Quamobrem falso quidam calumniantur Catholicos scriptores, quasi tradiderint, sacramentum pœnitentiæ absque bono motu suscipientium gratiam conferre, quod numquam Ecclesia Dei docuit, neque sensit; sed et falso docent, contritionem esse extortam et coactam, non liberam et voluntariam.}

that it is commonly conceived either from the consideration of the turpitude of sin, or from the fear of hell and of punishment, it declares that if, with the hope of pardon, it exclude the wish to sin, it not only does not make a man a hypocrite, and a greater sinner, but that it is even a gift of God, and an impulse of the Holy Ghost,—who does not indeed as yet dwell in the penitent, but only moves him,—whereby the penitent being assisted prepares a way for himself unto justice. And although this [attrition] can not of itself, without the sacrament of Penance, conduct the sinner to justification, yet does it dispose him to obtain the grace of God in the sacrament of Penance. For, smitten profitably with this fear, the Ninivites, at the preaching of Jonas, did fearful penance, and obtained mercy from the Lord. Wherefore falsely do some calumniate Catholic writers, as if they had maintained that the sacrament of Penance confers grace without any good motion on the part of those who receive it: a thing which the Church of God never taught, or thought; and falsely also do they assert that contrition is extorted and forced, not free and voluntary.

Caput V.

Chapter V.

De Confessione.

On Confession.

\emph{Ex institutione sacramenti pœnitentiæ jam explicata universa Ecclesia semper intellexit, institutam etiam esse a Domino integram peccatorum confessionem, et omnibus post baptismum lapsis jure divino necessariam existere, quia Dominus noster Iesus Christus, e terris ascensurus ad cœlos, sacerdotes sui ipsius vicarios reliquit, tamquam præsides et judices, ad quos omnia mortalia crimina deferantur, in quæ Christi fideles ceciderint, quo, pro potestate clavium, remissionis aut retentionis peccatorum sententiam pronuncient. Constat enim, sacerdotes judicium hoc incognita causa exercere non potuisse, nec æquitatem quidem illos in pænis injungendis servare potuisse, si in genere tumtaxat, et non potius in specie, ac sigillatim sua ipsi peccata declarassent. Ex his colligitur, oportere a pœnitentibus omnia peccata mortalia, quorum post diligentem sui discussionem conscientiam habent, in confessione recenseri, etiam si occultissima illa sint et tantum adversus duo ultima decalogi præcepta commissa, quæ nonnunquam animum }

From the institution of the sacrament of Penance, as already explained, the universal Church has always understood that the entire confession of sins was also instituted by the Lord, and is of divine right necessary for all who have fallen after baptism; because that our Lord Jesus Christ, when about to ascend from earth to heaven, left priests his own vicars, as presidents and judges, unto whom all the mortal crimes, into which the faithful of Christ may have fallen, should be carried, in order that, in accordance with the power of the keys, they may pronounce the sentence of forgiveness or retention of sins. For it is manifest that priests could not have exercised this judgment without knowledge of the cause; neither indeed could they have observed equity in enjoining punishments, if the said faithful should have declared their sins in general only, and not rather specifically, and one by one. Whence it is gathered that all the mortal sins, of which, after a diligent examination of themselves, they are conscious, must needs be by penitents enumerated in confession, even though those sins be most hidden, and committed only against the two last precepts of the

\emph{gravius sauciant, et periculosiora sunt iis, quæ in manifesto admittuntur. Nam venialia, quibus a gratia Dei non excludimur et in quæ frequentius labimur, quamquam recte et utiliter citraque omnem præsumptionem in confessione dicantur, quod piorum hominum usus demonstrat, taceri tamen citra culpam multisque aliis remediis expiari possunt. Verum, cum universa mortalia peccata, etiam cogitationis, homines iræ filios et Dei inimicos reddant, necessum est, omnium etiam veniam cum aperta et verecunda confessione, a Deo quærere. Itaque dum omnia, quæ memoriæ occurrunt, peccata Christi fideles confiteri student, procul dubio omnia divinæ misericordiæ ignoscenda exponunt. Qui vero secus faciunt et scienter aliqua retinent, nihil divinæ bonitati per sacerdotem remittendum proponunt. Si enim erubescat ægrotus vulnus medico detegere, quod ignorat, medicina non curat. Colligitur præterea, etiam eas circumstantias in confessione explicandas esse, quæ speciem peccati mutant, quod sine illis peccata ipsa neque a pœnitentibus }

decalogue,—sins which sometimes wound the soul more grievously, and are more dangerous, than those which are committed outwardly. For venial sins, whereby we are not excluded from the grace of God, and into which we fall more frequently, although they be rightly and profitably, and without any presumption, declared in confession, as the custom of pious persons demonstrates, yet may they be omitted without guilt, and be expiated by many other remedies. But, whereas all mortal sins, even those of thought, render men \emph{children of wrath},\footnote{\par  Ephes. ii. 3.} and enemies of God, it is necessary to seek also for the pardon of them all from God, with an open and modest confession. Wherefore, while the faithful of Christ are careful to confess all the sins which occur to their memory, they without doubt lay them all bare before the mercy of God to be pardoned: whereas they who act otherwise, and knowingly keep back certain sins, such set nothing before the divine bounty to be forgiven through the priest; for if the sick be ashamed to show his wound to the physician, his medical art cures not that which it knows not of. We gather, furthermore, that those circumstances which change the species

\emph{integre exponantur, nec judicibus innotescant; et fieri nequeat, ut de gravitate criminum recte censere possint et pœnam, quam oportet, pro illis pœnitentibus imponere. Unde alienum a ratione est docere, circumstantias has ab hominibus otiosis excogitatas fuisse, aut unam tantum circumstantiam confitendam esse, nempe peccasse in fratrem. Sed et impium est, confessionem, quæ hac ratione fieri præcipitur, impossibilem dicere, aut carnificinam, illam conscientiarum appellare; constat enim, nihil aliud in Ecclesia a pœnitentibus exigi, quam ut, postquam quisque diligentius se excusserit et conscientiæ suæ sinus omnes et latebras exploraverit, ea peccata confiteatur, quibus se Dominum et Deum suum mortaliter offendissi meminerit; reliqua autem peccata, quæ diligenter cogitanti non occurrunt, in universum eadem confessione inclusa esse intelliguntur; pro quibus fideliter cum propheta dicimus: Ab occultis meis munda me, Domine. Ipsa vero hujusmodi confessionis difficultas ac peccata detegendi verecundia, gravis quidem videri }

of the sin are also to be explained in confession, because that, without them, the sins themselves are neither entirely set forth by the penitents, nor are they known clearly to the judges; and it can not be that they can estimate rightly the grievousness of the crimes, and impose on the penitents the punishment which ought to be inflicted on account of them. Whence it is unreasonable to teach that these circumstances have been invented by idle men; or that one circumstance only is to be confessed, to wit, that one has sinned against a brother. But it is also impious to assert, that confession, enjoined to be made in this manner, is impossible, or to call it a slaughter-house of consciences; for it is certain, that in the Church nothing else is required of penitents, but that, after each has examined himself diligently, and searched all the folds and recesses of his conscience, he confess those sins by which he shall remember that he has mortally offended his Lord and God: whilst the other sins, which do not occur to him after diligent thought, are understood to be included as a whole in that same confession; for which sins we confidently say with the prophet: \emph{From my secret sins cleanse me, O Lord.}\footnote{\par   Psa. xviii. 13.} Now, the

\emph{posset, nisi tot tantisque commodis et consolationibus levaretur, quæ omnibus digne ad hoc sacramentum accedentibus per absolutionem certissime conferuntur. Ceterum, quoad modum confitendi secreto apud solum sacerdotem, etsi Christus non vetuerit, quin aliquis in vindictam suorum scelerum et sui humiliationem, cum ob aliorum exemplum, tum ob Ecclesiæ offensæ ædificationem delicta sua publice confiteri possit: non est tamen hoc divino præcepto mandatum, nec satis consulte humana aliqua lege præciperetur, ut delicta, præsertim secreta, publica essent confessione aperienda; unde cum a sanctissimis et antiquissimis patribus magno unanimique consensu secreta confessio sacramentalis, qua ab initio Ecclesia sancta usa est et modo etiam utitur, fuerit semper commendata, manifeste refellitur inanis eorum calumnia, qui eam a divino mandato alienam et inventum humanum esse, atque a patribus in concilio lateranensi congregatis initium habuisse, docere non verentur; neque enim per lateranense concilium Ecclesia statuit, ut Christi fideles confiterentur, }

very difficulty of a confession like this, and the shame of making known one's sins, might indeed seem a grievous thing, were it not alleviated by the so many and so great advantages and consolations, which are most assuredly bestowed by absolution upon all who worthily approach to this sacrament. For the rest, as to the manner of confessing secretly to a priest alone, although Christ has not forbidden that a person may,—in punishment of his sins, and for his own humiliation, as well for an example to others as for the edification of the Church that has been scandalized,—confess his sins publicly, nevertheless this is not commanded by a divine precept; neither would it be very prudent to enjoin by any human law, that sins, especially such as are secret, should be made known by a public confession. Wherefore, whereas the secret sacramental confession, which was in use from the beginning in holy Church, and is still also in use, has always been commended by the most holy and the most ancient Fathers with a great and unanimous consent, the vain calumny of those is manifestly refuted, who are not ashamed to teach that confession is alien from the divine command, and is a human invention, and that it took its

\emph{quod jure divino necessarium et institutum esse intellexerat, sed ut præceptum confessionis, saltem semel in anno, ab omnibus et singulis, cum ad annos discretionis pervenissent, impleretur; unde jam in universa Ecclesia cum ingenti animarum fidelium fructu observatur mos ille salutaris confitendi sacro illo et maxime acceptabili tempore quadragesimæ: quem morem hæc sancta synodus maxime probat et amplectitur, tamquam pium et merito retinendum.}

rise from the Fathers assembled in the Council of Lateran: for the Church did not, through the Council of Lateran, ordain that the faithful of Christ should confess,—a thing which it knew to be necessary, and to be instituted of divine right,—but that the precept of confession should be complied with, at least once a year, by all and each, when they have attained to years of discretion. Whence, throughout the whole Church, the salutary custom is, to the great benefit of the souls of the faithful, now observed, of confessing at that most sacred and most acceptable time of Lent,—a custom which this holy Synod most highly approves of and embraces, as pious and worthy of being retained.

Caput VI.

Chapter VI.

De ministro hujus sacramenti et Absolutione.

On the ministry of this sacrament, and on Absolution.

\emph{Circa ministrum autem hujus sacramenti declarat sancta synodus, falsas esse et a veritate evangelii penitus alienas doctrinas omnes, quæ ad alios quosvis homines, præter episcopos et sacerdotes clavium ministerium perniciose extendunt, putantes verba illa Domini: Quæcumque alligaveritis super terram, erunt alligata et in cœlo, et quæcumque solveritis }

But, as regards the minister of this sacrament, the holy Synod declares all those doctrines to be false, and utterly alien from the truth of the Gospel, which perniciously extend the ministry of the keys to any others soever besides bishops and priests; imagining, contrary to the institution of this sacrament, that those words of our Lord, \emph{Whatsoever you shall bind, upon earth, shall be bound also in heaven, and}

\emph{super terram, erunt soluta et in cœlo; et: Quorum remiseritis peccata, remittuntur eis, et quorum retinueritis, retenta sunt: ad omnes Christi fideles, indifferenter et promiscue, contra institutionem hujus sacramenti ita fuisse dicta, ut quivis potestatem habeat remittendi peccata, publica quidem per correptionem, si correptus acquieverit, secreta vero per spontaneam confessionem cuicumque factam. Docet quoque, etiam sacerdotes, qui peccato mortali tenentur, per virtutem Spiritus Sancti in ordinatione collatam, tamquam Christi ministros, functionem remittendi peccata exercere, eosque prave sentire, qui in malis sacerdotibus hanc potestatem non esse contendunt. Quamvis autem absolutio sacerdotis alieni beneficii sit dispensatio, tamen non est solum nudum ministerium vel annuntiandi evangelium, vel declarandi remissa esse peccata; sed ad instar actus judicialis, quo ab ipso, velut a judice, sententia pronuntiatur. Atque ideo non debet pænitens adeo sibi de sua ipsius fide blandiri, ut, etiam si nulla illi adsit contritio, aut sacerdote animus serio }

\emph{whatsoever you shall loose upon earth shall be loosed also in heaven},\footnote{\par   Matt. xviii. 18.} and, \emph{Whose sins you shall forgive, they are forgiven them, and whose sins you shall retain, they are retained},\footnote{\par   John xx. 23.} were in such wise addressed to all the faithful of Christ indifferently and indiscriminately, as that every one has the power of forgiving sins,—public sins to wit by rebuke, provided he that is rebuked shall acquiesce, and secret sins by a voluntary confession made to any individual whatsoever. It also teaches, that even priests, who are in mortal sin, exercise, through the virtue of the Holy Ghost which was bestowed in ordination, the office of forgiving sins, as the ministers of Christ; and that their sentiment is erroneous who contend that this power exists not in bad priests. But although the absolution of the priest is the dispensation of another's bounty, yet is it not a bare ministry only, whether of announcing the Gospel, or of declaring that sins are forgiven, but is after the manner of a judicial act, whereby sentence is pronounced by the priest as by a judge; and therefore the penitent ought not, so to confide in his own personal faith as to think that,—even though there be no contrition on his part, or no intention on the

\emph{agendi et vere absolvendi desit, putet tamen se propter suam solam fidem vere et coram Deo esse absolutum. Nec enim fides sine pœnitentia remissionem ullam peccatorum præstaret; nec is esset nisi salutis suæ negligentissimus, qui sacerdotem joco se absolventem cognosceret, et non alium serio agentem sedulo requireret.}

part of the priest of acting seriously and absolving truly,—he is nevertheless truly and in God's sight absolved, on account of his faith alone. For neither would faith without penance bestow any remission of sins, nor would he be otherwise than most careless of his own salvation, who, knowing that a priest but absolved him in jest, should not carefully seek for another who would act in earnest.

Caput VII.

Chapter VII.

De casuum reservatione.

On the reservation of cases.

\emph{Quoniam igitur natura et ratio judicii illud exposcit, ut sententia in subditos dumtaxat feratur, persuasum semper in Ecclesia Dei fuit, et verissimum esse synodus hæc confirmat, nullius momenti absolutionem eam esse debere, quam sacerdos in eum profert, in quem ordinariam aut subdelegatam non habet jurisdictionem. Magnopere vero ad Christiani populi disciplinam pertinere sanctissimis patribus nostris visum est, ut atrociora quædam et graviora crimina non a quibusvis, sed a summis dumtaxat sacerdotibus absolverentur; unde merito Pontifices maximi pro suprema potestate sibi in Ecclesia universa tradita causas aliquas criminum }

Wherefore, since the nature and order of a judgment require this, that sentence be passed only on those subject [to that judicature], it has ever been firmly held in the Church of God, and this Synod ratifies it as a thing most true, that the absolution, which a priest pronounces upon one over whom he has not either an ordinary or a delegated jurisdiction, ought to be of no weight whatever. And it hath seemed to our most holy Fathers to be of great importance to the discipline of the Christian people, that certain more atrocious and more heinous crimes should be absolved, not by all priests, but only by the highest priests; whence the Sovereign Pontiffs, in virtue of the supreme power delivered to them in

\emph{graviores suo potuerunt peculiari judicio reservare. Neque dubitandum esset, quando omnia, quæ a Deo sunt, ordinata sunt, quin hoc idem episcopis omnibus in sua cuique diæcesi, in ædificationem tamen, non in destructionem liceat, pro illis in subditos tradita supra reliquos inferiores sacerdotes auctoritate, præsertim quoad illa, quibus excommunicationis censura annexa est. Hanc autem delictorum reservationem consonum est divinæ auctoritati non tantum in externa politia, sed etiam coram Deo vim habere. Verumtamen pie admodum, ne hac ipsa occasione aliquis pereat, in eadem Ecclesia Dei custoditum semper fuit, ut nulla sit reservatio in articulo mortis; atque ideo omnes sacerdotes quoslibet pœnitentes a quibusvis peccatis et censuris absolvere possunt; extra quem articulum sacerdotes cum nihil possint in casibus reservatis, id unum pœnitentibus persuadere nitantur, ut ad superiores et legitimos judices pro beneficio absolutionis accedant.}

the universal Church, were deservedly able to reserve, for their special judgment, certain more grievous cases of crimes. Neither is it to be doubted,—seeing that all things, that are from God, are well ordered,—but that this same may be lawfully done by all bishops, each in his own diocese, unto edification, however, not unto destruction, in virtue of the authority, above [that of] other inferior priests, delivered to them over their subjects, especially as regards those crimes to which the censure of excommunication is annexed. But it is consonant to the divine authority, that this reservation of cases have effect, not merely in external polity, but also in God's sight. Nevertheless, for fear lest any may perish on this account, it has always been very piously observed in the said Church of God, that there be no reservation at the point of death, and that therefore all priests may absolve all penitents whatsoever from every kind of sins and censures whatever: and as, save at that point of death, priests have no power in reserved cases, let this alone be their endeavor, to persuade penitents to repair to superior and lawful judges for the benefit of absolution.

Caput VIII.

Chapter VIII.

De Satisfactionis necessitate et fructu.

On the necessity and on the fruit of Satisfaction.

\emph{Demum quoad satisfactionem, quæ ex omnibus pœnitentiæ, partibus, quemadmodum a patribus nostris Christiano populo fuit perpetuo tempore commendata, ita una maxime nostra ætate summo pietatis prætextu impugnatur ab iis, qui speciem pietatis habent, virtutem autem ejus abnegarunt: sancta synodus declarat, falsum omnino esse et a verbo Dei alienum, culpam a Domino nunquam remitti, quin universa etiam pœna condonetur. Perspicua enim et illustria in sacris litteris exempla reperiuntur, quibus, præter divinam traditionem, hic error quam manifestissime revincitur. Sane et divinæ justitiæ ratio exigere videtur, ut aliter ab eo in gratiam recipiantur, qui ante baptismum per ignorantiam deliquerint; aliter vero qui semel a peccati et dæmonis servitute liberati, et accepto Spiritus Sancti dono, scientes templum Dei violare et Spiritum Sanctum contristare non formidaverint. Et divinam clementiam decet, ne ita nobis absque ulla }

Finally, as regards satisfaction,—which as it is, of all the parts of penance, that which has been at all times recommended to the Christian people by our Fathers, so is it the one especially which in our age is, under the loftiest pretext of piety, impugned by those who have \emph{an appearance of godliness, but have denied the power thereof},\footnote{\par   2 Tim. iii. 5.} —the holy Synod declares, that it is wholly false, and alien from the Word of God, that the guilt is never forgiven by the Lord, without the whole punishment also being therewith pardoned. For clear and illustrious examples are found in the sacred writings, whereby, besides by divine tradition, this error is refuted in the plainest manner possible. And truly the nature of divine justice seems to demand, that they, who through ignorance have sinned before baptism, be received into grace in one manner; and in another those who, after having been freed from the servitude of sin and of the devil, and after having received the gift of the Holy Ghost, have not feared, knowingly \emph{to violate the temple of God},\footnote{\par   1 Cor. iii. 17.} \emph{and to grieve the Holy Spirit.}\footnote{\par   Ephes. iv. 30.} And it

\emph{satisfactione peccata dimittantur, ut, occasione accepta, peccata leviora putantes, velut injurii et contumeliosi Spiritui Sancto in graviora labamur, thesaurizantes nobis iram in die iræ. Procul dubio enim magnopere a peccato revocant et quasi fræno quodam cærcent hæ satisfactoriæ pœnæ, cautioresque et vigilantiores in futurum pænitentes efficiunt; medentur quoque peccatorum reliquiis et vitiosos habitus male vivendo comparatos contrariis virtutum actionibus tollunt. Neque vero securior ulla via in Ecclesia Dei umquam existimata fuit ad amovendam imminentem a Domino pœnam, quam ut hæc pœnitentiæ opera homines cum vero animi dolore frequentent. Accedit ad hæc, quod, dum satisfaciendo patimur pro peccatis, Christo Jesu, qui pro peccatis nostris satisfecit, ex quo omnis nostra sufficientia est, conformes efficimur, certissimam quoque inde arrham habentes, quod, si compatimur et conglorificabimur. Neque vero ita nostra est satisfactio hæc, quam pro peccatis nostris exsolvimus, ut }

beseems the divine clemency, that sins be not in such wise pardoned us without any satisfaction, as that, taking occasion therefrom, thinking sins less grievous, we, offering as it were an insult and an \emph{outrage to the Holy Ghost},\footnote{\par   Heb. x. 29.} should fall into more grievous sins, \emph{treasuring up wrath against the day of wrath.}\footnote{\par   Rom. ii. 4.} For, doubtless, these satisfactory punishments greatly recall from sin, and check as it were with a bridle, and make penitents more cautious and watchful for the future; they are also remedies for the remains of sin, and, by acts of the opposite virtues, they remove the habits acquired by evil living. Neither indeed was there ever in the Church of God any way accounted surer to turn aside the impending chastisement of the Lord, than that men should, with true sorrow of mind, practice these works of penitence. Add to these things, that, whilst we thus, by making satisfaction, suffer for our sins, we are made conformable to Jesus Christ, who satisfied for our sins, from whom all our \emph{sufficiency is;}\footnote{\par   2 Cor. iii. 5.} having also thereby a most sure pledge, that \emph{if we suffer with him, we shall also be glorified with him.}\footnote{\par   Rom. viii. 17.} But neither is this satisfaction, which we

\emph{non sit per Christum Iesum, nam qui ex nobis, tamquam ex nobis, nihil possumus, eo cooperante, qui nos confortat, omnia possumus. Ita non habet homo, unde glorietur; sed omnis gloriatio nostra in Christo est; in quo vivimus, in quo meremur, in quo satisfacimus, facientes fructus dignos pœnitentiæ, qui ex illo vim habent, ab illo offeruntur Patri, et per illum acceptantur a Patre. Debent ergo sacerdotes Domini, quantum Spiritus et prudentia suggesserit, pro qualitate criminum et pœnitentium facultate, salutares et convenientes satisfactiones injungere; ne, si forte peccatis conniveant et indulgentius cum pœnitentibus agant, levissima quædam opera pro gravissimis delictis injungendo, alienorum peccatorum participes efficiantur. Habeant autem præ oculis, ut satisfactio, quam imponunt, non sit tantum ad novæ vitæ custodiam et infirmitatis medicamentum, sed etiam ad præteritorum peccatorum vindictam et castigationem: nam claves sacerdotum, non ad solvendum dumtaxat, sed et ad ligandum concessas etiam antiqui patres }

discharge for our sins, so our own, as not to be through Jesus Christ. For us who can do nothing of ourselves, as of ourselves, can do all things, he co-operating, who strengthens us. Thus, man has not wherein to glory, but all our glorying is in Christ: in whom we live; in whom we merit; in whom we satisfy; \emph{bringing forth fruits worthy of penance},\footnote{\par   Matt. iii. 18.} which from him have their efficacy; by him are offered to the Father; and through him are accepted by the Father. Therefore the priests of the Lord ought, as far as the Spirit and prudence shall suggest, to enjoin salutary and suitable satisfactions, according to the quality of the crimes and the ability of the penitent; lest, if haply they connive at sins, and deal too indulgently with penitents, by enjoining certain very light works for very grievous crimes, they be made partakers of other men's sins. But let them have in view, that the satisfaction, which they impose, be not only for the preservation of a new life and a medicine of infirmity, but also for the avenging and punishing of past sins. For the ancient Fathers likewise both believe and teach, that the keys of the priests were given, not to \emph{loose} only, but also to \emph{bind.}\footnote{\par   Matt. xvi. 19.;  John xx. 23.} But not therefore

\emph{et credunt et docent. Nec propterea existimarunt, sacramentum pœnitentiæ esse forum iræ vel pœnarum, sicut nemo umquam Catholicus sensit, ex hujusmodi nostris satisfactionibus vim meriti et satisfactionis Domini nostri Iesu Christi vel obscurari vel aliqua ex parte imminui; quod dum novatores intelligere nolunt, ita optimam pœnitentiam novam vitam esse docent, ut omnem satisfactionis vim et usum tollant.}

did they imagine that the sacrament of Penance is a tribunal of wrath or of punishments; even as no Catholic ever thought, that, by this kind of satisfaction on our parts, the efficacy of the merit and of the satisfaction of our Lord Jesus Christ is either obscured or in any way lessened: which when the innovators seek to understand, they in such wise maintain a new life to be the best penance, as to take away the entire efficacy and use of satisfaction.

Caput IX.

Chapter IX.

De operibus Satisfactionis.

On works of Satisfaction.

\emph{Docet præterea, tantam esse divinæ munificentiæ largitatem, ut non solum pænis sponte a nobis pro vindicando peccato susceptis, aut sacerdotis arbitrio pro mensura delicti impositis, sed etiam, quod maximum amoris argumentum est, temporalibus flagellis a Deo inflictis et a nobis patienter toleratis apud Deum Patrem per Christum Iesum satisfacere valeamus.}

The Synod teaches furthermore, that so great is the liberality of the divine munificence, that we are able through Jesus Christ to make satisfaction to God the Father, not only by punishments voluntarily undertaken of ourselves for the punishment of sin, or by those imposed at the discretion of the priest according to the measure of our delinquency, but also, which is a very great proof of love, by the temporal scourges inflicted of God, and borne patiently by us.

DOCTRINA DE SACRAMENTO EXTREMÆ UNCTIONIS.

ON THE SACRAMENT OF EXTREME UNCTION.

\emph{Visum est autem sanctæ synodo, præcedenti doctrinæ de pœnitentia adjungere ea, quæ sequuntur de sacramento }

It hath also seemed good to the holy Synod, to subjoin to the preceding doctrine on Penance, the following

\emph{extremæ unctionis, quod non modo pœnitentiæ, sed et totius Christianæ vitæ, quæ perpetua pœnitentia esse debet, consummativum existimatum est a Patribus. Primum itaque circa illius institutionem declarat et docet, quod clementissimus Redemptor noster, qui servis suis quovis tempore voluit de salutaribus remediis adversus omnia omnium hostium tela esse prospectum, quemadmodum auxilia maxima in sacramentis aliis præparavit, quibus Christiani conservare se integros, dum viverent, ab omni graviori spiritus incommodo possint: ita extremæ unctionis sacramento finem vitæ, tamquam firmissimo quodam præsidio, munivit. Nam etsi adversarius noster occasiones per omnem vitam quærat et captet, ut devorare animas nostras quoquo modo possit: nullum tamen tempus est, quo vehementius ille omnes suæ versutiæ nervos intendat ad perdendos nos penitus, et a fiducia etiam, si possit, divinæ misericordiæ deturbandos, quam cum impendere nobis exitum vitæ prospicit.}

on the sacrament of Extreme Unction, which by the Fathers was regarded as being the completion, not only of penance, but also of the whole Christian life, which ought to be a perpetual penance. First, therefore, as regards its institution, it declares and teaches, that our most gracious Redeemer,—who would have his servants at all times provided with salutary remedies against all the weapons of all their enemies,—as, in the other sacraments, he prepared the greatest aids, whereby during life, Christians may preserve themselves whole from every more grievous spiritual evil, so did he guard the close of life, by the sacrament of Extreme Unction, as with a most firm defense. For though \emph{our adversary} seeks and seizes opportunities, all our life long, to be able in any way \emph{to devour}\footnote{\par   1 Pet. v. 8.} our souls; yet is there no time wherein he strains more vehemently all the powers of his craft to ruin us utterly, and, if he can possibly, to make us fall even from trust in the mercy of God, than when he perceives the end of our life to be at hand.

Caput I.

Chapter I.

De institutione sacramenti Extremæ Unctionis.

On the institution of the sacrament of Extreme Unction.

\emph{Instituta est autem sacra unctio infirmorum tamquam vere et proprie sacramentum novi testamenti, a Christo Domino nos tro apud Marcum quidem insinuatum, per Iacobum autem apostolum ac Domini fratrem, fidelibus commendatum ac promulgatum. Infirmatur, inquit, quis in vobis? inducat presbyteros Ecclesiæ, et orent super eum, ungentes eum oleo in nomine Domini; et oratio fidei salvabit infirmum; et alleviabit eum Dominus; et si in peccatis sit, dimittentur ei. Quibus verbis, ut ex apostolica traditione per manus accepta Ecclesia didicit, docet materiam, formam, proprium ministrum, et effectum hujus salutaris sacramenti. Intellexit enim Ecclesia, materiam esse oleum ab episcopo benedictum; nam unctio aptissime Spiritus Sancti gratiam, qua invisibiliter anima ægrotantis inungitur, repræsentat; formam deinde esse illa verba: Per istam unctionem, etc.}

Now, this sacred unction of the sick was instituted by Christ our Lord, as truly and properly a sacrament of the new law, insinuated indeed in Mark, but recommended and promulgated to the faithful by James the Apostle, and brother of the Lord. \emph{Is any man}, he saith, \emph{sick among you? Let him, bring in the priests of the Church, and let them pray over him, anointing him with oil in the name of the Lord: and the prayer of faith shall save the sick man; and the Lord shall raise him up; and if he be in sins, they shall be forgiven him.}\footnote{\par   James v. 14, 15.} In which words, as the Church has learned from apostolic tradition, received from hand to hand, he teaches the matter, the form, the proper minister, and the effect of this salutary sacrament. For the Church has understood the matter thereof to be oil blessed by a bishop. For the unction very aptly represents the grace of the Holy Ghost, with which the soul of the sick person is invisibly anointed; and furthermore that those words, ``By this unction,'' etc., are the form.

Caput II.

Chapter II.

De effectu hujus Sacramenti.

On the effect of this Sacrament.

\emph{Res porro et effectus hujus sacramenti illis verbis explicatur: Et oratio fidei salvabit infirmum; et alleviabit eum Dominus; et si in peccatis sit, dimittentur ei. Res etenim hæc gratia est Spiritus Sancti, cujus unctio delicta, si quæ sint adhuc expianda, ac peccati reliquias abstergit; et ægroti animam alleviat et confirmat, magnam in eo divinæ misericordiæ fiduciam excitando; qua infirmus sublevatus et morbi incommoda ac labores levius fert, et tentationibus dæmonis, calcaneo insidiantis, facilius resistit, et sanitatem corporis interdum, ubi saluti animæ expedierit, consequitur.}

Moreover, the thing signified, and the effect of this sacrament, are explained in those words: \emph{And the prayer of faith shall save the sick man, and the Lord shall raise him up, and if he be in sins they shall be forgiven him.} For the thing here signified is the grace of the Holy Ghost; whose anointing cleanses away sins, if there be any still to be expiated, as also the remains of sins; \emph{and raises up} and strengthens the soul of the sick person, by exciting in him a great confidence in the divine mercy; whereby the sick being supported, bears more easily the inconveniences and pains of his sickness; and more readily resists the temptations of the devil who \emph{lies in wait for his heel;}\footnote{\par   Gen. iii. 15.} and at times obtains bodily health, when expedient for the welfare of the soul.

Caput III.

Chapter III.

De ministro hujus Sacramenti, et tempore, quo dari debeat.

On the minister of this Sacrament, and on the time when it ought to be administered.

\emph{Jam vero, quod attinet ad præscriptionem eorum, qui et suscipere et ministrare hoc sacramentum debent, haud obscure fuit illud etiam in verbis prædictis traditum. Nam et ostenditur illic, proprios hujus sacramenti }

And now as to prescribing who ought to receive, and who to administer this sacrament, this also was not obscurely delivered in the words above cited. For it is there also shown, that the proper ministers of this sacrament are the \emph{Presbyters}

\emph{ministros esse Ecclesia Presbyteros; quo nomine eo loco, non ætate seniores, aut primores in populo intelligendi veniunt, sed aut episcopi, aut sacerdotes ab ipsis rite ordinati per impositionem manuum presbyterii. Declaratur etiam, esse hanc unctionem infirmis adhibendam, illis vero præsertim, qui tam periculose decumbunt, ut in exitu vitæ constituti videantur; unde et sacramentum exeuntium nuncupatur. Quod si infirmi post susceptam hanc unctionem convaluerint, iterum hujus sacramenti subsidio juvari poterunt, cum in aliud simile vitæ discrimen inciderint. Quare nulla ratione audiendi sunt, qui contra tam apertam et dilucidam apostoli Iacobi sententiam docent, hanc unctionem vel figmentum esse humanum, vel ritum a patribus acceptum, nec mandatum Dei, nec promissionem gratiæ habentem; et qui illam jam cessasse asserunt, quasi ad gratiam curationum dumtaxat in primitiva Ecclesia referenda esset; et qui dicunt, ritum et usum, quem sancta Romana Ecclesia in hujus sacramenti administratione observat, Iacobi apostoli sententiæ repugnare, }

\emph{of the Church;} by which name are to be understood, in that place, not the elders by age, or the foremost in dignity amongst the people, but either bishops, or priests by bishops rightly ordained \emph{by the imposition of the hands of the priesthood.}\footnote{\par   1 Tim. iv. 14.} It is also declared, that this unction is to be applied to the sick, but to those especially who lie in such danger as to seem to be about to depart this life: whence also it is called the sacrament of the departing. And if the sick should, after having received this unction, recover, they may again be aided by the succor of this sacrament, when they fall into another like danger of death. Wherefore, they are on no account to be hearkened to, who, against so manifest and clear a sentence of the Apostle James, teach, either that this unction is a human figment or is a rite received from the Fathers, which neither has a command from God, nor a promise of grace: nor those who assert that it has already ceased, as though it were only to be referred to the grace of healing in the primitive Church; nor those who say that the rite and usage which the holy Roman Church observes in the administration of this sacrament is repugnant to the sentiment of the Apostle

\emph{atque ideo in alium commutandum esse; et denique, qui hanc extremam unctionem a fidelibus sine peccato contemni posse affirmant. Hæc enim omnia manifestissime pugnant cum perspicuis tanti apostoli verbis. Nec profecto Ecelesia Romana, aliarum omnium mater et magistra, aliud in hac administranda unctione, quantum ad ea, quæ hujus sacramenti substantiam perficiunt, observat, quam quod beatus Iacobus præscripsit. Neque vero tanti sacramenti contemptus absque ingenti scelere et ipsius Spiritus Sancti injuria esse posset.}

James, and that it is therefore to be changed into some other; nor finally those who affirm that this Extreme Unction may without sin be contemned by the faithful; for all these things are most manifestly at variance with the perspicuous words of so great an apostle. Neither assuredly does the Roman Church, the mother and mistress of all other churches, observe aught in administering this unction,—as regards those things which constitute the substance of this sacrament,—but what blessed James has prescribed. Nor indeed can there be contempt of so great a sacrament without a heinous sin, and an injury to the Holy Ghost himself.

\emph{Hæc sunt, quæ de pœnitentiæ et extremæ unctionis sacramentis sancta hæc œcumenica synodus profitetur et docet atque omnibus Christi fidelibus credenda et tenenda proponit. Sequentes autem canones inviolabiliter servandos esse tradit, et asserentes contrarium perpetuo damnat et anathematizat.}

These are the things which this holy œcumenical Synod professes and teaches and proposes to all the faithful of Christ, to be believed and held, touching the sacraments of Penance and Extreme Unction. And it delivers the following canons to be inviolably preserved; and condemns and anathematizes those who assert what is contrary thereto.

DE SANCTISSIMO PŒNITENTIÆ SACRAMENTO.

ON THE MOST HOLY SACRAMENT OF PENANCE.

Canon I.—\emph{Si quis dixerit, in Catholica Ecclesia pœnitentiam non esse vere et proprie sacramentum pro fidelibus, quoties }

Canon I.—If any one saith, that in the Catholic Church Penance is not truly and properly a sacrament, instituted by Christ our Lord

\emph{post baptismum in peccata labuntur, ipsi Deo reconciliandis a Christo Domino nostro institutum: anathema sit.}

for reconciling the faithful unto God, as often as they fall into sin after baptism: let him be anathema.

Canon II.—\emph{Si quis sacramenta confundens, ipsum baptismum pœnitentiæ sacramentum esse dixerit, quasi hæc duo sacramenta distincta non sint, atque ideo pœnitentiam non recte secundam post naufragium tabulam appellari: anathema sit.}

Canon II.—If any one, confounding the sacraments, saith that baptism is itself the sacrament of Penance, as though these two sacraments were not distinct, and that therefore Penance is not rightly called a second plank after shipwreck: let him be anathema.

Canon III.—\emph{Si quis dixerit, verba illa Domini Salvatoris: Accipite Spiritum Sanctum; quorum remiseritis peccata, remittuntur eis; et quorum retinueritis, retenta sunt: non esse intelligenda de potestate remittendi et retinendi peccata in sacramento pœnitentiæ, sicut Ecclesia Catholica ab initio semper intellexit; detorserit autem, contra institutionem hujus sacramenti, ad auctoritatem prædicandi evangelium: anathema sit.}

Canon III.—If any one saith, that those words of the Lord the Saviour, \emph{Receive ye the Holy Ghost, whose sins you shall forgive, they are forgiven them, and whose sins you shall retain, they are retained},\footnote{\par   John xx. 22, 23.} are not to be understood of the power of forgiving and of retaining sins in the sacrament of Penance, as the Catholic Church has always from the beginning understood them; but wrests them, contrary to the institution of this sacrament, to the power of preaching the gospel: let him be anathema.

Canon IV.—\emph{Si quis negaveret, ad integram et perfectam peccatorum remissionem requiri tres actus in pœnitente, quasi materiam sacramenti pœnitentiæ, videlicet, contritionem, confessionem, et satisfactionem, quæ tres pœnitentiæ}

Canon IV.—If any one denieth, that, for the entire and perfect remission of sins, there are required three acts in the penitent, which are as it were the matter of the sacrament of Penance, to wit, contrition, confession, and satisfaction, which are called the three parts of

\emph{partes dicuntur; aut dixerit, duas tantum esse pœnitentiæ partes, terrores scilicet incussos conscientiæ, agnito peccato, et fidem conceptam ex evangelio vel absolutione, qua credit quis sibi per Christum remissa peccata: anathema sit.}

penance; or saith that there are two parts only of penance, to wit, the terrors with which the conscience is smitten upon being convinced of sin, and the faith, generated by the gospel, or by the absolution, whereby one believes that his sins are forgiven him through Christ: let him be anathema.

Canon V.—\emph{Si quis dixerit eam contritionem, quæ paratur per discussionem, collectionem et detestationem peccatorum, qua quis recogitat annos suos in amaritudine animæ suæ, ponderando peccatorum suorum gravitatem, multitudinem, fœditatem, amissionem æternæ beatitudinis, et æternæ damnationis incursum, cum proposito melioris vitæ, non esse verum et utilem dolorem, nec præparare ad gratiam, sed facere hominem hypocritam et magis peccatorem; demum, illum esse dolorem coactum et non liberum ac voluntarium: anathema sit.}

Canon V.—If any one saith, that the contrition which is acquired by means of the examination, collection, and detestation of sins,—whereby one \emph{thinks over his years in the bitterness of his soul},\footnote{\par   Isa. xxxviii. 15.} by pondering on the grievousness, the multitude, the filthiness of his sins, the loss of eternal blessedness, and the eternal damnation which he has incurred, having therewith the purpose of a better life,—is not a true and profitable sorrow, does not prepare for grace, but makes a man a hypocrite and a greater sinner; in fine, that this [contrition] is a forced and not free and voluntary sorrow: let him be anathema.

Canon VI.—\emph{Si quis negaverit, confessionem sacramentalem vel institutam, vel ad salutem necessariam esse jure divino; aut dixerit, modum secrete confitendi soli sacerdoti, quem Ecclesia Cathotica ab initio semper observavit et observat, alienum }

Canon VI.—If any one denieth, either that sacramental confession was instituted, or is necessary to salvation, of divine right; or saith, that the manner of confessing secretly to a priest alone, which the Church hath ever observed from the beginning, and doth observe, is alien

\emph{esse ab institutione et mandato Christi, et inventum esse humanum: anathema sit.}

from the institution and command of Christ, and is a human invention: let him be anathema.

Canon VII.—\emph{Si quis dixerit, in sacramento pœnitentiæ ad remissionem peccatorum necessarium non esse jure divino confiteri omnia et singula peccata mortalia, quorum memoria cum debita et diligenti præmeditatione habeatur, etiam occulta, et quæ sunt contra duo ultima Decalogi præcepta, et circumstantias, quæ peccati speciem mutant, sed eam confessionem tantum esse utilem ad erudiendum et consolandum pœnitentem, et olim observatam fuisse tantum ad satisfactionem canonicam imponendam; aut dixerit eos, qui omnia peccata confiteri student, nihil relinquere velle divinæ misericordiæ ignoscendum; aut demum, non licere confiteri peccata venialia: anathema sit.}

Canon VII.—If any one saith, that, in the sacrament of Penance, it is not necessary, of divine right, for the remission of sins, to confess all and singular the mortal sins which after due and diligent previous meditation are remembered, even those [mortal sins] which are secret, and those which are opposed to the two last commandments of the Decalogue, as also the circumstances which change the species of a sin; but [saith] that such confession is only useful to instruct and console the penitent, and that it was of old only observed in order to impose a canonical satisfaction; or saith that they, who strive to confess all their sins, wish to leave nothing to the divine mercy to pardon; or, finally, that it is not lawful to confess venial sins: let him be anathema.

Canon VIII.—\emph{Si quis dixerit, confessionem omnium peccatorum, qualem Ecclesia servat, esse impossibilem et traditionem humanam a piis abolendam; aut ad eam non teneri omnes et singulos utriusque Christi fideles, juxta magni Concilii Lateranensis constitutionem, semel in anno et }

Canon VIII.—If any one saith, that the confession of all sins, such as it is observed in the Church, is impossible, and is a human tradition to be abolished by the godly; or that all and each of the faithful of Christ, of either sex, are not obliged thereunto once a year, conformably to the constitution of the great Council of Lateran, and that,

\emph{ob id suadendum esse Christi fidelibus, ut non confiteantur tempore quadragesimæ: anathema sit.}

for this cause, the faithful of Christ are to be persuaded not to confess during Lent: let him be anathema.

Canon IX.—\emph{Si quis dixerit, absolutionem sacramentalem sacerdotis, non esse actum judicialem, sed nudum ministerium pronunciandi et declarandi, remissa esse peccata confitenti, modo tantum credat, se esse absolutum; aut sacerdos non serio, sed joco absolvat; aut dixerit, non requiri confessionem pœnitentis, ut sacerdos ipsum absolvere possit: anathema sit.}

Canon IX.—If any one saith, that the sacramental absolution of the priest is not a judicial act, but a bare ministry of pronouncing and declaring sins to be forgiven to him who confesses; provided only he believe himself to be absolved, or [even though] the priest absolve not in earnest, but in joke; or saith, that the confession of the penitent is not required, in order that the priest may be able to absolve him: let him be anathema.

Canon X.—\emph{Si quis dixerit, sacerdotes, qui in peccato mortali sunt, potestatem ligandi et solvendi non habere; aut non solos sacerdotes esse ministros absolutionis, sed omnibus et singulis Christi fidelibus esse dictum: Quæcumque ligaveritis super terram, erunt ligata et in cœlo; et quæcumque solveritis super terram, erunt soluta et in cælo; et: Quorum remiseritis peccata, remittuntur eis; et quorum retinueritis, retenta sunt: quorum verborum virtute quilibet absolvere possit peccata, publica quidem per correptionem dumtaxat, si correptus acquieverit, }

Canon X.—If any one saith, that priests, who are in mortal sin, have not the power of binding and loosing; or, that not priests alone are the ministers of absolution, but that, to all and each of the faithful of Christ is it said: \emph{Whatsoever you shall bind upon earth shall be bound also in heaven; and whatsoever you shall loose upon earth, shall be loosed also in heaven;}\footnote{\par   Matt. xviii. 15.} \emph{and, whose sins you shall forgive, they are forgiven them; and whose sins you shall retain, they are retained;}\footnote{\par   John xx. 23.} by virtue of which words every one is able to absolve from sins, to wit, from public sins by reproof only, provided he who is

\emph{secreta vero per spontaneam confessionem: anathema sit.}

reproved yield thereto, and from secret sins by a voluntary confession: let him be anathema.

Canon XI.—\emph{Si quis dixerit, episcopos non habere jus reservandi sibi casus, nisi quoad externam politiam, atque ideo casuum reservationem non prohibere, quo minus sacerdos a reservatis vere absolvat: anathema sit.}

Canon XI.—If any one saith, that bishops have not the right of reserving cases to themselves, except as regards external polity, and that therefore the reservation of cases hinders not, but that a priest may truly absolve from reserved cases: let him be anathema.

Canon XII.—\emph{Si quis dixerit, totam pœnam simul cum culpa remitti semper a Deo, satisfactionemque pœnitentium non esse aliam quam fidem, qua apprehendunt Christum pro eis satisfecisse: anathema sit.}

Canon XII.—If any one saith, that God always remits the whole punishment together with the guilt, and that the satisfaction of penitents is no other than the faith whereby they apprehend that Christ has satisfied for them: let him be anathema.

Canon XIII.—\emph{Si quis dixerit, pro peccatis, quoad pœnam temporalem, minime Deo pro Christi merita satisfieri pœnis ab eo inflictis et patienter toleratis, vel a sacerdote injunctis, sed neque sponte susceptis, ut jejuniis, orationibus, eleemosynis, vel aliis etiam pietatis operibus, atque ideo optimam pœnitentiam esse tantum novam vitam: anathema sit.}

Canon XIII.—If any one saith, that satisfaction for sins, as to their temporal punishment, is nowise made to God, through the merits of Jesus Christ, by the punishments inflicted by him, and patiently borne, or by those enjoined by the priest, nor even by those voluntarily undertaken, as by fastings, prayers, alms-deeds, or by other works also of piety; and that, therefore, the best penance is merely a new life: let him be anathema.

Canon XIV.—\emph{Si quis dixerit, satisfactiones, quibus pœnitentes per Christum Iesum peccata redimunt, non esse cultus }

Canon XIV.—If any one saith, that the satisfactions, by which penitents redeem their sins through Jesus Christ, are not a worship of

\emph{Dei, sed traditiones hominum, doctrinam de gratia, et verum Dei cultum atque ipsum beneficium mortis Christi obscurantes: anathema sit.}

God, but traditions of men, which obscure the doctrine of grace, and the true worship of God, and the benefit itself of the death of Christ: let him be anathema.

Canon XV.—\emph{Si quis dixerit, claves Ecclesiæ esse datas tantum ad solvendum, non etiam ad ligandum, et propterea sacerdotes, dum imponunt pœnas confitentibus, agere contra finem clavium et contra institutionem Christi; et fictionem esse, quod, virtute clavium sublata pœna æterna, pœna temporalis plerumque exsolvenda remaneat: anathema sit.}

Canon XV.—If any one saith, that the keys are given to the Church, only \emph{to loose}, not also \emph{to bind;} and that, therefore, priests act contrary to the purpose of the keys, and contrary to the institution of Christ, when they impose punishments on those who confess; and that it is a fiction, that, after the eternal punishment has, by virtue of the keys, been removed, there remains for the most part a temporal punishment to be discharged: let him be anathema.

DE SACRAMENTO EXTREMÆ UNCTIONIS.

ON THE SACRAMENT OF EXTREME UNCTION.

Canon I.—\emph{Si quis dixerit, extremam unctionem non esse vere et proprie sacramentum a Christo domino nostro institutum et a beato Iacobo apostolo promulgatum; sed ritum tantum acceptum a patribus aut figmentum humanum: anathema sit.}

Canon I.—If any one saith, that Extreme Unction is not truly and properly a sacrament, instituted by Christ our Lord, and promulgated by the blessed Apostle James; but is only a rite received from the Fathers, or a human figment: let him be anathema.

Canon II.—\emph{Si quis dixerit, sacram infirmorum unctionem non conferre gratiam, nec remittere peccata, nec alleviare infirmos, sed jam cessasse, quasi olim tantum fuerit gratia curationum: anathema sit.}

Canon II.—If any one saith, that the sacred unction of the sick does not confer grace, nor remit sin, nor comfort the sick; but that it has already ceased, as though it were of old only the grace of working cures: let him be anathema.

Canon III.—\emph{Si quis dixerit, extremæ unctionis ritum et usum, quem observat sancta Romana Ecclesia, repugnare sententiæ beati Iacobi apostoli, ideoque eum mutandum, posseque a Christianis absque peccato contemni: anathema sit.}

Canon III.—If any one saith, that the right and usage of Extreme Unction, which the holy Roman Church observes, is repugnant to the sentiment of the blessed Apostle James, and that is therefore to be changed, and may, without sin, be contemned by Christians: let him be anathema.

Canon IV.—\emph{Si quis dixerit, Presbyteros Ecclesiæ, quos beatus Iacobus adducendos esse ad infirmum, inungendum hortatur, non esse sacerdotes ab episcopo ordinatos, sed ætate seniores in quavis communitate, ob idque proprium extremæ unctionis ministrum non esse solum sacerdotem: anathema sit.}

Canon IV.—If any one saith, that the \emph{Presbyters of the Church}, whom blessed James exhorts to be brought to anoint the sick, are not the priests who have been ordained by a bishop, but the elders in each community, and that for this cause a priest alone is not the proper minister of Extreme Unction: let him be anathema.

\xchapter{Twenty-first Session, held July 16, 1562.}

Sessio Vigesimaprima,

Twenty-first Session,

\emph{celebrata die XVI. Iulii} 1562.

\emph{held July} 16, 1562.

DOCTRINA DE COMMUNIONE SUB UTRAQUE SPECIE, ET PARVULORUM.

DOCTRINE CONCERNING THE COMMUNION UNDER BOTH SPECIES, AND OF LITTLE CHILDREN.

Caput I.

Chapter I.

\emph{Laicos et clericos non conficientes non adstringi jure divino ad communionem sub utraque specie.}

\emph{That laymen and clerics, when not sacrificing, are not bound, of divine right, to communion under both species.}

\emph{Itaque sancta ipsa synodus, a Spiritu Sancto, qui spiritus est sapientiæ et intellectus, spiritus consilii et pietatis, edocta, atque ipsius Ecclesiæ judicium }

Wherefore, this holy Synod,—instructed by the Holy \emph{Spirit}, who is \emph{the spirit of wisdom and of understanding, the spirit of counsel and of godliness},\footnote{\par   Isa. xi. 2.} and following the

\emph{et consuetudinem secuta, declarat, ac docet, nullo divino præcepto laicos et clericos non conficientes, obligari ad Eucharistiæ sacramentum sub utraque specie sumendum; neque ullo pacto, salva fide, dubitari posse, quin illis alterius speciei communio ad salutem sufficiat: nam, etsi Christus Dominus in ultima cæna venerabile hoc sacramentum in panis, et vini speciebus instituit et apostolis tradidit; non tamen illa institutio et traditio eo tendunt, ut omnes Christi fideles statuto Domini ad utramque speciem accipiendam adstringantur. Sed neque ex sermone illo, apud Ioannem VI., recte colligitur, utriusqme speciei communionem a Domino præceptam esse: utcumque juxta varias sanctorum patrum et doctorum interpretationes intelligatur: namque, qui dixit: Nisi manducaveritis carnem filii hominis et biberitis ejus sanguinem, non habebitis vitam in vobis: dixit quoque: Si quis manducaverit ex hoc pane, vivet in æternum. Et qui dixit: Qui manducat meam carnem, et bibit meum sanguinem, habet vitam æternam: dixit etiam: Panis, quem ego dabo, caro mea est pro mundi }

judgment and usage of the Church itself,—declares and teaches, that laymen, and clerics when not consecrating, are not obliged, by any divine precept, to receive the sacrament of the Eucharist under both species; and that neither can it by any means be doubted, without injury to faith, that communion under either species is sufficient for them unto salvation. For, although Christ, the Lord, in the Last Supper, instituted and delivered to the apostles, this venerable sacrament in the species of bread and wine; not therefore do that institution and delivery tend thereunto, that all the faithful of the Church be bound, by the institution of the Lord, to receive both species. But neither is it rightly gathered, from that discourse which is in the sixth of John,—however according to the various interpretations of holy Fathers and Doctors it be understood,—that the communion of both species was enjoined by the Lord; for he who said, \emph{Except you eat the flesh of the Son of man and drink his blood, you shall not have life in you} (v. 54), also said: \emph{He that eateth this bread shall live forever} (v. 59); and he who said, \emph{He that eateth my flesh and drinketh my blood hath everlasting life} (v. 55), also said: \emph{The bread that I will}

\emph{vita. Et denique qui dixit: Qui manducat meam carnem et bibit meum sanguinem, in me manet et ego in illo: dixit nihilominus: Qui manducat hunc panem, vivet in æternum.}

\emph{give is my flesh for the life of the world} (v. 52); and, in fine, he who said, \emph{He that eateth my flesh and drinkeih my blood, abideth in me and I in him} (v. 57), said, nevertheless, \emph{He that eateth this bread shall live forever} (v. 59).

Caput II.

Chapter II.

\emph{Ecclesiæ potestas circa dispensationem sacramenti Eucharistiæ.}

\emph{The power of the Church as regards the dispensation of the Sacrament of the Eucharist.}

\emph{Præterea declarat, hanc potestatem perpetuo in Ecclesia fuisse, ut in sacramentorum dispensatione, salva illorum substantia, ea statueret vel mutaret, quæ suscipientium utilitati seu ipsorum sacramentorum venerationi, pro rerum, temporum et locorum veritate, magis expedire judicaret. Id autem apostolus non obscure visus est innuisse, cum ait: Sic nos existimet homo, ut ministros Christi et dispensatores mysteriorum Dei; atque ipsum quidem hac potestate usum esse satis constat cum in multis aliis, tum in hoc ipso sacramento, cum, ordinatis nonnullis circa ejus usum, Cetera, inquit, cum venero, disponam. Quare agnoscens sancta mater Ecclesia hanc suam in administratione sacramentorum auctoritatem, licet ab }

It furthermore declares, that this power has ever been in the Church, that, in the dispensation of the sacraments, their substance being untouched, it may ordain, or change, what things soever it may judge most expedient, for the profit of those who receive, or for the veneration of the said sacraments, according to the difference of circumstances, times, and places. And this the Apostle seems not obscurely to have intimated, when he says: \emph{Let a man so account of us, as of the ministers of Christ, and the dispensers of the mysteries of God.}\footnote{\par   1 Cor. iv. i.} And, indeed, it is sufficiently manifest that he himself exercised this power, as in many other things, so in regard of this very sacrament; when, after having ordained certain things touching the use thereof, he says: \emph{The rest I will set in order when I come.}\footnote{\par   1 Cor. xi. 34.} Wherefore, holy

\emph{initio Christianæ religionis non infrequens utriusque speciei usus fuisset, tamen progressu temporis, latissime jam mutata illa consuetudine, gravibus et justis causis adducta hanc consuetudinem sub altera specie communicandi approbabit, et pro lege habendam decrevit, quam reprobare aut sine ipsius Ecclesiæ auctoritate pro libito mutare non licet.}

Mother Church, knowing this her authority in the administration of the sacraments, although the use or both species has, from the beginning of the Christian religion, not been unfrequent, yet, in progress of time, that custom having been already very widely changed, she, induced by weighty and just reasons, has approved of this custom of communicating under one species, and decreed that it was to be held as a law; which it is not lawful to reprobate, or to change at pleasure, without the authority of the Church itself.

Caput III.

Chapter III.

\emph{Totum et integrum Christum ac verum sacramentum sub qualibet specie sumi.}

\emph{That Christ whole and entire and a true Sacrament are received under either species.}

\emph{Insuper declarat, quamvis Redemptor noster, ut antea dictum est, in suprema illa cæna hoc sacramentum in duabus speciebus instituerit et apostolis tradiderit, tamen fatendum esse, etiam sub altera tantum specie totum atque integrum Christum verumque sacramentum sumi; ac propterea, quod ad fructum attinet nulla gratia necessaria ad salutem eos defraudari, qui unam speciam solam accipiunt.}

It moreover declares, that although, as hath been already said, our Redeemer, in that last supper, instituted, and delivered to the apostles, this sacrament in two species, yet is to be acknowledged, that Christ whole and entire and a true sacrament are received under either species alone; and that therefore, as regards the fruit thereof, they, who receive one species alone are not defrauded of any grace necessary to salvation.

Caput IV.

Chapter IV.

\emph{Parvulos non obligari ad communionem sacramentalem.}

\emph{That little Children are not bound to sacramental Communion.}

\emph{Denique eadem sancta synodus docet, parvulos usu rationis carentes nulla obligari necessitate ad sacramentalem Eucharistiæ communionem, siquidem, per baptismi lavacrum regenerati et Christo incorporati, adeptam jam filiorum Dei gratiam in illa ætate amittere non possunt. Neque ideo tamen damnanda est antiquitas, si eum morem in quibusdam locis aliquando servavit. Ut enim sanctissimi illi patres sui facti probabilem causam pro illius temporis ratione habuerunt, ita certe eos nulla salutis necessitate id fecisse sine controversia credendum est.}

Finally, this same holy Synod teaches, that little children, who have not attained to the use of reason, are not by any necessity obliged to the sacramental communion of the Eucharist: forasmuch as, having been regenerated by the laver of baptism, and being incorporated with Christ, they can not, at that age, lose the grace which they have already acquired of being the sons of God. Not therefore, however, is antiquity to be condemned, if, in some places, it, at one time, observed that custom; for as those most holy Fathers had a probable cause for what they did in respect of their times, so, assuredly, is it to be believed without controversy, that they did this without any necessity thereof unto salvation.

DE COMMUNIONE SUB UTRAQUE SPECIE ET PARVULORUM.

ON COMMUNION UNDER BOTH SPECIES, AND ON THE COMMUNION OF INFANTS.

Canon I.—\emph{Si quis dixerit, ex Dei præcepto vel necessitate salutis omnes et singulos Christi fideles utramque speciem sanctissimi Eucharistiæ sacramenti sumere debere: anathema sit.}

Canon I.—If any one saith, that, by the precept of God, or by necessity of salvation, all and each of the faithful of Christ ought to receive both species of the most holy sacrament of the Eucharist: let him be anathema.

Canon II.—\emph{Si quis dixerit, sanctum Ecclesiam Catholicam non justis }

Canon II.—If any one saith, that the holy Catholic Church was not

\emph{causis et rationibus adductam fuisse, ut laicos atque etiam clericos non conficientes sub panis tantummodo specie communicaret, aut in eo errasse: anathema sit.}

induced, by just causes and reasons, to communicate, under the species of bread only, laymen, and also clerics when not consecrating: let him be anathema.

Canon III.—\emph{Si quis negaverit, totum et integrum Christum, omnium gratiarum fontem et auctorem, sub una panis specie sumi, quia, ut quidam falso asserunt, non secundum ipsius Christi institutionem sub utraque specie sumatur: anathema sit.}

Canon III.—If any one denieth, that Christ whole and entire,—the fountain and author of all graces,—is received under the one species of bread; because that, as some falsely assert, he is not received, according to the institution of Christ himself, under both species: let him be anathema.

Canon IV.—\emph{Si quis dixerit, parvulis, antequam ad annos discretionis pervenerint, necessariam esse Eucharistiæ communionem: anathema sit.}

Canon IV.—If any one saith, that the communion of the Eucharist is necessary for little children, before they have arrived at years of discretion: let him be anathema.

\emph{Duos vero articulos alias propositos nondum tamen excussos, videlicet: an rationes, quibus sancta Catholica Ecclesia adducta fuit, ut communicaret laicos atque etiam non celebrantes sacerdotes, sub una tantum panis specie, ita sint retinendæ, ut nulla ratione calicis usus cuiquam sit permittendus; et: an, si honestis et Christianæ caritati consentaneis rationibus concedendus alicui vel nationi vel regno calicis usus videatur, sub aliquibus conditionibus concedendus sit, et quænam sint illæ, eadem }

As regards, however, those two articles, proposed on another occasion, but which have not as yet been discussed: to wit, whether the reasons by which the holy Catholic Church was led to communicate, under the one species of bread only, laymen, and also priests when not celebrating, are in such wise to be adhered to, as that on no account is the use of the chalice to be allowed to any one soever; and whether, in case that, for reasons beseeming and consonant with Christian charity, it appears that the use of the chalice is to be granted to any nation or kingdom, it is to be conceded

\emph{sancta synodus in aliud tempus, oblata sibi quamprimum occasione, examinandos atque definiendos reservat.}

under certain conditions; and what are those conditions: this same holy Synod reserves the same to another time,—for the earliest opportunity that shall present itself,—to be examined and defined.

\xchapter{Twenty-second Session, held Sept. 17, 1562.}

Sessio Vigesimasecunda,

Twenty-second Session,

\emph{celebrata die XVII. Sept.} 1562.

\emph{held Sept.} 17, 1562.

DOCTRINA DE SACRIFICIO MISSÆ

DOCTRINE ON THE SACRIFICE OF THE MASS.

Caput I.

Chapter I.

\emph{De institutione sacrosancti missæ sacrificii.}

\emph{On the institution of the most holy Sacrifice of the Mass.}

\emph{Quoniam sub priori Testamento, teste Apostolo Paulo, propter Levitici sacerdotii imbecillitatem consummatio non erat, oportuit, Deo patre misericordiarum ita ordinante, sacerdotem alium secundum ordinem Melchisedech surgere, Dominum nostrum Iesum Christum, qui posset omnes, quotquot sanctificandi essent, consummare, et ad perfectum adducere. Is igitur Deus et Dominus noster, etsi semel se ipsum in ara crucis, morte intercedente, Deo patri oblaturus erat, ut æternam illic redemptionem operaretur, quia tamen per mortem sacerdotium ejus }

Forasmuch as, under the former Testament, according to the testimony of the Apostle Paul, there was no \emph{perfection, because of the weakness of the Levitical priesthood;}\footnote{\par   Heb. vii. 11, 18.} there was need, God, the Father of mercies, so ordaining, that \emph{another priest should rise, according to the order of Melchisedech},\footnote{\par   Heb. v. 10.} our Lord Jesus Christ, who might consummate, and lead to what is perfect, as many as were to be sanctified. He, therefore, our God and Lord, though he was about to offer himself once on the altar of the cross unto God the Father, \emph{by means of his death}, there to operate \emph{an eternal redemption;}\footnote{\par   Heb. ix. 12.} nevertheless, because that his priesthood was not

\emph{extinguendum non erat, in cæna novissima, qua nocte tradebatur, ut dilectæ sponsæ suæ Ecclesiæ visibile, sicut hominum natura exigit, relinqueret sacrificium, quo cruentum illud semel in cruce peragendum repræsentaretur, ejusque memoria in finem usque sæculi permaneret, atque illius salutaris virtus in remissionem eorum, quæ a nobis quotidie committuntur, peccatorum applicaretur, sacerdotem secundum ordinem Melchisedech se in æternum constitutum declarans, corpus et sanguinem suum sub speciebus panis et vini Deo Patri obtulit, ac sub earumdem rerum symbolis apostolis, quos tunc Novi Testamenti sacerdotes constituebat, ut sumerent, tradidit, et eisdem eorumque in sacerdotio successoribus, ut offerrent, præcepit per hæc verba: Hoc facite in meam commemorationem: uti semper Catholica Ecclesia intellexit et docuit. Nam celebrato veteri Pascha, quod in memoriam exitus de Aegypto multitudo filiorum Israel immolabat, novum instituit Pascha se ipsum ab Ecclesia per sacerdotes sub signis visibilibus immolandum in memoriam }

to be extinguished by his death, in the Last Supper, on the night in which he was betrayed,—that he might leave, to his own beloved Spouse the Church, a visible sacrifice, such as the nature of man requires, whereby that bloody sacrifice, once to be accomplished on the cross, might be represented, and the memory thereof remain even unto the end of the world, and its salutary virtue be applied to the remission of those sins which we daily commit,—declaring himself constituted \emph{a priest forever, according to the order of Melchisedech},\footnote{\par   Psa. cix. 4.} he offered up to God the Father his own body and blood under the species of bread and wine; and, under the symbols of those same things, he delivered [his own body and blood] to be received by his apostles, whom he then constituted priests of the New Testament; and by those words, \emph{Do this in commemoration of me},\footnote{\par   Luke xxii. 19.} he commanded them and their successors in the priesthood to offer [them]; even as the Catholic Church has always understood and taught. For, having celebrated the ancient Passover, which the multitude of the children of Israel immolated in memory of their going out of Egypt, he instituted the new Passover [to wit], himself to

\emph{transitus sui ex hoc mundo ad Patrem, quando per sui sanguinis effusionem nos redemit eripuitque de potestate tenebrarum, et in regnum suum transtulit. Et hæc quidem illa munda oblatio est, quæ nulla indignitate aut malitia offerentium inquinari potest; quam Dominus per Malachiam nomini suo, quod magnum futurum esset in gentibus, in omni loco mundam offerendam prædixit, et quam non obscure innuit Apostolus Paulus Corinthiis scribens, cum dicit, non posse eos, qui participatione mensæ dæmoniorum polluti sint, mensæ Domini participes fieri, per mensam altare utrobique intelligens. Hæc denique illa est, quæ per varias sacrificiorum, naturæ et legis tempore, similitudines figurabatur; utpote quæ bona omnia, per illa significata, velut illorum omnium consummatio et perfectio complectitur.}

be immolated, under visible signs, by the Church through [the ministry of] priests, in memory of his own passage from this world unto the Father, when by the effusion of his own blood he redeemed us, \emph{and delivered us from the power of darkness, and translated us into his kingdom.}\footnote{\par   Col. i. 13.} And this is indeed that clean oblation, which can not be defiled by any unworthiness, or malice of those that offer [it]; which the Lord foretold by Malachias was to be \emph{offered in every place, clean to his name, which was to be great amongst the Gentiles;}\footnote{\par   Mal. i. 11.} and which the Apostle Paul, writing to the Corinthians, has not obscurely indicated, when he says, that they who are defiled by \emph{the participation of the table of devils, can not be partakers of the table of the Lord;}\footnote{\par   1 Cor. x. 20 sqq.} by \emph{the table}, meaning in both places the altar. This, in fine, is that oblation which was prefigured by various types of sacrifices, during the period of nature, and of the law; inasmuch as it comprises all the good things signified by those sacrifices, as being the consummation and perfection of them all.

Caput II.

Chapter II.

\emph{Sacrificium missæ est propitiatorium, tam pro vivis, quam pro defunctis.}

\emph{That the Sacrifice of the Mass is propitiatory, both for the living and the dead.}

\emph{Et quoniam in divino hoc sacrificio, quod in missa peragitur, idem ille Christus continetur et incruente immolatur, qui in ara crucis semel se ipsum cruente obtulit, docet sancta synodus, sacrificium istud vere propitiatorium esse, per ipsumque fieri, ut, si cum vero corde et recta fide, cum metu et reverentia, contriti ac pœnitentes ad Deum accedamus, misericordiam consequamur et gratiam inveniamus in auxilio opportuno. Hujus quippe oblatione placatus Dominus gratiam et donum pœnitentiæ concedens, crimina et peccata etiam ingentia dimittit. Una enim eademque est hostia, idem nunc offerens sacerdotum ministerio, qui se ipsum tunc in cruce obtulit, sola offerendi ratione diversa. Cujus quidem oblationis cruentæ, inquam, fructus per hanc incruentam uberrime percipiuntur, tantum abest, ut illi per hanc quovis modo derogetur. Quare non solum pro fidelium vivorum peccatis, pœnis, satisfactionibus et aliis necessitatibus, sed pro defunctis }

And forasmuch as, in this divine sacrifice which is celebrated in the mass, that same Christ is contained and immolated in an unbloody manner who once offered himself in a bloody manner on the altar of the cross; the holy Synod teaches, that this sacrifice is truly propitiatory, and that by means thereof this is effected, that we obtain mercy, and find grace \emph{in seasonable aid},\footnote{\par   Heb. iv. 6.} if we draw nigh unto God, contrite and penitent, with a sincere heart and upright faith, with fear and reverence. For the Lord, appeased by the oblation thereof, and granting the grace and gift of penitence, forgives even heinous crimes and sins. For the victim is one and the same, the same now offering by the ministry of priests, who then offered himself on the cross, the manner alone of offering being different. The fruits indeed of which oblation, of that bloody one to wit, are received most plentifully through this unbloody one; so far is this [latter] from derogating in any way from, that [former oblation]. Wherefore, not only for the sins, punishments, satisfactions, and other necessities of the faithful who are living, but

\emph{in Christo nondum ad plenum purgatis rite juxta apostolorum traditionem offertur.}

also for those who are departed in Christ, and who are not as yet fully purified, is it rightly offered, agreeably to a tradition of the apostles.

Caput III.

Chapter III.

\emph{De missa in honorem sanctorum.}

\emph{On Masses in honor of the Saints.}

\emph{Et quamvis in honorem et memoriam sanctorum nonnullus interdum missas Ecclesia celebrare consueverit, non tamen illis sacrificium offerri docet, sed Deo soli, qui illos coronavit; unde nec sacerdos dicere solet: Offero tibi sacrificium, Petre vel Paule; sed, Deo de illorum victoriis gratias agens, eorum patrocinia implorat, ut ipsi pro nobis intercedere dignentur in cœlis, quorum memoriam facimus in terris.}

And although the Church has been accustomed at times to celebrate certain masses in honor and memory of the saints; not therefore, however, doth she teach that sacrifice is offered unto them, but unto God alone, who crowned them; whence neither is the priest wont to say, 'I offer sacrifice to thee, Peter or Paul;' but, giving thanks to God for their victories, he implores their patronage, that they may vouchsafe to intercede for us in heaven, whose memory we celebrate upon earth.

Caput IV.

Chapter IV.

\emph{De canone missæ.}

\emph{On the Canon of the Mass.}

\emph{Et cum sancta sancta administrari conveniat, sitque hoc omnium sanctissimum sacrificium, Ecclesia Catholica, ut digne reverenterque offerretur ac perciperetur, sacrum canonem multis ante sæculis instituit, ita ab omni errore purum, ut nihil in eo contineatur, quod non maxime sanctitatem ac pietatem quamdam redoleat, mentesque oferentium in Deum erigat.}

And whereas it beseemeth that holy things be administered in a holy manner, and of all holy things this sacrifice is the most holy; to the end that it might be worthily and reverently offered and received, the Catholic Church instituted, many years ago, the sacred Canon, so pure from every error, that nothing is contained therein which does not in the highest degree savor of a certain holiness and piety, and raise

\emph{Is enim constat cum ex ipsis Domini verbis, tum ex apostolorum, traditionibus ac sanctorum quoque pontificum piis institutionibus.}

up unto God the minds of those that offer. For it is composed out of the very words of the Lord, the traditions of the Apostles, and the pious institutions also of holy Pontiffs.

Caput V.

Chapter V.

\emph{De missæ ceremoniis et ritibus.}

\emph{On the solemn ceremonies of the Sacrifice of the Mass.}

\emph{Cumque natura hominum ea sit, ut non facile queat sine adminiculis exterioribus ad rerum divinarum meditationem sustolli, propterea pia mater Ecclesia ritus quosdam, ut scilicet quædam summissa voce, alia vero elatiore, in missa pronunciarentur, instituit. Cerimonias item, adhibuit, ut mysticas benedictiones, lumina, thymiamata, vestes, aliaque id genus multa ex apostolica disciplina et traditione, quo et majestas tanti sacrificii commendaretur, et mentes fidelium per hæc visibilia religionis et pietatis signa ad rerum altissimarum, quæ in hoc sacrificio latent, contemplationem excitarentur.}

And whereas such is the nature of man, that, without external helps, he can not easily be raised to the meditation of divine things; therefore has holy Mother Church instituted certain rites, to wit, that certain things be pronounced in the mass in a low, and others in a louder, tone. She has likewise employed ceremonies, such as mystic benedictions, lights, incense, vestments, and many other things of this kind, derived from an apostolical discipline and tradition, whereby both the majesty of so great a sacrifice might be recommended, and the minds of the faithful be excited, by those visible signs of religion and piety, to the contemplation of those most sublime things which are hidden in this sacrifice.

Caput VI.

Chapter VI.

\emph{De missa, in qua solus sacerdos communicat.}

\emph{On Mass wherein the priest alone communicates.}

\emph{Optaret quidem sacrosancta synodus, ut in singulis missis }

The sacred and holy Synod would fain indeed that, at each mass, the

\emph{fideles adstantes non solum spirituali affectu, sed sacramentali etiam Eucharistiæ perceptione communicarent, quod ad eos sanctissimi hujus sacrificii fructus uberior proveniret; nec tamen, si id non semper fiat, propterea missas illas, in quibus solus sacerdos sacramentaliter communicat, ut privatas et illicitas damnat, sed probat atque adeo commendat, siquidem illæ quoque missæ vere communes censeri debent, partim, quod in eis populus spiritualiter communicet, partim vero, quod a publico Ecclesiæ ministro non pro se tantum, sed pro omnibus fidelibus, qui ad corpus Christi pertinent, celebrentur.}

faithful who are present should communicate, not only in spiritual desire, but also by the sacramental participation of the Eucharist, that thereby a more abundant fruit might be derived to them from this most holy sacrifice: but not therefore, if this be not always done, does it condemn, as private and unlawful, but approves of and therefore commends, those masses in which the priest alone communicates sacramentally; since those masses also ought to be considered as truly common; partly because the people communicate spiritually thereat; partly also because they are celebrated by a public minister of the Church, not for himself only, but for all the faithful, who belong to the body of Christ.

Caput VII.

Chapter VII.

\emph{De aqua miscenda vino in calice offerendo.}

\emph{On the water that is to be mixed with the wine to be offered in the chalice.}

\emph{Monet deinde sancta synodus, præceptum esse ab Ecclesia sacerdotibus, ut aquam vino in calice offerendo miscerent, tum quod Christum Dominum ita fecisse credatur, tum etiam quia e latere ejus aqua simul cum sanguine exierit, quod sacramentum hac mixtione recolitur, et, }

The holy Synod notices, in the next place, that it has been enjoined by the Church on priests, to mix water with the wine that is to be offered in the chalice; as well because it is believed that Christ the Lord did this, as also because from \emph{his side there came out blood and water;}\footnote{\par   John xix. 34.} the memory of which mystery is

\emph{cum aquæ in apocalypsi beati Ioannis populi dicantur, ipsius populi fidelis cum capite Christo unio repræsentatur.}

renewed by this commixture; and, whereas in the apocalypse of blessed John \emph{the peoples} are called \emph{waters},\footnote{\par   Apoc. xvii. 15.} the union of that faithful people with Christ their head is thereby represented.

Caput VIII.

Chapter VIII.

\emph{Missa vulgari lingua non celebretur. Ejus mysteria populo explicentur.}

\emph{On not celebrating the Mass every where in the vulgar tongue; the mysteries of the Mass to be explained to the people.}

\emph{Etsi missa magnam contineat populi fidelis eruditionem; non tamen expedire visum est patribus, ut vulgari passim lingua celebraretur. Quamobrem, retento ubique cujusque Ecclesiæ antiquo et a sancta Romana Ecclesia, omnium ecclesiarum matre et magistra, probato ritu, ne oves Christi esuriant, neve parvuli panem petant et non sit qui frangat eis, mandat sancta synodus pastoribus et singulis curam animarum gerentibus, ut frequenter inter missarum celebrationem vel per se vel per alios ex iis, quæ in missa leguntur, aliquid exponant; atque inter cetera sanctissimi hujus sacrificii mysterium aliquod declarent, diebus præsertim dominicis et festis.}

Although the mass contains great instruction for the faithful people, nevertheless, it has not seemed expedient to the Fathers that it should be every where celebrated in the vulgar tongue. Wherefore, the ancient usage of each Church, and the rite approved of by the holy Roman Church, the mother and mistress of all churches, being in each place retained; and, that the sheep of Christ may not suffer hunger, \emph{nor the little ones ask for bread, and there be none to break it unto them},\footnote{\par   Lam. iv. 4.} the holy Synod charges pastors, and all who have the cure of souls, that they frequently, during the celebration of mass, expound either by themselves, or others, some portion of those things which are read at Mass, and that, amongst the rest, they explain some mystery of this most holy sacrifice, especially on the Lord's days and festivals.

Caput IX.

Chapter IX.

\emph{Prolegomenon canonum sequentium.}

\emph{Preliminary Remark on the following Canons.}

\emph{Quia vero adversus veterem hanc in sacrosancto evangelio, apostolorum traditionibus sanctorumque patrum doctrina fundatam fidem hoc tempore multi disseminati sunt errores, multaque a multis docentur et disputantur; sancta synodus, post multos gravesque his de rebus mature habitos tractatus, unanimi patrum omnium concensu quæ huic purissimæ fidei sacræque doctrinæ adversantur damnare et a sancta Ecclesia eliminare, per subjectos hos canones constituit.}

And because that many errors are at this time disseminated and many things are taught and maintained by divers persons, in opposition to this ancient faith, which is based on the sacred Gospel, the traditions of the Apostles, and the doctrine of the holy Fathers; the sacred and holy Synod, after many and grave deliberations maturely had touching these matters, has resolved, with the unanimous consent of all the Fathers, to condemn, and to eliminate from holy Church by means of the canons subjoined, whatsoever is opposed to this most pure faith and sacred doctrine.

DE SACRIFICIO MISSÆ.

ON THE SACRIFICE OF THE MASS.

Canon I.—\emph{Si quis dixerit, in missa non offerri Deo verum et proprium sacrificium, aut quod offerri non sit aliud quam nobis Christum ad manducandum dari: anathema sit.}

Canon I.—If any one saith, that in the mass a true and proper sacrifice is not offered to God; or, that to be offered is nothing else but that Christ is given us to eat: let him be anathema.

Canon II.—\emph{Si quis dixerit, illis verbis: Hoc facite in meam commemorationem, Christum non instituisse apostolos sacerdotes, aut non ordinasse, ut ipsi aliique sacerdotes offerrent corpus et sanguinem suum: anathema sit.}

Canon II.—If any one saith, that by those words, \emph{Do this for the commemoration of me} (Luke xxii. 19), Christ did not institute the apostles priests; or, did not ordain that they and other priests should offer his own body and blood: let him be anathema.

Canon III.—\emph{Si quis dixerit, }

Canon III.—If any one saith,

\emph{missæ sacrificium tantum esse laudis et gratiarum actionis, aut nudam commemorationem sacrificii in cruce peracti, non autem propitiatorium; vel soli prodesse sumenti; neque pro vivis et defunctis pro peccatis, pœnis, satisfactionibus et aliis necessitatibus offerri debere: anathema sit.}

that the sacrifice of the mass is only a sacrifice of praise and of thanksgiving; or, that it is a bare commemoration of the sacrifice consummated on the cross, but not a propitiatory sacrifice; or, that it profits him only who receives; and that it ought not to be offered for the living and the dead for sins, pains, satisfactions, and other necessities: let him be anathema.

Canon IV.—\emph{Si quis dixerit, blasphemiam irrogari sanctissimo Christi sacrificio in cruce peracto per missæ sacrificium, aut illi per hoc derogari: anathema sit.}

Canon IV.—If any one saith, that, by the sacrifice of the mass, a blasphemy is cast upon the most holy sacrifice of Christ consummated on the cross; or, that it is thereby derogated from: let him be anathema.

Canon V.—\emph{Si quis dixerit, imposturam esse, missas celebrare in honorem sanctorum et pro illorum intercessione apud Deum obtinenda, sicut Ecclesia intendit: anathema sit.}

Canon V.—If any one saith, that it is an imposture to celebrate masses in honor of the saints, and for obtaining their intercession with God, as the Church intends: let him be anathema.

Canon VI.—\emph{Si quis dixerit, canonem missæ errores continere, ideoque abrogandum esse: anathema sit.}

Canon VI.—If any one saith, that the canon of the mass contains errors, and is therefore to be abrogated: let him be anathema.

Canon VII.—\emph{Si quis dixerit, ceremonias, vestes et externa signa, quibus in missarum celebratione Ecclesia Catholica utitur, irritabula impietatis esse magis quam officia pietatis: anathema sit.}

Canon VII.—If any one saith, that the ceremonies, vestments, and outward signs, which the Catholic Church makes use of in the celebration of masses, are incentives to impiety, rather than offices of piety: let him be anathema.

Canon VIII.—\emph{Si quis dixerit, missas, in quibus solus sacerdos }

Canon VIII.—If any one saith, that masses, wherein the priest alone

\emph{sacramentaliter communicat, illicitas esse ideoque abrogandas: anathema sit.}

communicates sacramentally, are unlawful, and are, therefore, to be abrogated: let him be anathema.

Canon IX.—\emph{Si quis dixerit, Ecclesiæ Romanæ ritum, quo submissa voce pars canonis et verba consecrationis proferuntur, damnandum esse; aut lingua tantum vulgari missam celebrari debere; aut aquam non miscendam esse vino in calice offerendo, eo quod sit contra Christi institutionem: anathema sit.}

Canon IX.—If anyone saith, that the rite of the Roman Church, according to which a part of the canon and the words of consecration are pronounced in a low tone, is to be condemned; or, that the mass ought to be celebrated in the vulgar tongue only; or, that water ought not to be mixed with the wine that is to be offered in the chalice, for that it is contrary to the institution of Christ: let him be anathema.

\xchapter{Twenty-third Session, held July 15, 1563.}

Sessio Vigesimatertia,

Twenty-third Session,

\emph{celebrata die XV. Iulii} 1563.

\emph{held July} 15, 1563.

VERA ET CATHOLICA DOCTRINA DE SACRAMENTO ORDINIS.

THE TRUE AND CATHOLIC DOCTRINE CONCERNING THE SACRAMENT OF ORDER.

Caput I.

Chapter I.

\emph{De institutione sacerdoti novæ legis.}

\emph{ On the institution of the Priesthood of the New Law.}

\emph{Sacrificium et sacerdotium ita Dei ordinatione conjuncta sunt, ut utrumque in omni lege exstiterit. Cum igitur in Novo Testamento sanctum Eucharistiæ sacrificium visibile ex Domini institutione Catholica Ecclesia acceperit, fateri etiam oportet, in ea novum esse visible et externum sacerdotium, in quod vetus translatum est. Hoc autem ab eodem Domino }

Sacrifice and priesthood are, by the ordinance of God, in such wise conjoined, as that both have existed in every law. Whereas, therefore, in the New Testament, the Catholic Church has received, from the institution of Christ, the holy visible sacrifice of the Eucharist; it must needs also be confessed, that there is, in that Church, a new, visible, and external priesthood, into which the old has been \emph{translated.}\footnote{\par   Heb. vii. 12.}

\emph{Salvatore nostro institutum esse, atque apostolis eorumque successoribus in sacerdotio potestatem traditum consecrandi, offerendi et ministrandi corpus et sanguinem ejus, necnon et peccata dimittendi et retinendi, sacræ litteræ ostendunt et Catholicæ Ecclesiæ traditio semper docuit.}

And the sacred Scriptures show, and the tradition of the Catholic Church has always taught, that this priesthood was instituted by the same Lord our Saviour, and that to the Apostles, and their successors in the priesthood, was the power delivered of consecrating, offering, and administering his body and blood, as also of forgiving and of retaining sins.

Caput II.

Chapter II.

\emph{De septem ordinibus.}

\emph{On the Seven Orders.}

\emph{Cum autem divina res sit tam sancti sacerdotii ministerium, consentaneum fuit, quo dignius et majori cum veneratione exerceri posset, ut in Ecclesiæ ordinatissima dispositione plures et diversi essent ministrorum ordines, qui sacerdotio ex officio deservirent, ita distributi, ut, qui jam clericali tonsura insigniti essent, per minores ad majores ascenderent. Nam non solum de sacerdotibus, sed et de diaconis sacræ litteræ apertam mentionem faciunt, et quæ maxime in illorum ordinatione attendenda sunt gravissimis verbis docent; et ab ipso Ecclesiæ initio sequentium ordinum nomina, atque uniuscujusque eorum propria ministeria, subdiaconi scilicet, }

And whereas the ministry of so holy a priesthood is a divine thing; to the end that it might be exercised in a more worthy manner, and with greater veneration, it was suitable that, in the most well ordered settlement of the Church, there should be several and diverse orders of ministers to minister to the priesthood, by virtue of their office; orders so distributed as that those already marked with the clerical tonsure should ascend through the lesser to the greater orders. For the sacred Scriptures make open mention not only of priests, but also of deacons; and teach, in words the most weighty, what things are especially to be attended to in the Ordination thereof; and, from the very beginning of the Church, the names of the following orders, and

\emph{acolythi, exorcistæ, lectoris et ostiarii in usu fuisse cognoscuntur, quamvis non pari gradu; nam subdiaconatus ad majores ordines a patribus et sacris conciliis refertur, in quibus et de aliis inferioribus frequentissime legimus.}

the ministrations proper to each one of them, are known to have been in use; to wit, those of subdeacon, acolyth, exorcist, lector, and door-keeper; though these were not of equal rank; for the subdeaconship is classed amongst the greater orders by the Fathers and sacred Councils, wherein also we very often read of the other inferior orders.

Caput III.

Chapter III.

\emph{Ordinem vere esse sacramentum.}

\emph{That Order is truly and properly a Sacrament.}

\emph{Cum Scripturæ testimonio, apostolica traditione et patrum unanimi consensu perspicuum sit, per sacram ordinationem, quæ verbis et signis exterioribus perficitur, gratiam conferri, dubitare nemo debet, ordinem esse vere et proprie unum ex septem sanctæ Ecclesiæ sacramentis. Inquit enim apostolus: Admoneo te, ut resuscites gratiam Dei, quæ est in te, per impositionem manuum mearum. Non enim dedit nobis Deus spiritum timoris, sed virtutis et dilectionis et sobrietatis.}

Whereas, by the testimony of Scripture, by Apostolic tradition, and the unanimous consent of the Fathers, it is clear that grace is conferred by sacred ordination, which is performed by words and outward signs, no one ought to doubt that Order is truly and properly one of the seven sacraments of holy Church. For the Apostle says: \emph{I admonish thee that thou stir up the grace of God, which is in thee by the imposition of my hands. For God has not given us the spirit of fear, but of power, and of love, and of sobriety.}\footnote{\par   2 Tim. i. 6, 7.}

Caput IV.

Chapter IV.

\emph{De ecclesiastica hierarchia et ordinatione.}

\emph{On the Ecclesiastical hierarchy, and on Ordination.}

\emph{Quoniam vero in sacramento ordinis, sicut et in baptismo et }

But, forasmuch as in the sacrament of Order, as also in Baptism

\emph{confirmatione, character imprimitur, qui nec deleri nec auferri potest, merito sancta synodus damnat eorum sententiam, qui asserunt Novi Testamenti sacerdotes temporariam tantummodo potestatem habere, et semel rite ordinatos iterum laicos effici posse, si verbi Dei ministerium non exerceant. Quod si quis omnes Christianos promiscue Novi Testamenti sacerdotes esse, aut omnes pari inter se potestate spirituali præditos affirmet, nihil aliud facere videtur, quam ecclesiasticam hierarchiam, quæ est ut castrorum acies ordinata, confundere; perinde ac si contra beati Pauli doctrinam omnes apostoli, omnes prophetæ, omnes evangelistæ, omnes pastores, omnes sint doctores. Proinde sacrosancta synodus declarat, præter ceteros ecclesiasticos gradus episcopos, qui in apostolorum locum successerunt, ad hunc hierarchicum ordinem præcipue pertinere, et positos, sicut idem apostolus ait, a Spiritu Sancto regere Ecclesiam Dei; eosque presbyteris superiores esse, ac sacramentum confirmationis conferre, ministros Ecclesiæ ordinare, atque alia pleraque peragere }

and Confirmation, a character is imprinted which can neither be effaced nor taken away, the holy Synod with reason condemns the opinion of those who assert that the priests of the New Testament have only a temporary power; and that those who have once been rightly ordained can again become laymen, if they do not exercise the ministry of the Word of God. And if any one affirm, that all Christians indiscriminately are priests of the New Testament, or that they are all mutually endowed with an equal spiritual power, he clearly does nothing but confound the ecclesiastical hierarchy, which is \emph{as an army set in array;}\footnote{\par   Cant. vi. 3.} as if, contrary to the doctrine of blessed Paul, \emph{all} were \emph{apostles, all prophets, all evangelists, all pastors, all doctors.} \footnote{\par   Ephes. vi. 11, 12.} Wherefore, the holy Synod declares that, besides the other ecclesiastical degrees, bishops, who have succeeded to the place of the Apostles, principally belong to this hierarchical order; that they are \emph{placed}, as the same apostle says, \emph{by the Holy Ghost, to rule the Church of God;}\footnote{\par   Acts xx. 28.} that they are superior to priests; administer the sacrament of Confirmation; ordain the ministers of the Church; and that they can perform very many other things; over which

\emph{ipsos posse, quarum functionum potestatem reliqui inferioris ordinis nullam habent. Docet insuper sacrosancta synodus, in ordinatione episcoporum, sacerdotum et ceterorum ordinum nec populi nec cujusvis sæcularis potestatis et magistratus consensum sive vocationem sive auctoritatem ita requiri, ut sine ea irrita sit ordinatio; quin potius decernit, eos, qui tantummodo a populo aut sæculari potestate ac magistratu vocati et instituti ad hæc ministeria, exercenda adscendunt, et qui ea propria temeritate sibi sumunt, omnes non Ecclesiæ ministros sed fures et latrones per ostium non ingressos habendos esse. Hæo sunt, quæ generatim sacræ synodo visum est Christi fideles de sacramento ordinis docere. His autem contraria certis et propriis canonibus in hunc, qui sequitur, modum damnare constituit, ut omnes adjuvante Christo fidei regula utentes in tot errorum tenebris Catholicam veritatem facilius agnoscere et tenere possint.}

functions others of an inferior order have no power. Furthermore, the sacred and holy Synod teaches, that, in the ordination of bishops, priests, and of the other orders, neither the consent, nor vocation, nor authority, whether of the people, or of any civil power or magistrate whatsoever, is required in such wise as that, without this, the ordination is invalid: yea rather doth it decree, that all those who, being only called and instituted by the people, or by the civil power and magistrate, ascend to the exercise of these ministrations, and those who of their own rashness assume them to themselves, are not ministers of the Church, but are to be looked upon as \emph{thieves and robbers, who have not entered by the door.}\footnote{\par   John x. 1.} These are the things which it hath seemed good to the sacred Synod to teach the faithful of Christ, in general terms, touching the sacrament of Order. But it hath resolved to condemn whatsoever things are contrary thereunto, in express and specific canons, in the manner following; in order that all men, with the help of Christ, using the rule of faith, may, in the midst of the darkness of so many errors, more easily be able to recognize and to hold Catholic truth.

DE SACRAMENTO ORDINIS.

ON THE SACRAMENT OF ORDER.

Canon I.—\emph{Si quis dixerit, non esse in Novo Testamento sacerdotium visibile et externum, vel non esse potestatem aliquam consecrandi et offerendi verum corpus et sanguinem Domini, et peccata remittendi et retinendi, sed officium tantum et nudum ministerium prædicandi evangelium, vel eos, qui non prædicant, prorsus non esse sacerdotes: anathema sit.}

Canon I.—If any one saith, that there is not in the New Testament a visible and external priesthood; or, that there is not any power of consecrating and offering the true body and blood of the Lord, and of forgiving and retaining sins, but only an office and bare ministry of preaching the Gospel; or, that those who do not preach are not priests at all: let him be anathema.

Canon II.—\emph{Si quis dixerit, præter sacerdotium non esse in Ecclesia Catholica alios ordines et majores et minores, per quos, velut per gradus quosdam, in sacerdotium tendatur: anathema sit.}

Canon II.—If any one saith, that, besides the priesthood, there are not in the Catholic Church other orders, both greater and minor, by which, as by certain steps, advance is made unto the priesthood: let him be anathema.

Canon III.—\emph{Si quis dixerit, ordinem sive sacram ordinationem non esse vere et proprie sacramentum a Christo Domino institutum, vel esse figmentum quoddam humanum, excogitatum a viris rerum ecclesiasticarum imperitis, aut esse tantum ritum quemdam eligendi ministros verbi Dei et sacramentorum: anathema sit.}

Canon III.—If any one saith, that order, or sacred ordination, is not truly and properly a sacrament instituted by Christ the Lord; or, that it is a kind of human figment devised by men unskilled in ecclesiastical matters; or, that it is only a kind of rite for choosing ministers of the Word of God and of the sacraments: let him be anathema.

Canon IV.—\emph{Si quis dixerit, per sacram ordinationem non dari Spiritum Sanctum, ac proinde frustra episcopos dicere: }

Canon IV.—If any one saith, that, by sacred ordination, the Holy Ghost is not given; and that vainly therefore do the bishops say,

\emph{Accipe Spiritum Sanctum, aut per eam non imprimi characterem; vel eum, qui sacerdos semel fuit, laicum rursus fieri posse: anathema sit.}

\emph{Receive ye the Holy Ghost;} or, that a character is not imprinted by that ordination; or, that he who has once been a priest can again become a layman: let him be anathema.

Canon V.—\emph{Si quis dixerit, sacram unctionem, qua Ecclesia in sancta ordinatione utitur, non tantum non requiri, sed contemnendam et perniciosam esse, similiter et alias ordinis ceremonies: anathema sit.}

Canon V.—If any one saith, that the sacred unction which the Church uses in holy ordination is not only not required, but is to be despised and is pernicious, as likewise are the other ceremonies of order: let him be anathema.

Canon VI.—\emph{Si quis dixerit, in Ecclesia Catholica non esse hierarchiam divina ordinatione institutam, quæ constat ex episcopis, presbyteris et ministris: anathema sit.}

Canon VI.—If any one saith, that, in the Catholic Church there is not a hierarchy by divine ordination instituted, consisting of bishops, priests, and ministers: let him be anathema.

Canon VII.—\emph{Si quis dixerit, episcopos non esse presbyteris superiores, vel non habere potestatem confirmandi et ordinandi, vel eam, quam habent, illis esse cum presbyteris communem, vel ordines ab ipsis collatos sine populi vel potestatis sæcularis consensu aut vocatione irritos esse; aut eos qui nec ab ecclesiastica et canonica potestate rite ordinati, nec missi sunt, sed aliunde veniunt, legitimos esse verbi et sacramentorum ministros: anathema, sit.}

Canon VII.—If any one saith, that bishops are not superior to priests; or, that they have not the power of confirming and ordaining: or, that the power which they possess is common to them and to priests; or, that orders, conferred by them, without the consent or vocation of the people, or of the secular power, are invalid; or, that those who have neither been rightly ordained, nor sent, by ecclesiastical and canonical power, but come from elsewhere, are lawful ministers of the Word and of the sacraments: let him be anathema.

Canon VIII.—\emph{Si quis dixerit, episcopos, qui auctoritate Romani }

Canon VIII.—If any one saith, that the bishops, who are assumed

\emph{pontiftcis assumuntur, non esse legitimos et veros episcopos, sed figmentum humanum: anathema sit.}

by authority of the Roman Pontiff, are not legitimate and true bishops, but are a human figment: let him be anathema.

\xchapter{Twenty-fourth Session, held Nov. 11, 1563.}

Sessio Vigesimaquarta,

Twenty-fourth Session,

\emph{celebrata die XI. Nov.} 1563.

\emph{held Nov.} 11, 1563.

DOCTRINA DE SACRAMENTO MATRIMONII.

DOCTRINE ON THE SACRAMENT OF MATRIMONY.

\emph{Matrimonii perpetuum indissolubilemque nexum primus humani generis parens divini Spiritus instinctu pronuntiavit, cum dixit: Hoc nunc os ex ossibus meis et caro de carne mea; quamobrem relinquet homo patrem suum et matrem et adhærebit uxori suæ, et erunt duo in carne una.}

The first parent of the human race, under the influence of the Divine Spirit, pronounced the bond of matrimony perpetual and indissoluble, when he said: \emph{This now is bone of my bones, and flesh of my flesh. Wherefore a man shall leave father and mother, and shall cleave to his wife, and they shall be two in one flesh.}\footnote{\par   Gen. ii. 23, 24.}

\emph{Hoc autem vinculo duos tantummodo copulari et conjungi, Christus Dominus apertius docuit, cum postrema illa verba tamquam a Deo prolata referens dixit: Itaque jam non sunt duo, sed una caro; statimque ejusdem nexus firmitatem ab Adamo tanto ante pronuntiatam his verbis confirmavit: Quod ergo Deus conjunxit, homo non separet.}

But, that by this bond two only are united and joined together, our Lord taught more plainly, when, rehearsing those last words as having been uttered by God, he said: \emph{Therefore now they are not two, but one flesh;}\footnote{\par   Matt. xix. 6.} and straightway confirmed the firmness of that tie, proclaimed so long before by Adam, by these words: \emph{What therefore God hath joined together, let no man put asunder.}\footnote{\par   Matt. xix. 6.}

\emph{Gratiam vero, quæ naturalem illum amorem perficeret et indissolubilem }

But the grace which might perfect that natural love, and confirm

\emph{unitatem confirmaret conjugesque sanctificaret, ipse Christus, venerabilium sacramentorum institutor atque perfector, sua nobis passione promeruit; quod Paulus Apostolus innuit, dicens: Viri, diligite uxores vestras, sicut Christus dilexit Ecclesiam, et seipsum tradidit pro ea; mox subjungens: Sacramentum hoc magnum est, ego autem dico in Christo et in Ecclesia.}

that indissoluble union, and sanctify the married, Christ himself, the institutor and perfecter of the venerable sacraments, merited for us by his passion; as the Apostle Paul intimates, saying, \emph{Husbands love your wives, as Christ also loved the Church, and delivered himself up for it;} adding shortly after, \emph{This is a great sacrament, but I speak in Christ and in the Church.}\footnote{\par  Ephes. v. 25, 32.}

\emph{Cum igitur matrimonium in lege evangelica veteribus connubiis per Christum gratia præstet, merito inter novæ legis sacramenta adnumerandum, sancti patres nostri, concilia, et universalis Ecclesiæ traditio semper docuerunt, adversus quam impii homines hujus sæculi insanientes non solum perperam de hoc venerabili sacramento senserunt, sed de more suo prætextu evangelii libertatem carnis introducentes, multa ab Ecclesiæ Catholicæ sensu et ab apostolorum temporibus probata consuetudine aliena scripto et verbo asseruerunt non sine magna Christi fidelium jactura; quorum temeritati sancta et universalis synodus cupiens occurrere, insigniores prædictorum schismaticorum hæreses et errores, ne plures ad }

Whereas therefore matrimony, in the evangelical law, excels in grace, through Christ, the ancient marriages, with reason have our holy Fathers, the Councils, and the tradition of the universal Church, always taught, that it is to be numbered amongst the sacraments of the new law; against which, impious men of this age raging, have not only had false notions touching this venerable sacrament, but, introducing according to their wont, under the pretext of the Gospel, a carnal liberty, they have by word and writing asserted, not without great injury to the faithful of Christ, many things alien from the sentiment of the Catholic Church, and from the usage approved of since the times of the Apostles; the holy and universal Synod, wishing to meet the rashness of these men, has thought

\emph{se trahat perniciosa eorum contagio, exterminandos duxit, hos in ipsos hæreticos eorumque errores decernens anathematismos.}

it proper, lest their pernicious contagion may draw more after it, that the more remarkable heresies and errors of the above-named schismatics be exterminated, by decreeing against the said heretics and their errors the following anathemas.

DE SACRAMENTO MATRIMONII.

ON THE SACRAMENT OF MATRIMONY.

Canon I.—\emph{Si quis dixerit, matrimonium non esse vere et proprie unum ex septem legis evangelicæ sacramentis a Christo Domino institutum, sed ab hominibus in Ecclesia inventum, neque gratiam conferre: anathema sit.}

Canon I.—If any one saith, that matrimony is not truly and properly one of the seven sacraments of the evangelic law, [a sacrament] instituted by Christ the Lord; but that it has been invented by men in the Church; and that it does not confer grace: let him be anathema.

Canon II.—\emph{Si quis dixerit, licere Christianis plures simul habere uxores, et hoc nulla lege divina esse prohibitum: anathema sit.}

Canon II.—If any one saith, that it is lawful for Christians to have several wives at the same time, and that this is not prohibited by any divine law: let him be anathema.

Canon III.—\emph{Si quis dixerit, eos tantum consanguinitatis et affinitatis gradus, qui Levitico exprimuntur, posse impedire matrimonium contrahendum et dirimere contractum, nec posse Ecclesiam in nonnullis illorum dispensare aut constituere, ut plures impediant et dirimant: anathema sit.}

Canon III.—If any one saith, that those degrees only of consanguinity and affinity which are set down in Leviticus can hinder matrimony from being contracted, and dissolve it when contracted; and that the Church can not dispense in some of those degrees, or establish that others may hinder and dissolve it: let him be anathema.

Canon IV.—\emph{Si quis dixerit, Ecclesiam non potuisse constituere impedimenta matrimonium dirimentia, vel in iis constituendis }

Canon IV.—If any one saith, that the Church could not establish impediments dissolving marriage; or, that she has erred in establishing

\emph{errasse: anathema sit.}

them: let him be anathema.

Canon V.—\emph{Si quis dixerit, propter hæresim, aut molestam cohabitationem, aut affectatam absentiam a conjuge, dissolvi posse matrimonii vinculum : anathema sit.}

Canon V.—If any one saith, that on account of heresy, or irksome cohabitation, or the affected absence of one of the parties, the bond of matrimony may be dissolved: let him be anathema.

Canon VI.—\emph{Si quis dixerit, matrimonium ratum non consummatum per solemnem religionis professionem alterius conjugum non dirimi: anathema sit.}

Canon VI.—If any one saith, that matrimony contracted, but not consummated, is not dissolved by the solemn profession of religion by one of the parties: let him be anathema.

Canon VII.—\emph{Si quis dixerit, Ecclesiam errare, cum docuit et docet juxta evangelicam et apostolicam doctrinam, propter adulterium alterius conjugum matrimonii vinculum non posse dissolvi, et utrumque, vel etiam innocentem, qui causam adulterio non dedit, non posse, altero conjuge vivente, aliud matrimonium contrahere, mæcharique eum, qui, dimissa adultera, aliam duxerit, et eam, quæ, dimisso adultero, alii nupserit: anathema sit.}

Canon VII.—If any one saith, that the Church has erred, in that she hath taught, and doth teach, in accordance with the evangelical and apostolical doctrine, that the bond of matrimony can not be dissolved on account of the adultery of one of the married parties; and that both, or even the innocent one who gave not occasion to the adultery, can not contract another marriage during the lifetime of the other; and, that he is guilty of adultery, who, having put away the adulteress, shall take another wife, as also she, who, having put away the adulterer, shall take another husband: let him be anathema.

Canon VIII.—\emph{Si quis dixerit, Ecclesiam errare, cum ob multas causas separationem inter conjuges quoad thorum seu }

Canon VIII.—If any one saith, that the Church errs, in that she declares that, for many causes, a separation may take place between

\emph{quoad cohabitationem ad certum incertumve tempus fieri posse decernit: anathema sit.}

husband and wife, in regard of bed, or in regard of cohabitation, for a determinate or for an indeterminate period: let him be anathema.

Canon IX.—\emph{Si quis dixerit, clericos in sacris ordinibus constitutos, vel regulares castitatem solemniter professos posse matrimonium contrahere, contractumque validum esse non obstante lege ecclesiastica vel voto; et oppositum nil aliud esse quam damnare matrimonium, posseque omnes contrahere matrimonium, qui non sentiunt se castitatis, etiam si eam voverint, habere donum: anathema sit; cum Deus id recte petentibus non deneget, nec patiatur nos supra id quod possumus, tentari.}

Canon IX.—If anyone saith, that clerics constituted in sacred orders, or regulars, who have solemnly professed chastity, are able to contract marriage, and that being contracted it is valid, notwithstanding the ecclesiastical law, or vow; and that the contrary is nothing else than to condemn marriage; and, that all who do not feel that they have the gift of chastity, even though they have made a vow thereof, may contract marriage: let him be anathema; seeing that God refuses not that gift to those who ask for it rightly, neither does \emph{he suffer us to be tempted above that which we are able.}\footnote{\par   1 Cor. x. 13.}

Canon X.—\emph{Si quis dixerit, statum, conjugalem anteponendum esse statui virginitatis vel cælibatus, et non esse melius ac beatius manere in virginitate aut cælibatu, quam jungi matrimonio: anathema sit.}

Canon X.—If any one saith, that the marriage state is to be placed above the state of virginity, or of celibacy, and that it is not better and more blessed to remain in virginity, or in celibacy, than to be united in matrimony: let him be anathema.

Canon XI.—\emph{Si quis dixerit, prohibitionem solemnitatis nuptiarum certis anni temporibus superstitionem esse tyrannicam ab ethnicorum superstitione profectam, }

Canon XI.—If any one saith, that the prohibition of the solemnization of marriages at certain times of the year is a tyrannical superstition, derived from the superstition of the

\emph{aut benedictiones et alias ceremonias, quibus Ecclesia in illis utitur, damnaverit: anathema sit.}

heathen; or condemn the benedictions and other ceremonies which the Church makes use of therein: let him be anathema.

Canon XII.—\emph{Si quis dixerit, causas matrimoniales non spectare ad judices ecclesiasticos: anathema sit.}

Canon XII.—If any one saith, that matrimonial causes do not belong to ecclesiastical judges: let him be anathema.

\xchapter{Twenty-fifth Session, December 3–4, 1563.}

Sessio Vigesimaquinta,

Twenty-fifth Session,

\emph{cœpta die III. absoluta die IV. Decembris} 1563.

\emph{begun on the third, and terminated on the fourth of December}, 1563

DECRETUM DE PURGATORIO.

DECREE CONCERNING PURGATORY.

\emph{Cum Catholica Ecclesia, Spiritu Sancto edocta ex sacris litteris et antiqua patrum traditione, in sacris conciliis et novissime in hoc œcumenica synodo docuerit, purgatorium esse, animasque ibi detentas, fidelium suffragiis, potissimum vero acceptabili altaris sacrificio, juvari; præcipit sancta synodus episcopis, ut sanam de purgatorio doctrinam a sanctis patribus et sacris conciliis traditam, a Christi fidelibus credi, teneri, doceri et ubique prædicari diligenter studeant.}

Whereas the Catholic Church, instructed by the Holy Ghost, has, from the Sacred Writings and the ancient tradition of the Fathers, taught, in sacred Councils, and very recently in this œcumenical Synod, that there is a Purgatory, and that the souls there detained are helped by the suffrages of the faithful, but principally by the acceptable sacrifice of the altar,—the holy Synod enjoins on bishops that they diligently endeavor that the sound doctrine concerning Purgatory, transmitted by the holy Fathers and sacred Councils, be believed, maintained, taught, and every where proclaimed by the faithful of Christ.

\emph{Apud rudem vero plebem difficiliores ac subtiliores quæstiones, quaæque ædificationem non faciunt, }

But let the more difficult and subtle questions, and which tend not to edification, and from which for the

\emph{et ex quibus plerumque nulla fit pietatis accessio, a popularibus concionibus secludantur. Incerta item, vel quæ specie falsi laborant, evulgari ac tractari non permittant. Ea vero, quæ ad curiositatem quamdam aut superstitionem spectant, vel turpe lucrum sapiunt, tamquam scandala et fidelium offendicula prohibeant.}

most part there is no increase of piety, be excluded from popular discourses before the uneducated multitude. In like manner, such things as are uncertain, or which labor under an appearance of error, let them not allow to be made public and treated of. While those things which tend to a certain kind of curiosity or superstition, or which savor of filthy lucre, let them prohibit as scandals and stumbling-blocks of the faithful. But let the bishops take care that the suffrages of the faithful who are living, to wit, the sacrifices of masses, prayers, alms, and other works of piety, which have been wont to be performed by the faithful for the other faithful departed, be piously and devoutly performed, in accordance with the institutes of the Church; and that whatsoever is due on their behalf, from the endowments of testators, or in other way, be discharged, not in a perfunctory manner, but diligently and accurately, by the priests and ministers of the Church, and others who are bound to render this [service].

\emph{Curent autem episcopi, ut fidelium vivorum suffragia, missarum scilicet sacrificia, orationes, eleemosynæ, aliaque pietatis opera, quæ a fidelibus pro aliis fidelibus defunctis fieri consueverunt, secundum Ecclesiæ instituta pie et devote fiant; et quæ pro illis ex testatorum fundationibus vel alia ratione debentur, non perfunctorie, sed a sacerdotibus et Ecclesia ministris et aliis, qui hoc præstare tenentur, diligenter et accurate persolvantur.}

DE INVOCATIONE, VENERATIONE, ET RELIQUIIS SANCTORUM, ET SACRIS IMAGINIBUS.

ON THE INVOCATION, VENERATION, AND RELICS OF SAINTS, AND ON SACRED IMAGES.

\emph{Mandat sancta synodus omnibus episcopis et ceteris docendi }

The holy Synod enjoins on all bishops, and others who sustain the

\emph{munus curamque sustinentibus, ut juxta Catholicæ et Apostolicæ Ecclesiæ usum a primævis Christianæ religionis temporibus receptum sanctorumque patrum consensionem et sacrorum conciliorum decreta in primis de sanctorum intercessione, invocatione, reliquiarum honore et legitimo imaginum usu, fideles diligenter instruant, docentes eos, sanctos una cum Christo regnantes orationes suas pro hominibus Deo offerre; bonum, atque utile esse, suppliciter eos invocare; et ob beneficia impenetranda a Deo per filium ejus Iesum Christum Dominum nostrum, qui solus noster redemptor et salvator est, ad eorum orationes, opem, auxiliumque confugere; illos vero, qui negant, sanctos æterna felicitate in cœlo fruentes invocandos esse; aut qui asserunt, vel illos pro hominibus non orare, vel eorum, ut pro nobis etiam singulis orent, invocationem esse idololatriam, vel pugnare cum verbo Dei, adversarique honori unius mediatoris Dei et hominum Iesu Christi, vel stultum esse, in cœlo regnantibus voce vel mente supplicare, impie sentire.}

office and charge of teaching, that, agreeably to the usage of the Catholic and Apostolic Church, received from the primitive times of the Christian religion, and agreeably to the consent of the holy Fathers, and to the decrees of sacred Councils, they especially instruct the faithful diligently concerning the intercession and invocation of saints; the honor [paid] to relics; and the legitimate use of images: teaching them, that the saints, who reign together with Christ, offer up their own prayers to God for men; that it is good and useful suppliantly to invoke them, and to have recourse to their prayers, aid, [and] help for obtaining benefits from God, through his Son, Jesus Christ our Lord, who is our alone Redeemer and Saviour; but that they think impiously who deny that the saints, who enjoy eternal happiness in heaven, are to be invocated; or who assert either that they do not pray for men; or that the invocation of them to pray for each of us even in particular is idolatry; or that it is repugnant to the Word of God, and is opposed to the honor of the \emph{one mediator of God and men, Christ Jesus;}\footnote{\par   1 Tim. ii. 5.} or that it is foolish to supplicate, vocally or mentally, those who reign in heaven.

\emph{Sanctorum quoque martyrum et aliorum cum Christo viventium sancta corpora, quæ viva membra fuerunt Christi et templum Spiritus Sancti, ab ipso ad æternam vitam suscitanda et glorificanda, a fidelibus veneranda esse, per quæ multa beneficia a Deo hominibus præstantur; ita ut affirmantes, sanctorum reliquiis venerationem atque honorem non deberi; vel eas aliaque sacra monumenta a fidelibus inutiliter honorari, atque eorum opis impetrandæ causa sanctorum memorias frustra frequentari; omnino damnandos esse, prout jam pridem eos damnavit, et nunc etiam damnat Ecclesia.}

Also, that the holy bodies of holy martyrs, and of others now living with Christ,—which bodies were the living members of Christ, and \emph{the temple of the Holy Ghost},\footnote{\par   1 Cor. iii. 6.} and which are by him to be raised unto eternal life, and to be glorified,—are to be venerated by the faithful; through which [bodies] many benefits are bestowed by God on men; so that they who affirm that veneration and honor are not due to the relics of saints; or that these, and other sacred monuments, are uselessly honored by the faithful; and that the places dedicated to the memories of the saints are in vain visited with the view of obtaining their aid, are wholly to be condemned, as the Church has already long since condemned, and now also condemns them.

\emph{Imagines porro Christi, Deiparæ Virginis et aliorum sanctorum in templis præsertim habendas et retinendas, eisque debitum honorem et venerationem impertiendam; non quod credatur in esse aliqua in iis divinitas vel virtus, propter quam sint colendæ, vel quod ab eis sit aliquid petendum, vel quod fiducia in imaginibus sit figenda veluti olim fiebat a gentibus, quæ in idolis spem suam}

Moreover, that the images of Christ, of the Virgin Mother of God, and of the other saints, are to be had and retained particularly in temples, and that due honor and veneration are to be given them; not that any divinity, or virtue, is believed to be in them, on account of which they are to be worshipped; or that any thing is to be asked of them; or that trust is to be reposed in images, as was of old done by the Gentiles, who placed their hope

\emph{collocabant; sed quoniam honos, qui eis exhibetur, refertur ad prototypa, quæ illæ repræsentant, ita ut per imagines, quas osculamur et coram quibus caput aperimus et procumbimus, Christum adoremus, et sanctos, quorum illæ similitudinem gerunt, veneremur: id quod conciliorum præsertim vero secundæ Nicænæ Synodi decretis contra imaginum oppugnatores est sancitum.}

in idols; but because the honor which is shown them is referred to the prototypes which those images represent; in such wise that by the images which we kiss, and before which we uncover the head, and prostrate ourselves, we adore Christ, and we venerate the saints, whose similitude they bear: as, by the decrees of Councils, and especially of the second Synod of Nicæa, has been defined against the opponents of images.

\emph{Illud vero diligenter doceant episcopi, per historias mysteriorum nostræ redemptionis picturis vel aliis similitudinibus expressas erudiri et confirmari populum in articulis fidei commemorandis et assidue recolendis; tum vero ex omnibus sacris imaginibus magnum fructum percipi, non solum quia admonetur populus beneficiorum et munerum, quæ a Christo sibi collata sunt, sed etiam quia Dei per sanctos miracula et salutaria exempla oculis fidelium subjiciuntur, ut pro iis Deo gratias agant, ad sanctorumque imitationem vitam moresque suos componant, excitenturque ad adorandum ac diligendum Deum et ad pietatem colendam. Si quis autem his decretis contraria docuerit }

And the bishops shall carefully teach this,—that, by means of the histories of the mysteries of our Redemption, portrayed by paintings or other representations, the people is instructed, and confirmed in [the habit of] remembering, and continually revolving in mind the articles of faith; as also that great profit is derived from all sacred images, not only because the people are thereby admonished of the benefits and gifts bestowed upon them by Christ, but also because the miracles which God has performed by means of the saints, and their salutary examples, are set before the eyes of the faithful; that so they may give God thanks for those things; may order their own lives and manners in imitation of the saints; and may be excited to adore and love God, and to cultivate piety. But if any one

\emph{aut senserit: anathema sit.}

shall teach or entertain sentiments contrary to these decrees: let him be anathema.

\emph{In has autem sanctas et salutares observationes si qui abusus irrepserint, eos prorsus aboleri sancta synodus vehementer cupit; ita ut nullæ falsi dogmatis imagines et rudibus periculosi erroris occasionem præbentes, statuantur. Quod si aliquando historias et narrationes sacræ scripturæ, cum id indoctæ plebi expediet, exprimi et figurari contigerit, doceatur populus, non propterea divinitatem figurari, quasi corporeis oculis conspici vel coloribus, aut figuris exprimi possit.}

And if any abuses have crept in amongst these holy and salutary observances, the holy Synod ardently desires that they be utterly abolished; in such wise that no images [suggestive] of false doctrine, and furnishing occasion of dangerous error to the uneducated, be set up. And if at times, when expedient for the unlettered people, it happen that the facts and narratives of sacred Scripture are portrayed and represented, the people shall be taught, that not thereby is the Divinity represented, as though it could be seen by the eyes of the body, or be portrayed by colors or figures.

\emph{Omnis porro superstitio in sanctorum invocatione, reliquiarum veneratione et imaginum sacro usu tollatur, omnis turpis quæstus eliminetur, omnis denique lascivia vitetur; ita ut procaci venustate imagines non pingantur nec ornentur, et sanctorum celebratione ac reliquiarum visitatione homines ad commessationes atque ebrietates non abutantur, quasi festi dies in honorem sanctorum per luxum ac lasciviam agantur.}

Moreover, in the invocation of saints, the veneration of relics, and the sacred use of images, every superstition shall be removed, all filthy lucre be abolished; finally, all lasciviousness be avoided; in such wise that figures shall not be painted or adorned with a beauty exciting to lust; nor the celebration of the saints and the visitation of relics be by any perverted into revelings and drunkenness; as if festivals were celebrated to the honor of the saints by luxury and wantonness.

\emph{Postremo, tanta circa hæc digentia }

In fine, let so great care and diligence

\emph{et cura ab episcopis adhibeatur, ut nihil inordinatum aut præpostere et tumultuarie accomodatum, nihil profanum nihilque inhonestum appareat, cum domum Dei deceat sanctitudo.}

be used herein by bishops, as that there be nothing seen that is disorderly, or that is unbecomingly or confusedly arranged, nothing that is profane, nothing indecorous, seeing that \emph{holiness becometh the house of God.}\footnote{\par   Psa. xcii. 5.}

\emph{Hæc ut fidelius observentur, statuit sancta synodus, nemini licere ullo in loco vel ecclesia, etiam quomodolibet exempta, ullam insolitam ponere vel ponendam curare imaginem, nisi ab episcopo approbata fuerit; nulla etiam admittenda esse nova miracula, nec novas reliquias recipiendas, nisi eodem recognoscente et approbante episcopo, qui, simul atque de iis aliquid compertum habuerit, adhibitis in consilium theologis et aliis piis viris, ea faciat, quæ veritati et pietati consentanea judicaverit.}

And that these things may be the more faithfully observed, the holy Synod ordains, that no one be allowed to place, or cause to be placed, any unusual image, in any place or church, howsoever exempted, except that image has been approved of by the bishop; also, that no new miracles are to be acknowledged, or new relics recognized, unless the said bishop has taken cognizance and approved thereof; who, as soon as he has obtained some certain information in regard of these matters, shall, after having taken the advice of theologians, and of other pious men, act therein as he shall judge to be consonant with truth and piety. But if any doubtful or difficult abuse has to be extirpated; or, in fine, if any more grave question shall arise touching these matters, the bishop, before deciding the controversy, shall await the sentence of the metropolitan and of the bishops of the province, in a provincial Council; yet so that nothing new, or that previously has

\emph{Quod si aliquis dubius, aut difficilis abusus sit exstirpandus, vel omnino aliqua de iis rebus gravior quæstio incidat, episcopus, antequam controversiam dirimat, metropolitani et comprovincialium episcoporum in concilio provinciali sententiam exspectet, ita tamen, ut nihil inconsulto }

\emph{sanctissimo Romano pontifice novum aut in Ecclesia hactenus inusitatum decernatur.}

not been usual in the Church, shall be resolved on without having first consulted the most holy Roman Pontiff.

\xchapter{Continuation of the Twenty-fifth Session, December 4, 1563.}

Continuatio Sessionis

Continuation of the Session,

\emph{die IV. Decembris.}

\emph{on the fourth day of December.}

DECRETUM DE INDULGENTIIS.

DECREE CONCERNING INDULGENCES.

\emph{Cum potestas conferendi indulgentias a Christo Ecclesiæ concessa sit, atque hujusmodi potestate divinitus sibi tradita antiquissimis etiam temporibus illa usa fuerit, sacrosancta synodus indulgentiarum usum, Christiano populo maxime salutarem et sacrorum conciliorum auctoritate probatum, in Ecclesia retinendum esse docet et præcipit, eosque anathemate damnat, qui aut inutiles esse asserunt, vel eas concedendi in Ecclesia potestatem esse negant. In his tamen concedendis moderationem juxta veterem et probatam in Ecclesia consuetudinem adhiberi cupit, ne nimia facilitate ecclesiastica disciplina enervetur.}

Whereas the power of conferring Indulgences was granted by Christ to the Church, and she has, even in the most ancient times, used the said power delivered unto her of God, the sacred holy Synod teaches and enjoins that the use of Indulgences, for the Christian people most salutary, and approved of by the authority of sacred Councils, is to be retained in the Church; and it condemns with anathema those who either assert that they are useless, or who deny that there is in the Church the power of granting them. In granting them, however, it desires that, in accordance with the ancient and approved custom in the Church, moderation be observed; lest, by excessive facility, ecclesiastical discipline be enervated. And being desirous that the abuses which have crept therein, and by occasion of which this honorable name of Indulgences is blasphemed by heretics, be amended and corrected, it

\emph{Abusus vero, qui in his irrepserunt, et quorum occasione insigne hoc indulgentiarum nomen ab hæreticis blasphematur, emendatos et correctos }

\emph{cupiens, præsenti decreto generaliter statuit, pravos quæstus omnes pro his consequendis, unde plurima in Christiano populo abusuum causa fluxit, omnino abolendos esse. }

ordains generally by this decree, that all evil gains for the obtaining thereof,—whence, a most prolific cause of abuses amongst the Christian people has been derived,—be wholly abolished. But as regards the other abuses which have proceeded from superstition, ignorance, irreverence, or from whatsoever other source, since, by reason of the manifold corruptions in the places and provinces where the said abuses are committed, they can not conveniently be specially prohibited, it commands all bishops diligently to collect, each in his own Church, all abuses of this nature, and to report them in the first provincial Synod; that, after having been reviewed by the opinions of the other bishops also, they may forthwith be referred to the Sovereign Roman Pontiff, by whose authority and prudence that which may be expedient for the universal Church will be ordained; that thus the gift of holy Indulgences may be dispensed to all the faithful, piously, holily, and incorruptly.

\emph{Ceteros vero, qui ex superstitione, ignorantia, irreverentia, aut aliunde quomodocumque provenerunt, cum ob multiplices locorum et provinciarum, apud quas hi committuntur, corruptelas commode nequeant specialiter prohiberi; mandat omnibus episcopis, ut diligenter quisque hujusmodi abusus Ecclesiæ suæ colligat, eosque in prima synodo provinciali referat; ut, aliorum quoque episcoporum sententia cognita, statim ad summum Romanum pontificem deferantur, cujus auctoritate et prudentia, quod universali Ecclesiæ expediet, statuatur; ut ita sanctarum indulgentiarum munus pie, sancte et incorrupte omnibus fidelibus dispensetur.}

\xchapter{Profession of the Tridentine Faith. A.D. 1564}

[From the bulls of Pope Pius IV., '\emph{Injunctum nobis},' Nov. 13, 1564, and '\emph{In sacrosancta},' Dec. 9, 1564 (in the \emph{Bullar. Rom.}, also in Streitwolf and Klener, \emph{Libri Symb. Eccles. Cath.} Tom. II. pp. 315–321). The Latin text of the Creed is given also by Streitwolf and Klener (Tom. I. p. 98, sub tit.: \emph{Forma juramenti professionis fidei}), by Denzinger, and in other collections of Roman Symbols. See Vol. I. § 25, pp. 96–99.]

I. \emph{Ego —— firma fide credo et profiteor omnia et singula, quæ continentur in symbolo fidei, quo sancta Romana Ecclesia utitur, videlicet:}

I. I, —— with a firm faith believe and profess all and every one of the things contained in that creed which the holy Roman Church makes use of:

\emph{'Credo in unum Deum, Patrem omnipotontem},' etc. [\emph{Symbolum Nicenum.} See p. 27.]

'I believe in one God, the Father Almighty,' etc. [The Nicene Creed. See pp. 27 and 98.]

II. \emph{Apostolicas et ecclesiasticas traditiones, reliquasque ejusdem Ecclesiæ observationes et constitutiones firmissime admitto et amplector.}

II. I most steadfastly admit and embrace apostolic and ecclesiastic traditions, and all other observances and constitutions of the same Church.

III. \emph{Item sacram Scripturam juxta eum sensum, quem tenuit et tenet sancta mater Ecclesia, cujus est judicare de vero sensu et interpretatione sacrarum Scripturarum, admitto; nec eam unquam, nisi juxta unanimem consensum patrum accipiam et interpretabor.}

III. I also admit the holy Scriptures, according to that sense which our holy mother Church has held and does hold, to which it belongs to judge of the true sense and interpretation of the Scriptures; neither will I ever take and interpret them otherwise than according to the unanimous consent of the Fathers.

IV. \emph{Profiteor quoque, septem esse vere et proprie sacramenta novæ legis a Jesu Christo Domino nostro instituta, atque ad salutem humani generis, licet non omnia singulis, necessaria: scilicet baptismum, confirmationem, }

IV. I also profess that there are truly and properly seven sacraments of the new law, instituted by Jesus Christ our Lord, and necessary for the salvation of mankind, though not all for every one, to wit: baptism, confirmation,

\emph{eucharistiam, pœnitentiam, extremam unctionem, ordinem et matrimonium; illaque gratiam conferre; et ex his baptismum, confirmationem et ordinem sine sacrilegio reiterare non posse. Receptos quoque et approbatos Ecclesiæ Catholicæ ritus in supradictorum omnium sacramentorum solemni administratione recipio et admitto.}

the eucharist, penance, extreme unction, holy orders, and matrimony; and that they confer grace; and that of these, baptism, confirmation, and ordination can not be reiterated without sacrilege. I also receive and admit the received and approved ceremonies of the Catholic Church, used in the solemn administration of the aforesaid sacraments.

V. \emph{Omnia et singula, quæ de peccato originali et de justificatione in sacrosancta Tridentina synodo definita et declarata fuerunt, amplector et recipio.}

V. I embrace and receive all and every one of the things which have been defined and declared in the holy Council of Trent concerning original sin and justification.

VI. \emph{Profiteor pariter, in missa offerri Deo verum, proprium et propitiatorium sacrificium pro vivis et defunctis; atque in sanctissimo eucharistiæ sacramento esse vere, realiter et substantialiter corpus et sanguinem, una cum anima et divinitate Domini nostri Jesu Christi, fierique conversionem totius substantiæ panis in corpus et totius substantiæ vini in sanguinem; quam conversionem Catholica Ecclesia transsubstantiationem appellat.}

VI. I profess, likewise, that in the mass there is offered to God a true, proper, and propitiatory sacrifice for the living and the dead; and that in the most holy sacrament of the eucharist there is truly, really, and substantially, the body and blood, together with the soul and divinity of our Lord Jesus Christ; and that there is made a change of the whole essence of the bread into the body, and of the whole essence of the wine into the blood; which change the Catholic Church calls transubstantiation.

VII. \emph{Fateor etiam, sub altera tantum specie totum atque integrum Christum, verumque sacramentum sumi.}

VII. I also confess that under either kind alone Christ is received whole and entire, and a true sacrament.

VIII. \emph{Constanter teneo, purgatorium }

VIII. I firmly hold that there is

\emph{esse, animasque ibi detentas fidelium suffragiis juvari. Similiter et sanctos una cum Christo regnantes venerandos atque invocandos esse, eosque orationes Deo pro nobis offerre, atque eorum reliquias esse venerandas.}

a purgatory, and that the souls therein detained are helped by the suffrages of the faithful. Likewise, that the saints reigning with Christ are to be honored and invoked, and that they offer up prayers to God for us, and that their relics are to be had in veneration.

IX. \emph{Firmissime} \footnote{\par Bullarium Rom.: \emph{firmiter.}}  \emph{assero, imagines Christi ac Deiparæ semper Virginis, nec non aliorum sanctorum habendas et retinendas esse, atque eis debitum honorem ac venerationem impertiendam. Indulgentiarum etiam potestatem a Christo in Ecclesia relictam fuisse, illarumque usum Christiano populo maxime salutarem esse affirmo.}

IX. I most firmly assert that the images of Christ, and of the perpetual Virgin the Mother of God, and also of other saints, ought to be had and retained, and that due honor and veneration are to be given them. I also affirm that the power of indulgences was left by Christ in the Church, and that the use of them is most wholesome to Christian people.

X. \emph{Sanctam Catholicam et Apostolicam Romanam Ecclesiam omnium ecclesiarum matrem et magistram agnosco, Romano que pontifici, beati Petri apostolorum principis successori ac Jesu Christi vicario veram obedientiam spondeo ac juro.}

X. I acknowledge the holy Catholic Apostolic Roman Church for the mother and mistress of all churches; and I promise and swear true obedience to the Bishop of Rome, successor to St. Peter, Prince of the Apostles, and Vicar of Jesus Christ.

XI. \emph{Cætera item omnia a sacris canonibus et œcumenicis conciliis, ac præcipue a sacrosancta Tridentina synodo tradita, definita et declarata indubitanter recipio atque profiteor; simulque contraria omnia, atque hæreses quascumque ab Ecclesia }

XI. I likewise undoubtingly receive and profess all other things delivered, defined, and declared by the Sacred Canons and General Councils, and particularly by the holy Council of Trent; and I condemn, reject, and anathematize all things contrary thereto, and all

\emph{damnatas, rejectas et anathematizatas ego pariter damno, rejicio et anathematizo.}

heresies which the Church has condemned, rejected, and anathematized.

XII. \emph{Hanc veram Catholicam fidem, extra quam nemo salvus esse potest, quam in præsenti sponte profiteor et veraciter teneo, eundem integram et inviolatam} \footnote{\par [Note.—As it was promulgated by Pius IX., Jan. 20, 1877—\emph{Acta sedis sanc.} X., 382—and is now offered to Catholic priests and professors, Pius IV.'s Profession contains in article XI, after the words \emph{Tridentino synodo}, the clause \emph{et ab œcumenico concilio Vaticano} (\emph{tradita, definita et declarata}) \emph{præsertim de Romani pontificis primatu ac infallibili magisterio.} The insertion conforms to Pius IX.'s letter to a German bishop, Nov. 6, 1876, that it is altogether necessary that priests with full and unreserved assent of will accept the definition of papal infallibility unless they want to abandon the right faith,  \emph{pleno et absoluto intellectus et voluntatis assensu definitionem complectantur, nisi a recta fide aberrare velint.}  In the same letter, Pius wrote that 'nothing could be more absurd than to think that the Holy Spirit would vouchsafe truths and that, at the same time, it might be inopportune to teach them.' The Profession is printed with the insertion in Benedict's Code of Canon Law.—Ed.]}  \emph{usque ad extremum vitæ spiritum constantissime, Deo adjuvante, retinere et confiteri, atque a meis subditis vel illis, quorum cura ad me in munere meo spectabit, teneri, doceri et prædicari, quantum in me erit, curaturum. Ita ego idem —— spondeo, voveo ac juro. Sic me Deus adjuvet, et hæc sancta Dei Evangelia.}

XII. I do, at this present, freely profess and truly hold this true Catholic faith, without which no one can be saved; and I promise most constantly to retain and confess the same entire and inviolate, with God's assistance, to the end of my life. And I will take care, as far as in me lies, that it shall be held, taught, and preached by my subjects, or by those the care of whom shall appertain to me in my office. This I promise, vow, and swear—so help me God, and these holy Gospels of God.
